
\begin{frame}
\frametitle{{Index}}
$\bullet$ \hyperlink{main.cpp}{main.cpp}\newline
$\bullet$ \hyperlink{Board.hpp}{Board.hpp}\newline
$\bullet$ \hyperlink{Board.tpp}{Board.tpp}\newline
$\bullet$ \hyperlink{mesure.tpp}{mesure.tpp}\newline
$\bullet$ \hyperlink{move.tpp}{move.tpp}\newline
$\bullet$ \hyperlink{getprobamove.tpp}{get-proba-move.tpp}\newline
$\bullet$ \hyperlink{Piece.hpp}{Piece.hpp}\newline
$\bullet$ \hyperlink{Piece.tpp}{Piece.tpp}\newline
$\bullet$ \hyperlink{getlistmove.tpp}{get-list-move.tpp}\newline
$\bullet$ \hyperlink{TypePiece.hpp}{TypePiece.hpp}\newline
$\bullet$ \hyperlink{Coord.hpp}{Coord.hpp}\newline
$\bullet$ \hyperlink{Coord.cpp}{Coord.cpp}\newline
$\bullet$ \hyperlink{Move.hpp}{Move.hpp}\newline
$\bullet$ \hyperlink{Move.cpp}{Move.cpp}\newline

\end{frame}
\begin{frame}
$\bullet$ \hyperlink{Qubit.hpp}{Qubit.hpp}\newline
$\bullet$ \hyperlink{Qubit.tpp}{Qubit.tpp}\newline
$\bullet$ \hyperlink{Random.hpp}{Random.hpp}\newline
$\bullet$ \hyperlink{Random.cpp}{Random.cpp}\newline
$\bullet$ \hyperlink{mathutility.hpp}{math-utility.hpp}\newline
$\bullet$ \hyperlink{mathutility.tpp}{math-utility.tpp}\newline
$\bullet$ \hyperlink{checkpath.hpp}{check-path.hpp}\newline
$\bullet$ \hyperlink{checkpath.tpp}{check-path.tpp}\newline
$\bullet$ \hyperlink{Unitary.hpp}{Unitary.hpp}\newline
$\bullet$ \hyperlink{CMatrix.hpp}{CMatrix.hpp}\newline
$\bullet$ \hyperlink{Complexprinter.hpp}{Complex-printer.hpp}\newline
$\bullet$ \hyperlink{Constexpr.hpp}{Constexpr.hpp}\newline
$\bullet$ \hyperlink{ConsoleInterface.hpp}{ConsoleInterface.hpp}\newline
$\bullet$ \hyperlink{ConsoleInterface.tpp}{ConsoleInterface.tpp}\newline
$\bullet$ \hyperlink{finalsolver.hpp}{final-solver.hpp}\newline

\end{frame}
\begin{frame}
$\bullet$ \hyperlink{finalsolver.tpp}{final-solver.tpp}\newline
$\bullet$ \hyperlink{Matrix.hpp}{Matrix.hpp}\newline
$\bullet$ \hyperlink{Matrix.tpp}{Matrix.tpp}\newline
$\bullet$ \hyperlink{testmovepromotion.cpp}{test-move-promotion.cpp}\newline
$\bullet$ \hyperlink{testmoveenpassantwithcapture.cpp}{test-move-enpassant-with-capture.cpp}\newline
$\bullet$ \hyperlink{testmovepawntwostep.cpp}{test-move-pawn-two-step.cpp}\newline
$\bullet$ \hyperlink{testmovepromotionwithcapture.cpp}{test-move-promotion-with-capture.cpp}\newline
$\bullet$ \hyperlink{testmovepromotionwithsplitbefore.cpp}{test-move-promotion-with-split-before.cpp}\newline
$\bullet$ \hyperlink{testmovesplitwithmesure.cpp}{test-move-split-with-mesure.cpp}\newline
$\bullet$ \hyperlink{testsplitinnonemptycase.cpp}{test-split-in-non-empty-case.cpp}\newline
$\bullet$ \hyperlink{testslidemergemovethroughasplitpiece.cpp}{test-slide-merge-move-through-a-split-piece.cpp}\newline
$\bullet$ \hyperlink{testgetprobamoveslidecapture.cpp}{test-get-proba-move-slide-capture.cpp}\newline
$\bullet$ \hyperlink{testgetprobamoveslideonsamecolor.cpp}{test-get-proba-move-slide-on-same-color.cpp}\newline
$\bullet$ \hyperlink{testgetprobamoveslidewithouttarget.cpp}{test-get-proba-move-slide-without-target.cpp}\newline
$\bullet$ \hyperlink{testgetprobamovejumpsamecolor.cpp}{test-get-proba-move-jump-same-color.cpp}\newline

\end{frame}
\begin{frame}
$\bullet$ \hyperlink{testgetprobamovejumpcapture.cpp}{test-get-proba-move-jump-capture.cpp}\newline
$\bullet$ \hyperlink{testcaptureslidemovetrue.cpp}{test-capture-slide-move-true.cpp}\newline
$\bullet$ \hyperlink{testcaptureslidemovefalse.cpp}{test-capture-slide-move-false.cpp}\newline
$\bullet$ \hyperlink{testmovemergeaftersplit.cpp}{test-move-merge-after-split.cpp}\newline
$\bullet$ \hyperlink{testsplitaftersplit.cpp}{test-split-after-split.cpp}\newline
$\bullet$ \hyperlink{testsplitaftersplit2.cpp}{test-split-after-split2.cpp}\newline
$\bullet$ \hyperlink{testmultiplesplit.cpp}{test-multiple-split.cpp}\newline
$\bullet$ \hyperlink{testsplittakes.cpp}{test-split-takes.cpp}\newline

\end{frame}{\tiny
\begin{frame}[fragile]
\label{main.cpp}
\frametitle{main.cpp}
\hyperlink{conclusion}{Retour au début}\newline
\begin{minted}[obeytabs=true,tabsize=2,linenos=true, breaklines]{cpp}
#include <iostream>
#include <complex>
#include <Complex_printer.hpp>
#include <CMatrix.hpp>
#include <Unitary.hpp>
#include <Qubit.hpp>
#include <Board.hpp>
#include <Piece.hpp>
#include <Move.hpp>
#include <observer_ptr.hpp>
#include <check_path.hpp>
#include <Constexpr.hpp>
#include <TypePiece.hpp>
#include <ConsoleInterface.hpp>
#include <final_solver.hpp>
#include <functional>


int main()
{
    Board<4> B{
        {{B_KING, B_KNIGHT, Piece(), Piece()},
         {B_BISHOP, Piece(), Piece(), Piece()},
         {Piece(), Piece(),  Piece(), Piece()},
         {Piece(), Piece(), W_QUEEN, W_KING}}};
Board<4> B1{
\end{minted}
\end{frame}

\begin{frame}[fragile]
\hyperlink{conclusion}{Retour au début}\newline
\begin{minted}[obeytabs=true,tabsize=2,linenos=true,breaklines,firstnumber=27]{cpp}
        {{B_KING, Piece(), Piece(), B_BISHOP},
         {Piece(), B_KNIGHT, Piece(), Piece()},
         {Piece(), Piece(),  Piece(), Piece()},
         {Piece(), Piece(), W_QUEEN, W_KING}}};

         Board<4> B3{
        {{Piece(), Piece(), B_KNIGHT, B_BISHOP},
         {Piece(), B_KING, Piece(), Piece()},
         {Piece(), Piece(),  Piece(), Piece()},
         {Piece(), Piece(), W_QUEEN, W_KING}}};


         Board<4> Bt{
        {{B_KING, B_BISHOP, Piece(), Piece()},
         {Piece(), Piece(), Piece(), Piece()},
         {Piece(), Piece(),  Piece(), W_KNIGHT},
         {Piece(), Piece(), W_KNIGHT, W_KING}}};

         Board<3> Bt1{
        {{B_KING, Piece(), Piece()},
         {Piece(), W_ROOK, Piece()},
         {Piece(), Piece(),  W_KING}}
         };
         Board<3> Bt2{
        {{B_KING, Piece(), Piece()},
         {Piece(), W_KING, Piece()},
         {Piece(), Piece(), Piece()}}
         };


\end{minted}
\end{frame}

\begin{frame}[fragile]
\hyperlink{conclusion}{Retour au début}\newline
\begin{minted}[obeytabs=true,tabsize=2,linenos=true,breaklines,firstnumber=57]{cpp}
    Board<1,6> B2{
        {
            {W_KING, W_ROOK, Piece(), Piece(), B_KING, Piece()}
        }
    };
    /*bool b = Final::brut_force_quantum_chess(B2, 1, Color::WHITE);
    std::cout<<std::boolalpha<<b<<std::en
    dl;*/
    auto[res, res2] {Final::res_pos(B,4)};
    std::cout<<res<<std::endl;
    Final::print_stack<4>(res2);
    //bool b {Final::brut_force_classic_chess (B1, 10, Color::WHITE)};
    //std::cout<< std::boolalpha<< b<< std::endl;
   

\end{minted}
\end{frame}

\begin{frame}[fragile]
\label{Board.hpp}
\frametitle{Board.hpp}
\hyperlink{conclusion}{Retour au début}\newline
\begin{minted}[obeytabs=true,tabsize=2,linenos=true, breaklines]{cpp}
#ifndef GAME_BOARD_HPP
#define GAME_BOARD_HPP

#include <vector>
#include <array>
#include <utility>
#include <optional>
#include <complex>
#include <cstddef>
#include <initializer_list>
#include <functional>
#include <CMatrix.hpp>
#include <Piece.hpp>
#include <TypePiece.hpp>
#include <Color.hpp>
#include <Coord.hpp>
#include <forward_list>
#include <observer_ptr.hpp>
#include <Move.hpp>
#include <Constexpr.hpp>
#include <check_path.hpp>

class Piece;

/**
 * @brief La classe représentant le plateau de jeu
\end{minted}
\end{frame}

\begin{frame}[fragile]
\hyperlink{conclusion}{Retour au début}\newline
\begin{minted}[obeytabs=true,tabsize=2,linenos=true,breaklines,firstnumber=27]{cpp}
 *
 * @tparam N Le nombre de ligne du plateau
 * @tparam M Le nombre de colonne du plateau
 */
template <std::size_t N = 8, std::size_t M = N>
class Board final
{
public:
    // Constructeur
    CONSTEXPR Board();
    /**
     * @brief Initialise le plateau à partir d'une liste de pièce,
     * les pièces ne sont pas divisées initialement.
     *
     * @param[in] board Double initializer_list sur des pointeurs sur Piece
     */
    CONSTEXPR Board(std::initializer_list<
                    std::initializer_list<
                        Piece>> const &
                        board);

    // Copie
    CONSTEXPR Board(Board const &) = default;
    CONSTEXPR Board &operator=(Board const &) = default;

    // Mouvement
    CONSTEXPR Board(Board &&) = default;
    CONSTEXPR Board &operator=(Board &&) = default;

    // Destructeur
\end{minted}
\end{frame}

\begin{frame}[fragile]
\hyperlink{conclusion}{Retour au début}\newline
\begin{minted}[obeytabs=true,tabsize=2,linenos=true,breaklines,firstnumber=57]{cpp}
    CONSTEXPR ~Board() = default;

    /**
     * @brief Retourne le nombre de ligne du plateau
     *
     * @return std::size_t Le nombre de ligne
     */
    CONSTEXPR static std::size_t numberLines() noexcept;

    /**
     * @brief Retourne le nombre de colonne
     *
     * @return std::size_t Le nombre de colonne
     */
    CONSTEXPR static std::size_t numberColumns() noexcept;

    /**
     * @brief Retourne un pointeur sur la piece à l'emplacement cible
     *
     * @param[in] n L'indice de la ligne
     * @param[in] m L'indice de la colonne
     * @return Un pointeur observateur sur une piece
     * ou nullptr si la case est vide
     */
    CONSTEXPR Piece const &
    operator()(std::size_t n,
               std::size_t m) const noexcept;

    /**
     * @brief Renvoie la liste dans tous les mouvements
\end{minted}
\end{frame}

\begin{frame}[fragile]
\hyperlink{conclusion}{Retour au début}\newline
\begin{minted}[obeytabs=true,tabsize=2,linenos=true,breaklines,firstnumber=87]{cpp}
     * classiques légaux d'une pièce
     *
     * @param pos Position de la pièce
     *
     * @return std::forward_list<Coord> La liste de coordonnées des positions
     * d'arrivées de chaque mouvement possible
     */
    CONSTEXPR std::forward_list<Coord>
    get_list_normal_move(Coord const &pos) const;

    /**
     * @brief Renvoie la liste de toutes les cases qu'une pièce
     * peut atteindre lors d'un mouvement split, il faut choisir
     * deux éléments de la liste pour créer un mouvement split.
     *
     * @param[in] pos Position de la pièce.
     *
     * @return std::forward_list<Coord> La liste de coordonnées
     * des positions d'arrivées possible sachant que chaque élément
     * représente une seule des deux cases d'arrivées d'un split move
     */
    CONSTEXPR std::forward_list<Coord>
    get_list_split_move(Coord const &pos) const;

    /**
     * @brief Teste si les mouvement de la piece sont si la
     * pièce est un pion un mouvement de promotion
     *
     * @note peut etre utiliser sans verifier si la pièce est un pion
     *
\end{minted}
\end{frame}

\begin{frame}[fragile]
\hyperlink{conclusion}{Retour au début}\newline
\begin{minted}[obeytabs=true,tabsize=2,linenos=true,breaklines,firstnumber=117]{cpp}
     * @param[in] pos La position de la pièce
     * @return true si le mouvement possible est un mouvement
     * de promotion
     * @return false sinon
     */
    CONSTEXPR bool
    check_if_use_move_promote(Coord const &pos) const noexcept;

    /**
     * @brief Renvoie la liste de toutes les promotions
     *
     * @warning Ne verifie pas la validité du mouvement,
     * de la position ou du type de la pièce
     *
     * @param[in] pos La position du pion
     * @return La liste de tout les mouvements de promotion
     */
    CONSTEXPR std::forward_list<Move>
    get_list_promote(Coord const &pos) const noexcept;

    /**
     * @brief Applique une fonction à l'entiéreté des mouvements possible
     *
     * @param[in, out] func La fonction à appliquer sur tout les mouvements
     * Prototype : bool f(Move const&)
     * Si renvoie true alors le parcourt est intérompue
     * @param[in] color La couleur du joueur auquel on
     * récupère les mouvements
     */
    template <class UnitaryFunction>
\end{minted}
\end{frame}

\begin{frame}[fragile]
\hyperlink{conclusion}{Retour au début}\newline
\begin{minted}[obeytabs=true,tabsize=2,linenos=true,breaklines,firstnumber=147]{cpp}
    CONSTEXPR void
    all_move(
        UnitaryFunction func,
        std::optional<Color> color_player = std::nullopt) const noexcept;

    /**
     * @brief Test si un mouvement est réalisable
     *
     * @param[in] move Le mouvement à tester
     * @return true si le mouvement est légal
     */
    CONSTEXPR bool move_is_legal(Move const &move) const;

    /**
     * @brief Réalise un mouvement d'une pièce quelque soit le type
     * du mouvement (classic, split, merge)
     *
     * @warning Ne procède aucune vérification sur la validité du mouvement
     *
     * @param[in] movement Le mouvement à raliser
     * @param[in] val_mes Optionel peut determiner la mesure de
     * la présence d'une pièce ou non
     */
    CONSTEXPR void
    move(Move const &movement, std::optional<bool> val_mes = std::nullopt);

    /**
     * @brief Test si un plateau est gagnant pour une couleur
     *
     * @param[in] c Couleur
\end{minted}
\end{frame}

\begin{frame}[fragile]
\hyperlink{conclusion}{Retour au début}\newline
\begin{minted}[obeytabs=true,tabsize=2,linenos=true,breaklines,firstnumber=177]{cpp}
     * @return true si la position est gagnante pour la couleur c
     * @return false sinon
     */
    CONSTEXPR bool winning_position(Color c) const noexcept;

    /**
     * @brief Recupère la couleur du joueur au tour de jouer
     *
     * @return Color la couleur du joueur
     */
    CONSTEXPR Color get_current_player() const noexcept;

    /**
     * @brief Change le joueur actuel et passe la main à l'autre joueur
     */
    CONSTEXPR void change_player() noexcept;

    /**
     * @brief Renvoie la probabilité qu'il y ait une pièce à une position.
     *
     * @param pos La position
     */
    CONSTEXPR double get_proba(Coord const &pos) const noexcept;

    friend class Piece;

    template <std::size_t _N, std::size_t _M>
    friend CONSTEXPR bool
    check_path_straight_1_instance(
        Board<_N, _M> const &board,
\end{minted}
\end{frame}

\begin{frame}[fragile]
\hyperlink{conclusion}{Retour au début}\newline
\begin{minted}[obeytabs=true,tabsize=2,linenos=true,breaklines,firstnumber=207]{cpp}
        Coord const &dpt,
        Coord const &arv,
        std::size_t position,
        std::optional<Coord>);

    template <std::size_t _N, std::size_t _M>
    friend CONSTEXPR bool
    check_path_diagonal_1_instance(
        Board<_N, _M> const &board,
        Coord const &dpt,
        Coord const &arv,
        std::size_t position,
        std::optional<Coord>);

    template <std::size_t _N, std::size_t _M>
    friend CONSTEXPR bool
    check_path_queen_1_instance(
        Board<_N, _M> const &board,
        Coord const &dpt,
        Coord const &arv,
        std::size_t position,
        std::optional<Coord>);

    /**
     * @brief Fonction qui permet de faire un mouvement de pion quelconque avec une promotion
     *
     * @param s Coordonnées de la source
     * @param t Coordonnées de la cible
     * @param p Le type de la pièce de la promotion
     * @param[in] val_mes Optionel peut determiner la mesure de
\end{minted}
\end{frame}

\begin{frame}[fragile]
\hyperlink{conclusion}{Retour au début}\newline
\begin{minted}[obeytabs=true,tabsize=2,linenos=true,breaklines,firstnumber=237]{cpp}
     * la présence d'une pièce ou non
     */
    CONSTEXPR void move_promotion(
        Move const &move,
        std::optional<bool> val_mes = std::nullopt);

    /**
     * @brief Renvoie la probabilité que le mouvement soit réalisé
     * 
     * @param move Le mouvement à réaliser
     * @return Une probabilité comprise entre 0 et 1
     */
    CONSTEXPR double get_proba_move(Move const &move) const noexcept;

private:


    /**
     * @brief Renvoie la probabilité que la fonction mesure renvoie true,
     * c'est-à-dire la proba que le mouvement soit"réussi"
     *
     * @param p Coordonnées de la case
     * @return La probabilité
     */
    CONSTEXPR double
    get_proba_mesure(Coord const &p) const noexcept;

    /**
     * @brief Renvoie la probabilité que la fonction mesure_capture_slide renvoie true,
     * c'est-à-dire la proba que le mouvement de capture_slide soit"réussi"
\end{minted}
\end{frame}

\begin{frame}[fragile]
\hyperlink{conclusion}{Retour au début}\newline
\begin{minted}[obeytabs=true,tabsize=2,linenos=true,breaklines,firstnumber=267]{cpp}
     *
     * @param s Coordonnées de la source
     * @param t Coordonnées de la cible
     * @param check_path Fonction qui teste la présence d'une pièce 
     * entre une source et une cible
     * @return La probabilité
     */
    CONSTEXPR double get_proba_mesure_capture_slide(
    Coord const &s,
    Coord const &t,
    std::function<bool(Board<N, M> const &,
                       Coord const &,
                       Coord const &,
                       std::size_t,
                       std::optional<Coord>)>
        check_path) const noexcept;

    /**
     * @brief Renvoie laprobabilité que la fonction mesure_castle renvoie true,
     * c'est-à-dire la proba que le mouvement de roque soit"réussi"
     *
     * @param king Coordonnées du roi
     * @param rook Coordonnées de la tour
     * @return La probabilité
     */
    CONSTEXPR double get_proba_mesure_castle(
        Coord const &king,
        Coord const &rook) const noexcept;

    /**
\end{minted}
\end{frame}

\begin{frame}[fragile]
\hyperlink{conclusion}{Retour au début}\newline
\begin{minted}[obeytabs=true,tabsize=2,linenos=true,breaklines,firstnumber=297]{cpp}
     * @brief Vérifie si la case à une possibilité de contenir une pièce,
     * et si elle n'en a pas modifie m_piece_board en nullptr.
     *
     * @param pos Les coordonnées de a position de la case a vérifier.
     */
    void update_case(std::size_t pos) noexcept;
    /**
     * @brief Fonction qui permet de mattre à jour le plateau,
     * car ce mouvement entraine l'apparition de plusieurs instance de board
     * identique, il faut donc les concaténer en ajoutant les probas, si la
     * proba vaut 0, on supprime l'instance
     *
     * @warning Complexité en k² x N x M, avec k la taille de m_board
     * c'est à dire le nombre d'instance du plateau, N et M les dimensions
     * du plateau
     */
    void update_board() noexcept;
    void update_board_classic() noexcept;

    /**
     * @brief Mouvement classique d'une pièce qui "saute"
     * (cavalier, roi)
     *
     * @warning Ne procède aucune vérification sur la validité du mouvement
     *
     * @tparam N Le nombre de ligne du plateau
     * @tparam M Le nombre de colonnes du plateau
     * @param[in] s Coordonnées de la source du mouvement
     * @param[in] t Coordonnées de la cible du mouvement
     * @param[in] val_mes Optionel peut determiner la mesure de
\end{minted}
\end{frame}

\begin{frame}[fragile]
\hyperlink{conclusion}{Retour au début}\newline
\begin{minted}[obeytabs=true,tabsize=2,linenos=true,breaklines,firstnumber=327]{cpp}
     * la présence d'une pièce ou non
     */
    CONSTEXPR void move_classic_jump(
        Coord const &s,
        Coord const &t,
        std::optional<bool> val_mes = std::nullopt);

    /**
     * @brief Mouvement "split jump"
     *
     * @warning Aucun tests sur la validité du mouvement (cible vide, ect)
     *
     * @tparam N Le nombre de lignes du plateau
     * @tparam M Le nombre de colonnes du plateau
     * @param[in] s Coordonnées de la source du mouvement
     * @param[in] t1 Coordonnées de la cible 1 (qui doit être vide)
     * @param[in] t2 Coordonnées de la cible 2 (qui doit être vide)
     */
    CONSTEXPR void move_split_jump(
        Coord const &s,
        Coord const &t1,
        Coord const &t2);

    /**
     * @brief Mouvement de merge
     *
     * @warning Aucun test sur la validité du mouvement
     * (cible vide, pièce identique sur les sources, ect)
     *
     * @tparam N Le nombre de lignes du plateau
\end{minted}
\end{frame}

\begin{frame}[fragile]
\hyperlink{conclusion}{Retour au début}\newline
\begin{minted}[obeytabs=true,tabsize=2,linenos=true,breaklines,firstnumber=357]{cpp}
     * @tparam M Le nombre de colonnes du plateau
     * @param[in] s1 Coordonnées de la source 1
     * @param[in] s2 Coordonnées de la source 2
     * @param[in] t Coordonnées de la cible
     */
    CONSTEXPR void move_merge_jump(
        Coord const &s,
        Coord const &t1,
        Coord const &t2);
    /**
     * @brief Mouvement classique d'un pièce qui "glisse" (fou, dame, tour)
     *
     * @warning Ne procède aucune vérification sur la validité du mouvement
     *
     * @tparam N Le nombre de lignes
     * @tparam M Le nombre de colonnes
     * @param[in] s Coordonnées de la source
     * @param[in] t Coordonnées de la cible
     * @param[in] check_path  Fonction qui permet de vérifier la présence
     * d'une pièce entre la source et la cible sur une instance du plateau
     * @param[in] val_mes Optionel peut determiner la mesure de
     * la présence d'une pièce ou non
     */
    CONSTEXPR void
    move_classic_slide(Coord const &s, Coord const &t,
                       std::function<
                           bool(
                               Board<N, M> const &,
                               Coord const &,
                               Coord const &,
\end{minted}
\end{frame}

\begin{frame}[fragile]
\hyperlink{conclusion}{Retour au début}\newline
\begin{minted}[obeytabs=true,tabsize=2,linenos=true,breaklines,firstnumber=387]{cpp}
                               std::size_t,
                               std::optional<Coord>)>
                           check_path,
                       std::optional<bool> val_mes = std::nullopt);
    /**
     * @brief Mouvement "split slide"
     * @warning Aucun tests sur la validité du mouvement (cible vide, ect)
     * @tparam N Le nombre de lignes du plateau
     * @tparam M Le nombre de colonnes du plateau
     * @param[in] s Coordonnées de la source du mouvement
     * @param[in] t1 Coordonnées de la cible 1 (qui doit être vide)
     * @param[in] t2 Coordonnées de la cible 2 (qui doit être vide)
     * @param[in] check_path Fonction qui permet de vérifier la présence d'une
     * pièce entre la source et la cible sur une instance du plateau
     */
    CONSTEXPR void move_split_slide(
        Coord const &s,
        Coord const &t1,
        Coord const &t2,
        std::function<
            bool(
                Board<N, M> const &,
                Coord const &,
                Coord const &,
                std::size_t,
                std::optional<Coord>)
                >
            check_path);

    /**
\end{minted}
\end{frame}

\begin{frame}[fragile]
\hyperlink{conclusion}{Retour au début}\newline
\begin{minted}[obeytabs=true,tabsize=2,linenos=true,breaklines,firstnumber=417]{cpp}
     * @brief Mouvement de merge
     * @warning Aucun test sur la validité du mouvement
     * (cible vide, pièce identique sur les sources, ect)
     * @tparam N Le nombre de lignes du plateau
     * @tparam M Le nombre de colonnes du plateau
     * @param[in] s1 Coordonnées de la source 1
     * @param[in] s2 Coordonnées de la source 2
     * @param[in] t Coordonnées de la cible
     * @param[in] check_path Fonction qui permet de vérifier la présence d'une
     * pièce entre la source et la cible sur une instance du plateau
     */
    CONSTEXPR void move_merge_slide(
        Coord const &s1,
        Coord const &s2,
        Coord const &t,
        std::function<
            bool(
                Board<N, M> const &,
                Coord const &,
                Coord const &,
                std::size_t,
                std::optional<Coord>)>
            check_path);

    /**
     * @brief Le mouvement de piond'une case fonctionne comme
     * un mouvementde jump classique, à la différence qu'un
     * pion ne peut pas capturer une pièce en avançant.
     * On mesure de la même façon qu'un jump classique en considérant
     * toutes les pièces comme des pièces alliées
\end{minted}
\end{frame}

\begin{frame}[fragile]
\hyperlink{conclusion}{Retour au début}\newline
\begin{minted}[obeytabs=true,tabsize=2,linenos=true,breaklines,firstnumber=447]{cpp}
     *
     * @warning Ne procède aucune vérification sur la validité du mouvement
     *
     * @tparam N Le nombre de lignes du plateau
     * @tparam M Le nombre de colonnes du plateau
     * @param[in] s Coordonnées de la source
     * @param[in] t Coordonnées de la cible
     * @param[in] val_mes Optionel peut determiner la mesure de
     * la présence d'une pièce ou non
     * @return true si le mouvement à etait effectué
     * @return false sinon
     */
    CONSTEXPR bool
    move_pawn_one_step(Coord const &s, Coord const &t,
                       std::optional<bool> val_mes = std::nullopt);

    /**
     * @brief Le mouvement de pion de deux cases fonctionne
     * comme un mouvement de slide classique, à la différence
     * qu'un pion ne peut pas capturer une pièce en avançant,
     * on mesure de la même façon qu'un slide classique en considérant
     * toutes les pièces comme des pièces alliées.
     *
     * @warning Aucun test sur la possibilité de faire un mouvement de deux
     *  cases du pion, pas de mise a jour du board sur la prise en passant
     *
     * @tparam N Le nombre de lignes du plateau
     * @tparam M Le nombre de colonnes du plateau
     * @param[in] s Coordonnées de la source
     * @param[in] t Coordonnées de la cible
\end{minted}
\end{frame}

\begin{frame}[fragile]
\hyperlink{conclusion}{Retour au début}\newline
\begin{minted}[obeytabs=true,tabsize=2,linenos=true,breaklines,firstnumber=477]{cpp}
     * @param[in] val_mes Optionel peut determiner la mesure de
     * la présence d'une pièce ou non
     * @return true si le mouvement à etait effectué
     * @return false sinon
     */
    CONSTEXPR bool
    move_pawn_two_step(Coord const &s, Coord const &t,
                       std::optional<bool> val_mes = std::nullopt);
    /**
     * @brief Mouvement de capture dun pion. A la différence
     * d'un mouvement de "capture jump", on mesure la présence
     * de la cible car le pion a besoin de la cible pour se déplacer
     *
     * @warning Ne procède aucune vérification sur la validité du mouvement
     *
     * @param[in] s Coordonnées de la source
     * @param[in] t Coordonnées de la cible
     * @param[in] val_mes Optionel peut determiner la mesure de
     * la présence d'une pièce ou non
     * @return true si le mouvement à etait effectué
     * @return false sinon
     */
    CONSTEXPR bool capture_pawn(Coord const &s,
                                Coord const &t,
                                std::optional<bool> val_mes = std::nullopt);

    /**
     * @brief Permet d'effectuer un mouvement de prise en passant
     *
     * @warning Aucun test sur la validité de la cible et de la cible
\end{minted}
\end{frame}

\begin{frame}[fragile]
\hyperlink{conclusion}{Retour au début}\newline
\begin{minted}[obeytabs=true,tabsize=2,linenos=true,breaklines,firstnumber=507]{cpp}
     * en passant, ni sur le type de la pièce
     *
     * @tparam N Le nombre de lignes du plateau
     * @tparam M Le nombre de colonnnes du plateau
     * @param[in] s Les coordonnées de la source
     * @param[in] t Les coordonnées de la cible (l'endroit où arrive le pion)
     * @param[in] ep Les coordonnées du pion capturer "en passant"
     * @param[in] val_mes Optionel peut determiner la mesure de
     * la présence d'une pièce ou non
     * @return true si le mouvement à etait effectué
     * @return false sinon
     */
    CONSTEXPR bool
    move_enpassant(Coord const &s, Coord const &t, Coord const &ep,
                   std::optional<bool> val_mes = std::nullopt);

    /**
     * @brief Mesure la présence d'une pièce
     *
     * @tparam N Le nombre de lignes du plateau
     * @tparam M Le nombre de colonnes du plateau
     * @param[in] position La position de la pièce mesurée
     * @param[in] val_mes Optionel peut determiner la mesure de
     * la présence d'une pièce ou non
     * @return true Si la pièce est présente sur la case position
     * @return false Sinon
     */
    CONSTEXPR bool mesure(Coord const &p,
                          std::optional<bool> val_mes = std::nullopt);

\end{minted}
\end{frame}

\begin{frame}[fragile]
\hyperlink{conclusion}{Retour au début}\newline
\begin{minted}[obeytabs=true,tabsize=2,linenos=true,breaklines,firstnumber=537]{cpp}
    /**
     * @brief Fonction qui permet de faire une mesure dans le cas
     d'un mouvement "capture slide"
     *
     * @tparam N Le nombre de lignes du plateau
     * @tparam M Le nombre de colonnes du plateau
     * @param[in] s Coordonnées de la source du mouvement
     * @param[in] t Coordonnées de la cible du mouvement
     * @param[in] check_path Fonction qui permet de vérifier si il y a une pièce
     * entre la source et la cible sur une instance du plateau
     * @param[in] val_mes Optionel peut determiner la mesure de
     * la présence d'une pièce ou non
     * @return true Si la mesure indique de faire le mouvement
     * @return false Sinon
     */
    CONSTEXPR bool
    mesure_capture_slide(Coord const &s, Coord const &t,
                         std::function<bool(Board<N, M> const &,
                                            Coord const &, Coord const &,
                                            std::size_t,
                                            std::optional<Coord>)>
                             check_path,
                         std::optional<bool> val_mes = std::nullopt);

    /**
     * @brief Mouvement classique d'une pièce (pion exclu)
     * @warning Ne procède aucune vérification sur la validité du mouvement
     * @tparam N Le nombre de ligne du plateau
     * @tparam M Le nombre de colonnes du plateau
     * @param[in] s Coordonnées de la source du mouvement
\end{minted}
\end{frame}

\begin{frame}[fragile]
\hyperlink{conclusion}{Retour au début}\newline
\begin{minted}[obeytabs=true,tabsize=2,linenos=true,breaklines,firstnumber=567]{cpp}
     * @param[in] t Coordonnées de la cible du mouvement
     */
    CONSTEXPR void move_classic(Coord const &s, Coord const &t,
                                std::optional<bool> val_mes);

    /**
     * @brief Mouvement split d'une pièce (pion exclu)
     * @warning Ne procède aucune vérification sur la validité du mouvement
     * @tparam N Le nombre de ligne du plateau
     * @tparam M Le nombre de colonnes du plateau
     * @param[in] s Coordonnées de la source du mouvement
     * @param[in] t1 Coordonnées de la cible 1 (qui doit être vide)
     * @param[in] t2 Coordonnées de la cible 2 (qui doit être vide)
     */
    CONSTEXPR void move_split(Coord const &s,
                              Coord const &t1,
                              Coord const &t2);

    /**
     * @brief Mouvement de merge de pièce (pion exclu)
     *
     * @warning Aucun test sur la validité du mouvement
     * (cible vide, pièce identique sur les sources, ect)
     *
     * @tparam N Le nombre de lignes du plateau
     * @tparam M Le nombre de colonnes du plateau
     * @param[in] s1 Coordonnées de la source 1
     * @param[in] s2 Coordonnées de la source 2
     * @param[in] t Coordonnées de la cible
     */
\end{minted}
\end{frame}

\begin{frame}[fragile]
\hyperlink{conclusion}{Retour au début}\newline
\begin{minted}[obeytabs=true,tabsize=2,linenos=true,breaklines,firstnumber=597]{cpp}
    CONSTEXPR void move_merge(Coord const &s1,
                              Coord const &s2,
                              Coord const &t);

    /**
     * @brief Gère tout les mouvements d'un pion
     * @warning Ne procède aucune vérification sur la validité du mouvement
     * @tparam N Le nombre de lignes du plateau
     * @tparam M Le nombre de colonnes du plateau
     * @param[in] s Coordonnées de la source
     * @param[in] t Coordonnées de la cible
     * @param[in] val_mes Optionel peut determiner la mesure de
     * la présence d'une pièce ou non
     * @return true si le mouvement à etait effectué
     * @return false sinon
     */
    CONSTEXPR bool move_pawn(
        Coord const &s,
        Coord const &t,
        std::optional<bool> val_mes = std::nullopt) noexcept;

    /**
     * @brief Renvoie une position 1D d'une coordonnée 2D
     *
     * @param[in] ligne L'indice de ligne
     * @param[in] colonne L'indice de colonne
     * @return std::size_t La position en 1D
     */
    CONSTEXPR static std::size_t
    offset(std::size_t ligne,
\end{minted}
\end{frame}

\begin{frame}[fragile]
\hyperlink{conclusion}{Retour au début}\newline
\begin{minted}[obeytabs=true,tabsize=2,linenos=true,breaklines,firstnumber=627]{cpp}
           std::size_t colonne) noexcept;

    /**
     * @brief Initialise un plateau à l'aide des listes d'initialisations
     *
     * @param[in] board La liste d'initialisation en 2D
     * @param[out] first_instance Le tableau de la première
     * instance du plateau
     * @param[out] piece_board Le plateau contenant les pièces
     */
    CONSTEXPR static void
    initializer_list_to_2_array(
        std::initializer_list<
            std::initializer_list<
                Piece>> const &board,
        std::array<bool, N * M> &first_instance,
        std::array<Piece, N * M> &piece_board) noexcept;

    /**
     * @brief Construit les mailbox en fonction
     * des dimenssions du plateau
     *
     * @param[out] S_mailbox La petite mailbox de la taille du plateau
     * @param[out] L_mailbox La grande mailbox comportant 4 ligne de plus
     * et 2 colonne de plus
     * @return CONSTEXPR
     */
    CONSTEXPR static void init_mailbox(std::array<int, N * M>
                                           &S_mailbox,
                                       std::array<int, (N + 4) * (M + 2)>
\end{minted}
\end{frame}

\begin{frame}[fragile]
\hyperlink{conclusion}{Retour au début}\newline
\begin{minted}[obeytabs=true,tabsize=2,linenos=true,breaklines,firstnumber=657]{cpp}
                                           &L_mailbox) noexcept;

    /**
     * @brief fonction auxiliaire qui permet de modifier un plateau
     à l'aide d'un array de la forme du type de retour de la fonction
     qubitToArray, le deuxième éléments du tableau à une probalité non nulle
     lors d'un mouvement split
     *
     * @tparam Q La taille du qubit
     * @tparam N Le nombre de ligne du plateau
     * @tparam M Le nombre de collone du plateau
     * @param[in] arrayQubit Type de retour de la fonction quibitToArray,
     * représente la facon dont va être modifier m_board
     * @param[in] position_board  L'indice du plateau dans le tableau de
     * toutes les instances du plateau
     * @param[in] tab_positions Tableau des indices des cases modifiées,
     * on utilise N*M+1 pour signifier que le qubit en question
     * est un qubit auxiliaire qui ne modifie pas le board
     */
    template <std::size_t Q>
    CONSTEXPR void modify(std::array<std::pair<std::array<bool, Q>,
                                               std::complex<double>>,
                                     2> const &arrayQubit,
                          std::size_t position_board,
                          std::array<std::size_t, Q> const &tab_positions);

    /**
     * @brief Fonction qui permet d'effectuer un mouvement
     * sur une instance du plateau
     *
\end{minted}
\end{frame}

\begin{frame}[fragile]
\hyperlink{conclusion}{Retour au début}\newline
\begin{minted}[obeytabs=true,tabsize=2,linenos=true,breaklines,firstnumber=687]{cpp}
     * @tparam N Le nombre de ligne du plateau
     * @tparam M Le nombre de colonnes du plateau
     * @tparam Q La taille du qubit
     * @param[in] case_modif Permet d'initialiser le qubit
     * @param[in] position L'instance du plateau modifiée
     * @param[in] matrix La matrice du mouvement que l'on veut effectuer
     * @param[in] tab_positions La position sur le plateau des variables
     * utilisées dans le qubit, que l'on met à N*M+1 si les variables
     * ne représentent pas des pièces
     */
    template <std::size_t Q>
    CONSTEXPR void move_1_instance(std::array<bool, Q> const &case_modif,
                                   std::size_t position,
                                   CMatrix<_2POW(Q)> const &matrix,
                                   std::array<std::size_t, Q> const
                                       &tab_positions) noexcept;

    /**
     * @brief Mouvement du petit roque.
     * @warning Auncun test sur la validité du mouvement
     * @param[in] king Coordonnées du roi
     * @param[in] rook Coordonnées de la tour
     * @param[in] val_mes Optionel peut determiner la mesure de
     * la présence d'une pièce ou non
     */
    CONSTEXPR void king_side_castle(
        Coord const &king,
        Coord const &rook,
        std::optional<bool> val_mes = std::nullopt);

\end{minted}
\end{frame}

\begin{frame}[fragile]
\hyperlink{conclusion}{Retour au début}\newline
\begin{minted}[obeytabs=true,tabsize=2,linenos=true,breaklines,firstnumber=717]{cpp}
    /**
     * @brief Fonction qui permet de mesurer la possibilité de faire le roque
     *
     * @param[in] king Coordonnées du roi
     * @param[in] rook Coordonnées de la tour
     * @param[in] val_mes Optionel peut determiner la mesure de
     * la présence d'une pièce ou non
     * @return true si le roque est possible
     * @return false sinon
     */
    CONSTEXPR bool mesure_castle(
        Coord const &king,
        Coord const &rook,
        std::optional<bool> val_mes = std::nullopt);

    /**
     * @brief Mouvement du grand roque.
     *
     * @warning Ne procède aucune vérification sur la validité du mouvement
     *
     * @param[in] king Coordonnées du roi
     * @param[in] rook Coordonnées de la tour
     * @param[in] val_mes Optionel peut determiner la mesure de
     * la présence d'une pièce ou non
     */
    CONSTEXPR void queen_side_castle(
        Coord const &king,
        Coord const &rook,
        std::optional<bool> val_mes = std::nullopt);

\end{minted}
\end{frame}

\begin{frame}[fragile]
\hyperlink{conclusion}{Retour au début}\newline
\begin{minted}[obeytabs=true,tabsize=2,linenos=true,breaklines,firstnumber=747]{cpp}
    /**
     * @brief Le tableau de toutes les instances possibles
     * du plateau
     */
    std::vector<std::pair<std::array<bool, N * M>,
                          std::complex<double>>>
        m_board;

    /**
     * @brief Le plateau indiquant le type des pièces
     */
    std::array<Piece, N * M> m_piece_board;

    /**
     * @brief La petite mailbox
     */
    std::array<int, N * M> m_S_mailbox;

    /**
     * @brief La grande mailbox
     */
    std::array<int, (N + 4) * (M + 2)> m_L_mailbox;

    /**
     * @brief Color::WHITE si c'est aux blanc de jouer aux noirs sinon
     */
    Color m_color_current_player;

    /**
     * @brief Vrai si les noirs[0] / blancs[1] peuvent faire le petit roque
\end{minted}
\end{frame}

\begin{frame}[fragile]
\hyperlink{conclusion}{Retour au début}\newline
\begin{minted}[obeytabs=true,tabsize=2,linenos=true,breaklines,firstnumber=777]{cpp}
     */
    std::array<bool, 2> m_k_castle;

    /**
     * @brief Vrai si les noirs[0] / blancs[1] peuvent faire le grand roque
     */
    std::array<bool, 2> m_q_castle;

    /**
     * @brief Contient la case sur laquelle il est
     * possible de faire une prise en passant qui
     * est vide (la case occupé si le pion avait avancé
     * d'une seule case)
     */
    std::optional<Coord> m_ep;
};

#include "Board.tpp"
#include "mesure.tpp"
#include "get_proba_move.tpp"
#include "move.tpp"


\end{minted}
\end{frame}

\begin{frame}[fragile]
\label{Board.tpp}
\frametitle{Board.tpp}
\hyperlink{conclusion}{Retour au début}\newline
\begin{minted}[obeytabs=true,tabsize=2,linenos=true, breaklines]{cpp}
// Lib standard
#include <array>
#include <initializer_list>
#include <algorithm>
#include <cassert>
#include <utility>
#include <functional>
#include <cmath>
#include <optional>
#include <stdexcept>

// Inclusion projet
#include <Qubit.hpp>
#include <Piece.hpp>
#include <CMatrix.hpp>
#include <Unitary.hpp>
#include <TypePiece.hpp>
#include <math_utility.hpp>
#include <Constexpr.hpp>
#include <Move.hpp>
#include <Random.hpp>
#include "Board.hpp"

template <std::size_t N, std::size_t M>
CONSTEXPR Board<N, M>::Board()
    : m_board(),
\end{minted}
\end{frame}

\begin{frame}[fragile]
\hyperlink{conclusion}{Retour au début}\newline
\begin{minted}[obeytabs=true,tabsize=2,linenos=true,breaklines,firstnumber=27]{cpp}
      m_piece_board(),
      m_S_mailbox(),
      m_L_mailbox(),
      m_color_current_player(Color::WHITE),
      m_k_castle({true, true}),
      m_q_castle({true, true}),
      m_ep()
{
    m_board.push_back(std::pair<std::array<bool, 64>,
                                std::complex<double>>{{}, 0});
    init_mailbox(m_S_mailbox, m_L_mailbox);
};

template <std::size_t N, std::size_t M>
CONSTEXPR Board<N, M>::Board(std::initializer_list<
                             std::initializer_list<
                                 Piece>> const &board)
    : m_board(),
      m_piece_board(),
      m_S_mailbox(),
      m_L_mailbox(),
      m_color_current_player(Color::WHITE),
      m_k_castle({true, true}),
      m_q_castle({true, true}),
      m_ep()
{
    init_mailbox(m_S_mailbox, m_L_mailbox);
    m_board.push_back(std::pair<std::array<bool, N * M>,
                                std::complex<double>>{{false}, 1.});
    initializer_list_to_2_array(board, m_board[0].first, m_piece_board);
\end{minted}
\end{frame}

\begin{frame}[fragile]
\hyperlink{conclusion}{Retour au début}\newline
\begin{minted}[obeytabs=true,tabsize=2,linenos=true,breaklines,firstnumber=57]{cpp}
};

template <std::size_t N, std::size_t M>
CONSTEXPR std::size_t
Board<N, M>::offset(std::size_t ligne, std::size_t colonne) noexcept
{
    return ligne * M + colonne;
};

template <std::size_t N, std::size_t M>
CONSTEXPR std::size_t Board<N, M>::numberLines() noexcept
{
    return N;
};

template <std::size_t N, std::size_t M>
CONSTEXPR std::size_t Board<N, M>::numberColumns() noexcept
{
    return M;
};

template <std::size_t N, std::size_t M>
CONSTEXPR void
Board<N, M>::initializer_list_to_2_array(
    std::initializer_list<
        std::initializer_list<
            Piece>> const &board,
    std::array<bool, N * M> &first_instance,
    std::array<Piece, N * M> &piece_board) noexcept
{
\end{minted}
\end{frame}

\begin{frame}[fragile]
\hyperlink{conclusion}{Retour au début}\newline
\begin{minted}[obeytabs=true,tabsize=2,linenos=true,breaklines,firstnumber=87]{cpp}

    auto it_tab{std::begin(first_instance)};
    auto it_piece_dst{std::begin(piece_board)};
    assert(std::size(board) <= N && "Value entry out of Board");
    for (std::initializer_list<Piece> const &e : board)
    {
        auto const it_tab_begin_line{it_tab};
        assert(std::size(e) <= M && "Value entry out of Board");
        std::for_each(std::begin(e),
                      std::end(e),
                      [it_tab](Piece piece) mutable
                      -> void
                      {
                        if (piece.get_type() != TypePiece::EMPTY)
                        {
                            *it_tab = true;
                        }
                        else
                        {
                            *it_tab = false;
                        }
                        it_tab++; });
        std::move(std::begin(e), std::end(e), it_piece_dst);
        it_tab = it_tab_begin_line + M;
        it_piece_dst += M;
    }
}

template <std::size_t N, std::size_t M>
CONSTEXPR void
\end{minted}
\end{frame}

\begin{frame}[fragile]
\hyperlink{conclusion}{Retour au début}\newline
\begin{minted}[obeytabs=true,tabsize=2,linenos=true,breaklines,firstnumber=117]{cpp}
Board<N, M>::init_mailbox(
    std::array<int, N * M> &S_mailbox,
    std::array<int, (N + 4) * (M + 2)> &L_mailbox) noexcept
{
    std::fill(std::begin(L_mailbox),
              std::begin(L_mailbox) + (2 * (M + 2) + 1), -1);

    std::fill(std::end(L_mailbox) - (2 * (M + 2) + 1),
              std::end(L_mailbox), -1);
    auto offset_L_mailboard{
        [](std::size_t n, std::size_t m)
            -> std::size_t
        {
            return n * (M + 2) + m;
        }};
    for (std::size_t i{0}; i < N; i++)
    {
        L_mailbox[offset_L_mailboard(i + 2, 0)] = -1;
        L_mailbox[offset_L_mailboard(i + 2, M + 1)] = -1;
        for (std::size_t j{0}; j < M; j++)
        {
            L_mailbox[offset_L_mailboard(i + 2, j + 1)] =
                static_cast<int>(Board<N, M>::offset(i, j));
            S_mailbox[Board<N, M>::offset(i, j)] =
                static_cast<int>(offset_L_mailboard(i + 2, j + 1));
        }
    }
}

template <std::size_t N, std::size_t M>
\end{minted}
\end{frame}

\begin{frame}[fragile]
\hyperlink{conclusion}{Retour au début}\newline
\begin{minted}[obeytabs=true,tabsize=2,linenos=true,breaklines,firstnumber=147]{cpp}
CONSTEXPR double Board<N, M>::get_proba(Coord const &pos) const noexcept
{
    double acc{0.};
    for (auto const &e : m_board)
    {
        acc += pow(abs(e.second), 2.) * e.first[offset(pos.n, pos.m)];
    }
    return acc;
}

template <std::size_t N, std::size_t M>
CONSTEXPR std::forward_list<Coord>
Board<N, M>::get_list_normal_move(Coord const &pos) const
{
    if ((*this)(pos.n, pos.m).get_type() != TypePiece::EMPTY)
    {
        return (*this)(pos.n, pos.m).get_list_normal_move(*this, pos);
    }
    return std::forward_list<Coord>{};
}

template <std::size_t N, std::size_t M>
CONSTEXPR std::forward_list<Coord>
Board<N, M>::get_list_split_move(Coord const &pos) const
{
    if ((*this)(pos.n, pos.m).get_type() != TypePiece::EMPTY)
    {
        return (*this)(pos.n, pos.m).get_list_split_move(*this, pos);
    }
    return std::forward_list<Coord>{};
\end{minted}
\end{frame}

\begin{frame}[fragile]
\hyperlink{conclusion}{Retour au début}\newline
\begin{minted}[obeytabs=true,tabsize=2,linenos=true,breaklines,firstnumber=177]{cpp}
}

template <std::size_t N, std::size_t M>
CONSTEXPR Color Board<N, M>::get_current_player() const noexcept
{
    return m_color_current_player;
}

template <std::size_t N, std::size_t M>
CONSTEXPR void Board<N, M>::change_player() noexcept
{
    if (get_current_player() == Color::WHITE)
    {
        m_color_current_player = Color::BLACK;
    }
    else
    {
        m_color_current_player = Color::WHITE;
    }
}

template <std::size_t N, std::size_t M>
template <std::size_t Q>
CONSTEXPR void Board<N, M>::modify(
    std::array<
        std::pair<
            std::array<bool, Q>,
            std::complex<double>>,
        2> const &arrayQubit,
    std::size_t position_board,
\end{minted}
\end{frame}

\begin{frame}[fragile]
\hyperlink{conclusion}{Retour au début}\newline
\begin{minted}[obeytabs=true,tabsize=2,linenos=true,breaklines,firstnumber=207]{cpp}
    std::array<std::size_t, Q> const &tab_positions)

{
    if (!complex_equal(arrayQubit[1].second, 0i))
    {
        std::pair<std::array<bool, N * M>, std::complex<double>> new_b{};
        std::copy(std::begin(m_board[position_board].first),
                  std::end(m_board[position_board].first),
                  std::begin(new_b.first));
        new_b.second = m_board[position_board].second;
        for (std::size_t i{0}; i < Q; i++)
        {
            if (tab_positions[i] < N * M + 1)
            { /*Ca nous permet de mettre des variables
                dans les qubits qui ne sont pas prises
                en compte lors de la modif du plateau */
                new_b.first[tab_positions[i]] = arrayQubit[1].first[i];
            }
        }
        new_b.second *= arrayQubit[1].second;
        m_board.push_back(new_b);
    }
    for (std::size_t i{0}; i < Q; i++)
    {
        if (tab_positions[i] < N * M + 1)
        {
            m_board[position_board].first[tab_positions[i]] =
                arrayQubit[0].first[i];
        }
    }
\end{minted}
\end{frame}

\begin{frame}[fragile]
\hyperlink{conclusion}{Retour au début}\newline
\begin{minted}[obeytabs=true,tabsize=2,linenos=true,breaklines,firstnumber=237]{cpp}
    m_board[position_board].second *= arrayQubit[0].second;
}

template <std::size_t N, std::size_t M>
template <std::size_t Q>
CONSTEXPR void
Board<N, M>::move_1_instance(std::array<bool, Q> const &case_modif,
                             std::size_t position,
                             CMatrix<_2POW(Q)> const &matrix,
                             std::array<std::size_t, Q> const
                                 &tab_positions) noexcept
{
    Qubit<Q> q{case_modif};
    auto x{qubitToArray(matrix * q)};
    modify(std::move(x), position, tab_positions);
}

template <std::size_t N, std::size_t M>
CONSTEXPR Piece const &
Board<N, M>::operator()(std::size_t n, std::size_t m) const noexcept
{
    return m_piece_board[offset(n, m)];
}

template <std::size_t N, std::size_t M>
void Board<N, M>::update_board() noexcept
{
    std::size_t size_board = std::size(m_board);
    for (std::size_t i{size_board}; i > 0; i--)
    {
\end{minted}
\end{frame}

\begin{frame}[fragile]
\hyperlink{conclusion}{Retour au début}\newline
\begin{minted}[obeytabs=true,tabsize=2,linenos=true,breaklines,firstnumber=267]{cpp}
        using namespace std::complex_literals;
        if (complex_equal(m_board[i - 1].second, 0i))
        {
            m_board.erase(std::begin(m_board) + i - 1);
        }
        else
        {
            for (std::size_t j{i - 1}; j > 0; j--)
            {
                if (m_board[i - 1].first == m_board[j - 1].first)
                {
                    m_board[j - 1].second += m_board[i - 1].second;
                    m_board.erase(std::begin(m_board) + i - 1);
                    break;
                }
            }
        }
    }
}
template <std::size_t N, std::size_t M>
CONSTEXPR bool Board<N, M>::winning_position(Color c) const noexcept
{
    for (std::size_t i{0}; i < N * M; i++)
    {
        if (m_piece_board[i].get_type() == TypePiece::KING &&
            m_piece_board[i].get_color() != c)
        {
            return false;
        }
    }
\end{minted}
\end{frame}

\begin{frame}[fragile]
\hyperlink{conclusion}{Retour au début}\newline
\begin{minted}[obeytabs=true,tabsize=2,linenos=true,breaklines,firstnumber=297]{cpp}
    return true;
}

template <std::size_t N, std::size_t M>
CONSTEXPR bool
Board<N, M>::move_is_legal(Move const &move) const
{
    if (move.type == TypeMove::NORMAL)
    {
        auto list{get_list_normal_move(move.normal.src)};
        return std::find(
                   std::begin(list),
                   std::end(list),
                   move.normal.arv) !=
               std::end(list);
    }
    else if (move.type == TypeMove::SPLIT)
    {
        std::forward_list<Coord> list{get_list_split_move(move.split.src)};
        std::array<bool, 2> found{};
        std::array<Coord, 2> arv{{move.split.arv1, move.split.arv2}};
        for (Coord const &e : list)
        {
            for (std::size_t i{0}; i < std::size(arv); i++)
            {
                if (e == arv[i])
                {
                    found[i] = true;
                }
            }
\end{minted}
\end{frame}

\begin{frame}[fragile]
\hyperlink{conclusion}{Retour au début}\newline
\begin{minted}[obeytabs=true,tabsize=2,linenos=true,breaklines,firstnumber=327]{cpp}
            if (std::all_of(
                    std::begin(found),
                    std::end(found),
                    [](bool v) -> bool
                    { return v; }))
            {
                break;
            }
        }
    }
    else if (move.type == TypeMove::MERGE)
    {
        std::forward_list<Coord> list1{
            get_list_split_move(move.merge.src1)};
        std::forward_list<Coord> list2{
            get_list_split_move(move.merge.src2)};

        return std::find(
                   std::begin(list1),
                   std::end(list1),
                   move.merge.arv) !=
                   std::end(list1) &&
               std::find(
                   std::begin(list2),
                   std::end(list2),
                   move.merge.arv) !=
                   std::end(list2);
    }
}

\end{minted}
\end{frame}

\begin{frame}[fragile]
\hyperlink{conclusion}{Retour au début}\newline
\begin{minted}[obeytabs=true,tabsize=2,linenos=true,breaklines,firstnumber=357]{cpp}
template <std::size_t N, std::size_t M>
void Board<N, M>::update_case(std::size_t pos) noexcept
{
    std::size_t size_board{std::size(m_board)};
    for (std::size_t i{0}; i < size_board; i++)
    {
        if (m_board[i].first[pos])
        {
            return;
        }
    }
    m_piece_board[pos] = Piece();
}

template <std::size_t N, std::size_t M>
CONSTEXPR bool
Board<N, M>::check_if_use_move_promote(Coord const &pos) const noexcept
{
    return (*this)(pos.n, pos.m).check_if_use_move_promote(*this, pos);
}

template <std::size_t N, std::size_t M>
CONSTEXPR std::forward_list<Move>
Board<N, M>::get_list_promote(Coord const &pos) const noexcept
{
    return (*this)(pos.n, pos.m).get_list_promote(*this, pos);
}

template <std::size_t N, std::size_t M>
template <class UnitaryFunction>
\end{minted}
\end{frame}

\begin{frame}[fragile]
\hyperlink{conclusion}{Retour au début}\newline
\begin{minted}[obeytabs=true,tabsize=2,linenos=true,breaklines,firstnumber=387]{cpp}
CONSTEXPR void
Board<N, M>::all_move(
    UnitaryFunction func,
    std::optional<Color> color_player) const noexcept
{
    std::unordered_map<TypePiece,
                       std::unordered_map<
                           Coord,
                           std::vector<Coord>,
                           Coord_hash>>
        move_merge{};

    for (std::size_t i{0}; i < numberLines(); i++)
    {
        for (std::size_t j{0}; j < numberColumns(); j++)
        {
            Piece const &piece{(*this)(i, j)};
            double p_proba{get_proba(Coord(i, j))};
            if (piece.get_type() != TypePiece::EMPTY &&
                piece.get_color() == color_player
                                         .value_or(get_current_player()))
            {
                if (!check_if_use_move_promote(Coord(i, j)))
                {
                    std::forward_list<Coord> move_normal{
                        get_list_normal_move(Coord(i, j))};

                    for (Coord const &c : move_normal)
                    {
                        if (func(Move_classic(Coord(i, j), c)))
\end{minted}
\end{frame}

\begin{frame}[fragile]
\hyperlink{conclusion}{Retour au début}\newline
\begin{minted}[obeytabs=true,tabsize=2,linenos=true,breaklines,firstnumber=417]{cpp}
                        {
                            return;
                        }
                    }
                    
                        std::forward_list<Coord> move_split{
                            get_list_split_move(Coord(i, j))};

                        std::forward_list<Coord>::
                            const_iterator it_move{
                                std::cbegin(move_split)};
                        for (; it_move != std::cend(move_split);
                             it_move++)
                        {
                            for (std::forward_list<Coord>::
                                     const_iterator it2{
                                         std::cbegin(move_split)};
                                 it2 != cend(move_split); it2++)
                            {
                                if (it_move == it2)
                                {
                                    continue;
                                }
                                if (func(Move_split(Coord(i, j), *it_move, *it2)))
                                {
                                    return;
                                }
                            }
                            if (piece.get_type() != TypePiece::PAWN)
                            {
\end{minted}
\end{frame}

\begin{frame}[fragile]
\hyperlink{conclusion}{Retour au début}\newline
\begin{minted}[obeytabs=true,tabsize=2,linenos=true,breaklines,firstnumber=447]{cpp}
                                if (move_merge
                                        .contains(piece.get_type()))
                                {
                                    if (move_merge
                                            .at(piece.get_type())
                                            .contains(*it_move))
                                    {
                                        for (Coord const &c :
                                             move_merge
                                                 .at(piece.get_type())
                                                 .at(*it_move))
                                        {
                                            if (p_proba < 1. ||
                                                get_proba(
                                                    Coord(
                                                        c.n,
                                                        c.m)) < 1.)
                                            {
                                                if (func(Move_merge(
                                                        Coord(i, j),
                                                        c,
                                                        *it_move)))
                                                {
                                                    return;
                                                }
                                            }
                                        }
                                    }
                                    else
                                    {
\end{minted}
\end{frame}

\begin{frame}[fragile]
\hyperlink{conclusion}{Retour au début}\newline
\begin{minted}[obeytabs=true,tabsize=2,linenos=true,breaklines,firstnumber=477]{cpp}
                                        move_merge
                                            .at(piece.get_type())
                                            .insert({*it_move,
                                                     std::vector{
                                                         Coord(i, j)}});
                                    }
                                }
                                else
                                {
                                    move_merge
                                        .insert(
                                            std::make_pair(
                                                piece
                                                    .get_type(),
                                                std::unordered_map<
                                                    Coord,
                                                    std::vector<Coord>,
                                                    Coord_hash>{
                                                    std::make_pair(
                                                        *it_move,
                                                        std::vector{
                                                            Coord(i, j)})}));
                                }
                            }
                        }
                    
                }
                else
                {
                    std::forward_list<Move> list{
\end{minted}
\end{frame}

\begin{frame}[fragile]
\hyperlink{conclusion}{Retour au début}\newline
\begin{minted}[obeytabs=true,tabsize=2,linenos=true,breaklines,firstnumber=507]{cpp}
                        get_list_promote(Coord(i, j))};

                    for (Move const &move : list)
                    {
                        if (func(move))
                        {
                            return;
                        }
                    }
                }
            }
        }
    }
}

template <std::size_t N, std::size_t M>
void Board<N, M>::update_board_classic() noexcept
{
    std::size_t size_board = std::size(m_board);
    for (std::size_t i{size_board}; i > 0; i--)
    {
        using namespace std::complex_literals;
        if (complex_equal(m_board[i - 1].second, 0i))
        {
            m_board.erase(std::begin(m_board) + i - 1);
        }
        else
        {
            for (std::size_t j{i - 1}; j > 0; j--)
            {
\end{minted}
\end{frame}

\begin{frame}[fragile]
\hyperlink{conclusion}{Retour au début}\newline
\begin{minted}[obeytabs=true,tabsize=2,linenos=true,breaklines,firstnumber=537]{cpp}
                if (m_board[i - 1].first == m_board[j - 1].first)
                {
                    double inter = std::pow(std::abs(m_board[j - 1].second), 2)+ std::pow(std::abs(m_board[i - 1].second), 2);
                    m_board[j - 1].second =std::pow(inter, 1./2 );
                    m_board.erase(std::begin(m_board) + i - 1);
                    break;
                }
            }
        }
    }
    /*std::size_t n_size_board = std::size(m_board);
    double all_proba {0};
    for (std::size_t i{0}; i<n_size_board; i++)
    {
        all_proba += std::pow(std::abs(m_board[i].second), 2);
    }
    all_proba = std::pow(all_proba, 1./2);
    for (std::size_t i{0}; i<n_size_board; i++)
    {
        m_board[i].second /= all_proba;
    }*/

\end{minted}
\end{frame}

\begin{frame}[fragile]
\label{mesure.tpp}
\frametitle{mesure.tpp}
\hyperlink{conclusion}{Retour au début}\newline
\begin{minted}[obeytabs=true,tabsize=2,linenos=true, breaklines]{cpp}
// Lib standard
#include <cassert>
#include <utility>
#include <functional>
#include <cmath>
#include <optional>
#include <stdexcept>
#include <iostream>

// Inclusion projet
#include <Piece.hpp>
#include <TypePiece.hpp>
#include <Constexpr.hpp>
#include "Board.hpp"

template <std::size_t N, std::size_t M>
CONSTEXPR bool Board<N, M>::mesure(Coord const &p,
                                   std::optional<bool> val_mes)
{

    Piece p_actuelle = (*this)(p.n, p.m);
    if (p_actuelle.get_type() == TypePiece::EMPTY)
    {
        return false;
    }
    else
\end{minted}
\end{frame}

\begin{frame}[fragile]
\hyperlink{conclusion}{Retour au début}\newline
\begin{minted}[obeytabs=true,tabsize=2,linenos=true,breaklines,firstnumber=27]{cpp}
    {
        bool mes;
        if (val_mes == std::nullopt)
        {
            double x = rnd::randreal(0., 1.);
            std::size_t indice_mes = 0;
            double pow_coef{
                std::pow(
                    std::abs(m_board[indice_mes].second), 2)};
            std::size_t size_board{std::size(m_board)};
            while (x - pow_coef > 0 && indice_mes < size_board - 1)
            {
                x -= pow_coef;
                indice_mes++;
                if (indice_mes >= size_board)
                {
                    throw std::runtime_error("Indice mesuré de mesure trop grand");
                }
                pow_coef = std::pow(std::abs(m_board[indice_mes].second), 2);
            }
            mes = m_board[indice_mes].first[offset(p.n, p.m)];
        }
        else
        {
            mes = val_mes.value();
        }
        double proba_delete = 0;
        // std::size_t nbr_elt_suppr{0};
        for (std::size_t i{std::size(m_board)}; i > 0; i--)
        {
\end{minted}
\end{frame}

\begin{frame}[fragile]
\hyperlink{conclusion}{Retour au début}\newline
\begin{minted}[obeytabs=true,tabsize=2,linenos=true,breaklines,firstnumber=57]{cpp}
            if (m_board[i - 1].first[offset(p.n, p.m)] != mes) //?
            {
                proba_delete +=
                    std::pow(std::abs(m_board[i - 1].second), 2);

                m_board.erase(
                    std::begin(m_board) + i - 1);
            }
        }
        for (auto &e : m_board)
        {
            e.second /= std::sqrt(1. - proba_delete);
        }
        for (std::size_t i{0}; i < N * M; i++)
        {
            if (p_actuelle.get_type() != TypePiece::EMPTY)
            {
                update_case(i);
            }
        }
        if (std::empty(m_board))
    {
        bool erreur = p_actuelle.get_type() == TypePiece::ROOK;
        std::cout<<erreur<<std::endl;
        std::cout<<mes<<std::endl;
        throw std::runtime_error("La mesure fait nimp");
    }
        return mes;
    }
}
\end{minted}
\end{frame}

\begin{frame}[fragile]
\hyperlink{conclusion}{Retour au début}\newline
\begin{minted}[obeytabs=true,tabsize=2,linenos=true,breaklines,firstnumber=87]{cpp}

template <std::size_t N, std::size_t M>
CONSTEXPR bool
Board<N, M>::mesure_capture_slide(
    Coord const &s,
    Coord const &t,
    std::function<bool(Board<N, M> const &,
                       Coord const &,
                       Coord const &,
                       std::size_t,
                       std::optional<Coord>)>
        check_path,
    std::optional<bool> val_mes)
{
    std::size_t position = offset(s.n, s.m);
    Piece p_actuelle = (*this)(s.n, s.m);
    if (p_actuelle.get_type() == TypePiece::EMPTY)
    {
        return false;
    }
    else
    {
        bool mes;
        if (val_mes == std::nullopt)
        {
            double x = rnd::randreal(0., 1.);
            std::size_t indice_mes = 0;
            double pow_coef{
                std::pow(std::abs(m_board[0].second), 2)};
            std::size_t size_board{std::size(m_board)};
\end{minted}
\end{frame}

\begin{frame}[fragile]
\hyperlink{conclusion}{Retour au début}\newline
\begin{minted}[obeytabs=true,tabsize=2,linenos=true,breaklines,firstnumber=117]{cpp}
            while (x - pow_coef > 0 && indice_mes < size_board - 1)
            {
                x -= pow_coef;
                indice_mes++;
                if (indice_mes >= size_board)
                {
                    throw std::runtime_error("Indice mesuré de mesure trop grand");
                }
                pow_coef = std::pow(std::abs(m_board[indice_mes].second), 2);

                // indice_suppr++,
            }
            mes = m_board[indice_mes].first[position] &&
                  check_path(*this, s, t, indice_mes, std::nullopt);
        }
        else
        {
            mes = val_mes.value();
        }
        double proba_delete = 0;
        // std::size_t nbr_elt_suppr{0};
        for (std::size_t i{std::size(m_board)}; i > 0; i--)
        {
            if ((m_board[i - 1].first[position] &&
                 check_path(*this, s, t, i - 1, std::nullopt)) != mes)
            {
                proba_delete +=
                    std::pow(std::abs(m_board[i - 1].second), 2);

                m_board.erase(
\end{minted}
\end{frame}

\begin{frame}[fragile]
\hyperlink{conclusion}{Retour au début}\newline
\begin{minted}[obeytabs=true,tabsize=2,linenos=true,breaklines,firstnumber=147]{cpp}
                    std::begin(m_board) + i - 1);
            }
        }
        for (auto &e : m_board)
        {
            e.second /= std::sqrt(1. - proba_delete);
        }
        for (std::size_t i{0}; i < N * M; i++)
        {
            if (p_actuelle.get_type() != TypePiece::EMPTY)
            {
                update_case(i);
            }
        }
        if (std::empty(m_board))
    {
        throw std::runtime_error("La mesure_capture_slide fait nimp");
    }
        return mes;
    }
}

template <std::size_t N, std::size_t M>
CONSTEXPR bool
Board<N, M>::mesure_castle(
    Coord const &king,
    Coord const &rook,
    std::optional<bool> val_mes)
{
    std::size_t position_king = offset(king.n, king.m);
\end{minted}
\end{frame}

\begin{frame}[fragile]
\hyperlink{conclusion}{Retour au début}\newline
\begin{minted}[obeytabs=true,tabsize=2,linenos=true,breaklines,firstnumber=177]{cpp}
    std::size_t position_rook = offset(rook.n, rook.m);

    bool mes;
    if (val_mes == std::nullopt)
    {
        double x = rnd::randreal(0., 1.);
        std::size_t indice_mes = 0;
        std::size_t size_board{std::size(m_board)};
        double pow_coef{
            std::pow(std::abs(m_board[0].second), 2)};

        while (x - pow_coef > 0 && indice_mes < size_board - 1)
        {
            x -= pow_coef;
            indice_mes++;
            if (indice_mes >= size_board)
            {
                throw std::runtime_error("Indice mesuré de mesure trop grand");
            }
            pow_coef = std::pow(std::abs(m_board[indice_mes].second), 2);
        }
        mes = m_board[indice_mes].first[position_king] &&
              m_board[indice_mes].first[position_rook] &&
              check_path_straight_1_instance(*this, king, rook, indice_mes, std::nullopt);
    }
    else
    {
        mes = val_mes.value();
    }
    double proba_delete = 0;
\end{minted}
\end{frame}

\begin{frame}[fragile]
\hyperlink{conclusion}{Retour au début}\newline
\begin{minted}[obeytabs=true,tabsize=2,linenos=true,breaklines,firstnumber=207]{cpp}
    // std::size_t nbr_elt_suppr{0};
    for (std::size_t i{std::size(m_board)}; i > 0; i--)
    {
        if ((m_board[i - 1].first[position_king] &&
             m_board[i - 1].first[position_rook] &&
             check_path_straight_1_instance(*this, king, rook, i - 1, std::nullopt)) != mes)
        {
            proba_delete +=
                std::pow(std::abs(m_board[i - 1].second), 2);

            m_board.erase(
                std::begin(m_board) + i - 1);
        }
    }
    for (auto &e : m_board)
    {
        e.second /= std::sqrt(1. - proba_delete);
    }
    for (std::size_t i{0}; i < N * M; i++)
    {
        if (m_piece_board[i].get_type() != TypePiece::EMPTY)
        {
            update_case(i);
        }
    }
    return mes;

\end{minted}
\end{frame}

\begin{frame}[fragile]
\label{move.tpp}
\frametitle{move.tpp}
\hyperlink{conclusion}{Retour au début}\newline
\begin{minted}[obeytabs=true,tabsize=2,linenos=true, breaklines]{cpp}
// Lib standard
#include <array>
#include <algorithm>
#include <cassert>
#include <utility>
#include <functional>
#include <cmath>
#include <optional>
#include <stdexcept>

// Inclusion projet
#include <Piece.hpp>
#include <CMatrix.hpp>
#include <Unitary.hpp>
#include <TypePiece.hpp>
#include <math_utility.hpp>
#include <Constexpr.hpp>
#include <Move.hpp>
#include "Board.hpp"

template <std::size_t N, std::size_t M>
CONSTEXPR void Board<N, M>::king_side_castle(Coord const &k,
                                             Coord const &r,
                                             std::optional<bool> val_mes)
{
    std::size_t king = offset(k.n, k.m);
\end{minted}
\end{frame}

\begin{frame}[fragile]
\hyperlink{conclusion}{Retour au début}\newline
\begin{minted}[obeytabs=true,tabsize=2,linenos=true,breaklines,firstnumber=27]{cpp}
    std::size_t new_king = offset(k.n, k.m + 2);
    std::size_t rook = offset(r.n, r.m);
    std::size_t new_rook = offset(r.n, r.m - 2);
    if (mesure_castle(k, r, val_mes))
    {
        for (std::size_t i{0}; i < std::size(m_board); i++)
        {
            move_1_instance(
                std::array<bool, 2>{true, false}, i,
                MATRIX_ISWAP,
                std::array<std::size_t, 2>{king, new_king});
            move_1_instance(
                std::array<bool, 2>{true, false}, i,
                MATRIX_ISWAP,
                std::array<std::size_t, 2>{rook, new_rook});
        }
        m_piece_board[new_king] = std::move(m_piece_board[king]);
        m_piece_board[new_rook] = std::move(m_piece_board[rook]);
        m_piece_board[king] = Piece();
        m_piece_board[rook] = Piece();
    }
}
template <std::size_t N, std::size_t M>
CONSTEXPR void Board<N, M>::queen_side_castle(Coord const &k,
                                              Coord const &r,
                                              std::optional<bool> val_mes)
{
    std::size_t king = offset(k.n, k.m);
    std::size_t new_king = offset(k.n, k.m - 2);
    std::size_t rook = offset(r.n, r.m);
\end{minted}
\end{frame}

\begin{frame}[fragile]
\hyperlink{conclusion}{Retour au début}\newline
\begin{minted}[obeytabs=true,tabsize=2,linenos=true,breaklines,firstnumber=57]{cpp}
    std::size_t new_rook = offset(r.n, r.m + 3);
    if (mesure_castle(k, r, val_mes))
    {
        for (std::size_t i{0}; i < std::size(m_board); i++)
        {
            move_1_instance(
                std::array<bool, 2>{true, false}, i,
                MATRIX_ISWAP,
                std::array<std::size_t, 2>{king, new_king});
            move_1_instance(
                std::array<bool, 2>{true, false}, i,
                MATRIX_ISWAP,
                std::array<std::size_t, 2>{rook, new_rook});
        }
        m_piece_board[new_king] = std::move(m_piece_board[king]);
        m_piece_board[new_rook] = std::move(m_piece_board[rook]);
        m_piece_board[king] = Piece();
        m_piece_board[rook] = Piece();
    }
}
template <std::size_t N, std::size_t M>
CONSTEXPR void
Board<N, M>::move_classic_jump(Coord const &s, Coord const &t,
                               std::optional<bool> val_mes)
{
    std::size_t source = offset(s.n, s.m);
    std::size_t target = offset(t.n, t.m);

    if (m_piece_board[target].get_type() == TypePiece::EMPTY ||
        (m_piece_board[target].get_type() == m_piece_board[source].get_type() &&
\end{minted}
\end{frame}

\begin{frame}[fragile]
\hyperlink{conclusion}{Retour au début}\newline
\begin{minted}[obeytabs=true,tabsize=2,linenos=true,breaklines,firstnumber=87]{cpp}
         m_piece_board[target].get_color() == m_piece_board[source].get_color()))
    {
        std::size_t const size_board{std::size(m_board)};
        for (std::size_t i{0}; i < size_board; i++)
        {
            move_1_instance(
                std::array<bool, 2>{m_board[i].first[source], false}, i,
                MATRIX_ISWAP,
                std::array<std::size_t, 2>{source, target});
        }
        m_piece_board[target] = std::move(m_piece_board[source]);
        m_piece_board[source] = Piece();
    }
    else
    {
        if (m_piece_board[source].same_color(m_piece_board[target]))
        {
            std::optional<bool> negation ;
            if(val_mes.has_value())
            {
                negation = !val_mes.value();
            }
            if (!mesure(t, negation))
            {
                for (std::size_t i{0}; i < std::size(m_board); i++)
                {
                    move_1_instance(
                        std::array<bool, 2>{m_board[i].first[source], false},
                        i, MATRIX_ISWAP,
                        std::array<std::size_t, 2>{source, target});
\end{minted}
\end{frame}

\begin{frame}[fragile]
\hyperlink{conclusion}{Retour au début}\newline
\begin{minted}[obeytabs=true,tabsize=2,linenos=true,breaklines,firstnumber=117]{cpp}
                }
                m_piece_board[target] = std::move(m_piece_board[source]);
                m_piece_board[source] = Piece();
            }
        }
        else
        {
            if (mesure(s, val_mes))
            {
                for (std::size_t i{0}; i < std::size(m_board); i++)
                {
                    CONSTEXPR CMatrix<8> A{
                        MATRIX_JUMP
                            .tensoriel_product(
                                CMatrix<2>::identity()) *
                        CMatrix<2>::identity()
                            .tensoriel_product(MATRIX_JUMP)};

                    move_1_instance(
                        std::array<bool, 3>{m_board[i].first[source],
                                            m_board[i].first[target],
                                            false},
                        i,
                        A,
                        std::array<std::size_t, 3>{source,
                                                   target,
                                                   N * M + 1});
                }
                m_piece_board[target] = std::move(m_piece_board[source]);
                m_piece_board[source] = Piece();
\end{minted}
\end{frame}

\begin{frame}[fragile]
\hyperlink{conclusion}{Retour au début}\newline
\begin{minted}[obeytabs=true,tabsize=2,linenos=true,breaklines,firstnumber=147]{cpp}
            }
        }
    }
}

template <std::size_t N, std::size_t M>
CONSTEXPR bool
Board<N, M>::move_pawn_one_step(Coord const &s, Coord const &t,
                                std::optional<bool> val_mes)
{
    std::size_t source = offset(s.n, s.m);
    std::size_t target = offset(t.n, t.m);

    if (m_piece_board[target].get_type() == TypePiece::EMPTY ||
        (m_piece_board[target].get_type() == m_piece_board[source].get_type() &&
         m_piece_board[target].get_color() == m_piece_board[source].get_color()))
    {
        std::size_t const size_board{std::size(m_board)};
        for (std::size_t i{0}; i < size_board; i++)
        {
            move_1_instance(
                std::array<bool, 2>{m_board[i].first[source], false},
                i, MATRIX_ISWAP,
                std::array<std::size_t, 2>{source, target});
        }
        m_piece_board[target] = std::move(m_piece_board[source]);
        m_piece_board[source] = Piece();
        return true;
    }
    else
\end{minted}
\end{frame}

\begin{frame}[fragile]
\hyperlink{conclusion}{Retour au début}\newline
\begin{minted}[obeytabs=true,tabsize=2,linenos=true,breaklines,firstnumber=177]{cpp}
    {
        std::optional<bool> negation ;
            if(val_mes.has_value())
            {
                negation = !val_mes.value();
            }
        if (!mesure(t, negation))
        {
            for (std::size_t i{0}; i < std::size(m_board); i++)
            {
                move_1_instance(
                    std::array<bool, 2>{m_board[i].first[source], false},
                    i, MATRIX_ISWAP,
                    std::array<std::size_t, 2>{source, target});
            }
            m_piece_board[target] = std::move(m_piece_board[source]);
            m_piece_board[source] = Piece();
            return true;
        }
    }
    return false;
}

template <std::size_t N, std::size_t M>
CONSTEXPR bool
Board<N, M>::move_pawn_two_step(Coord const &s, Coord const &t,
                                std::optional<bool> val_mes)

{
    std::size_t source = offset(s.n, s.m);
\end{minted}
\end{frame}

\begin{frame}[fragile]
\hyperlink{conclusion}{Retour au début}\newline
\begin{minted}[obeytabs=true,tabsize=2,linenos=true,breaklines,firstnumber=207]{cpp}
    std::size_t target = offset(t.n, t.m);
    if (m_piece_board[target].get_type() == TypePiece::EMPTY ||
        (m_piece_board[target].get_type() == m_piece_board[source].get_type() &&
         m_piece_board[target].get_color() == m_piece_board[source].get_color()))
    {
        std::size_t const size_board{std::size(m_board)};
        for (std::size_t i{0}; i < size_board; i++)
        {
            move_1_instance(
                std::array<bool, 3>{!check_path_straight_1_instance(
                                        *this, s, t, i, std::nullopt),
                                    false,
                                    m_board[i].first[source]},
                i, MATRIX_SLIDE,
                std::array<std::size_t, 3>{N * M + 1, target, source});
        }
        m_piece_board[target] = m_piece_board[source];
        update_case(source);
        return true;
    }
    else
    {
        std::optional<bool> negation ;
            if(val_mes.has_value())
            {
                negation = !val_mes.value();
            }
        if (!mesure(t, negation))
        {
            for (std::size_t i{0}; i < std::size(m_board); i++)
\end{minted}
\end{frame}

\begin{frame}[fragile]
\hyperlink{conclusion}{Retour au début}\newline
\begin{minted}[obeytabs=true,tabsize=2,linenos=true,breaklines,firstnumber=237]{cpp}
            {
                move_1_instance(
                    std::array<bool, 3>{!check_path_straight_1_instance(
                                            *this, s, t, i, std::nullopt),
                                        false,
                                        m_board[i].first[source]},
                    i, MATRIX_SLIDE,
                    std::array<std::size_t, 3>{N * M + 1, target, source});
            }
            m_piece_board[target] = std::move(m_piece_board[source]);
            update_case(source);
            return true;
        }
    }
    return false;
}

template <std::size_t N, std::size_t M>
CONSTEXPR bool Board<N, M>::capture_pawn(
    Coord const &s,
    Coord const &t,
    std::optional<bool> val_mes)
{
    std::size_t source = offset(s.n, s.m);
    std::size_t target = offset(t.n, t.m);
    if (mesure(s, val_mes))
    {
        for (std::size_t i{0}; i < std::size(m_board); i++)
        {
            if (m_board[i].first[target])
\end{minted}
\end{frame}

\begin{frame}[fragile]
\hyperlink{conclusion}{Retour au début}\newline
\begin{minted}[obeytabs=true,tabsize=2,linenos=true,breaklines,firstnumber=267]{cpp}
            {
                CONSTEXPR CMatrix<8> A{
                    MATRIX_JUMP
                        .tensoriel_product(
                            CMatrix<2>::identity()) *
                    CMatrix<2>::identity()
                        .tensoriel_product(MATRIX_JUMP)};

                move_1_instance(
                    std::array<bool, 3>{
                        m_board[i].first[source],
                        m_board[i].first[target],
                        false},
                    i, A,
                    std::array<std::size_t, 3>{source, target, N * M + 1});
            }
        }
        m_piece_board[target] = m_piece_board[source];
        update_case(source);
        return true;
    }
    return false;
}

template <std::size_t N, std::size_t M>
CONSTEXPR bool
Board<N, M>::move_enpassant(Coord const &s, Coord const &t, Coord const &ep,
                            std::optional<bool> val_mes)
{
    std::size_t source = offset(s.n, s.m);
\end{minted}
\end{frame}

\begin{frame}[fragile]
\hyperlink{conclusion}{Retour au début}\newline
\begin{minted}[obeytabs=true,tabsize=2,linenos=true,breaklines,firstnumber=297]{cpp}
    std::size_t target = offset(t.n, t.m);
    std::size_t enpassant = offset(ep.n, ep.m);

    if (m_piece_board[target].get_type() == TypePiece::EMPTY ||
        (m_piece_board[target].get_type() == m_piece_board[source].get_type() &&
         m_piece_board[target].get_color() == m_piece_board[source].get_color()))
    {

        std::size_t const size_board{std::size(m_board)};
        for (std::size_t i{0}; i < size_board; i++)
        {
            if (m_board[i].first[enpassant])
            {
                // permet d'enlever le pion pris en passant
                move_1_instance(
                    std::array<bool, 2>{m_board[i].first[enpassant], false},
                    i, MATRIX_ISWAP,
                    std::array<std::size_t, 2>{enpassant, N * M + 1});
                move_1_instance(
                    std::array<bool, 2>{m_board[i].first[source], false},
                    i, MATRIX_ISWAP,
                    std::array<std::size_t, 2>{source, target});
            }
        }
        m_piece_board[target] = std::move(m_piece_board[source]);
        m_piece_board[source] = Piece();
        m_piece_board[enpassant] = Piece();
        return true;
    }
    else
\end{minted}
\end{frame}

\begin{frame}[fragile]
\hyperlink{conclusion}{Retour au début}\newline
\begin{minted}[obeytabs=true,tabsize=2,linenos=true,breaklines,firstnumber=327]{cpp}
    {
        if (m_piece_board[source].same_color(m_piece_board[target]))
        {
            std::optional<bool> negation ;
            if(val_mes.has_value())
            {
                negation = !val_mes.value();
            }
            if (!mesure(t, negation))
            {

                for (std::size_t i{0}; i < std::size(m_board); i++)
                {
                    if (m_board[i].first[enpassant])
                    {
                        // permet d'enlever le pion pris en passant
                        move_1_instance(
                            std::array<bool, 2>{
                                m_board[i].first[enpassant],
                                false},
                            i, MATRIX_ISWAP,
                            std::array<std::size_t, 2>{enpassant, N * M + 1});

                        move_1_instance(
                            std::array<bool, 2>{
                                m_board[i].first[source],
                                false},
                            i, MATRIX_ISWAP,
                            std::array<std::size_t, 2>{source, target});
                    }
\end{minted}
\end{frame}

\begin{frame}[fragile]
\hyperlink{conclusion}{Retour au début}\newline
\begin{minted}[obeytabs=true,tabsize=2,linenos=true,breaklines,firstnumber=357]{cpp}
                }
                m_piece_board[target] = m_piece_board[source];
                m_piece_board[source] = Piece();
                return true;
            }
            return false;
        }
        else
        {
            if (mesure(s, val_mes))
            {
                for (std::size_t i{0}; i < std::size(m_board); i++)
                {
                    if (m_board[i].first[enpassant] ||
                        m_board[i].first[target])
                    {
                        // permet d'enlever le pion pris en passant
                        move_1_instance(
                            std::array<bool, 2>{
                                m_board[i].first[enpassant], false},
                            i, MATRIX_ISWAP,
                            std::array<std::size_t, 2>{
                                enpassant, N * M + 1});
                        // ces deux instructions représentent un mouvement
                        // de capture jump
                        move_1_instance(
                            std::array<bool, 2>{
                                m_board[i].first[target], false},
                            i, MATRIX_ISWAP,
                            std::array<std::size_t, 2>{
\end{minted}
\end{frame}

\begin{frame}[fragile]
\hyperlink{conclusion}{Retour au début}\newline
\begin{minted}[obeytabs=true,tabsize=2,linenos=true,breaklines,firstnumber=387]{cpp}
                                target, N * M + 1});

                        move_1_instance(
                            std::array<bool, 2>{
                                m_board[i].first[source], false},
                            i, MATRIX_ISWAP,
                            std::array<std::size_t, 2>{
                                source, target});
                    }
                }
                m_piece_board[target] = std::move(m_piece_board[source]);
                m_piece_board[source] = Piece();
                m_piece_board[enpassant] = Piece();
                return true;
            }
        }
    }
    return false;
}

template <std::size_t N, std::size_t M>
CONSTEXPR void
Board<N, M>::move_split_jump(Coord const &s, Coord const &t1, Coord const &t2)
{
    std::size_t const source = offset(s.n, s.m);
    std::size_t const target1 = offset(t1.n, t1.m);
    std::size_t const target2 = offset(t2.n, t2.m);
    std::size_t const size_board{std::size(m_board)};
    for (std::size_t i{0}; i < size_board; i++)
    {
\end{minted}
\end{frame}

\begin{frame}[fragile]
\hyperlink{conclusion}{Retour au début}\newline
\begin{minted}[obeytabs=true,tabsize=2,linenos=true,breaklines,firstnumber=417]{cpp}
        move_1_instance(
            std::array<bool, 3>{
                m_board[i].first[target2],
                m_board[i].first[target1],
                m_board[i].first[source]},
            i, MATRIX_SPLIT,
            std::array<std::size_t, 3>{target2, target1, source});
    }
    m_piece_board[target1] = m_piece_board[source];

    if (m_piece_board[target1].get_type() == TypePiece::EMPTY &&
        m_piece_board[target2].get_type() == TypePiece::EMPTY)
    {
        m_piece_board[target2] = std::move(m_piece_board[source]);
        m_piece_board[source] = Piece();
        update_case(target1);
        update_case(target2);
    }
    else
    {
        m_piece_board[target2] = m_piece_board[source];
        update_case(source);
        update_case(target1);
        update_case(target2);
        update_board();
    }
}

template <std::size_t N, std::size_t M>
CONSTEXPR void
\end{minted}
\end{frame}

\begin{frame}[fragile]
\hyperlink{conclusion}{Retour au début}\newline
\begin{minted}[obeytabs=true,tabsize=2,linenos=true,breaklines,firstnumber=447]{cpp}
Board<N, M>::move_merge_jump(Coord const &s1, Coord const &s2, Coord const &t)
{
    std::size_t source1 = offset(s1.n, s1.m);
    std::size_t source2 = offset(s2.n, s2.m);
    std::size_t target = offset(t.n, t.m);
    std::size_t const size_board{std::size(m_board)};
    for (std::size_t i{0}; i < size_board; i++)
    {
        move_1_instance(
            std::array<bool, 3>{
                m_board[i].first[source1],
                m_board[i].first[source2],
                m_board[i].first[target]},
            i, MATRIX_MERGE,
            std::array<std::size_t, 3>{source1, source2, target});
    }
    m_piece_board[target] = m_piece_board[source1];
    update_board();
    update_case(target);
    update_case(source1);
    update_case(source2);
}

template <std::size_t N, std::size_t M>
CONSTEXPR void
Board<N, M>::move_classic_slide(
    Coord const &s,
    Coord const &t,
    std::function<
        bool(
\end{minted}
\end{frame}

\begin{frame}[fragile]
\hyperlink{conclusion}{Retour au début}\newline
\begin{minted}[obeytabs=true,tabsize=2,linenos=true,breaklines,firstnumber=477]{cpp}
            Board<N, M> const &,
            Coord const &,
            Coord const &,
            std::size_t,
            std::optional<Coord>)>
        check_path,
    std::optional<bool> val_mes)
{
    std::size_t source = offset(s.n, s.m);
    std::size_t target = offset(t.n, t.m);
    if (m_piece_board[target].get_type() == TypePiece::EMPTY ||
        (m_piece_board[target].get_type() == m_piece_board[source].get_type() &&
         m_piece_board[target].get_color() == m_piece_board[source].get_color()))
    {
        std::size_t const size_board{std::size(m_board)};
        for (std::size_t i{0}; i < size_board; i++)
        {
            move_1_instance(
                std::array<bool, 3>{
                    !check_path(
                        *this, s, t, i, std::nullopt),
                    false,
                    m_board[i].first[source]},
                i, MATRIX_SLIDE,
                std::array<std::size_t, 3>{N * M + 1, target, source});
        }
        m_piece_board[target] = m_piece_board[source];
        update_case(source);
    }
    else
\end{minted}
\end{frame}

\begin{frame}[fragile]
\hyperlink{conclusion}{Retour au début}\newline
\begin{minted}[obeytabs=true,tabsize=2,linenos=true,breaklines,firstnumber=507]{cpp}
    {
        if (m_piece_board[source].same_color(m_piece_board[target]))
        {
            std::optional<bool> negation ;
            if(val_mes.has_value())
            {
                negation = !val_mes.value();
            }
            if (!mesure(t, negation))
            {
                for (std::size_t i{0}; i < std::size(m_board); i++)
                {
                    move_1_instance(
                        std::array<bool, 3>{
                            !check_path(*this, s, t, i, std::nullopt),
                            false,
                            m_board[i].first[source]},
                        i, MATRIX_SLIDE,
                        std::array<std::size_t, 3>{N * M + 1, target, source});
                }
                m_piece_board[target] = std::move(m_piece_board[source]);
                m_piece_board[source] = Piece();
            }
        }
        else
        {
            if (mesure_capture_slide(s, t, check_path, val_mes))
            {
                for (std::size_t i{0}; i < std::size(m_board); i++)
                {
\end{minted}
\end{frame}

\begin{frame}[fragile]
\hyperlink{conclusion}{Retour au début}\newline
\begin{minted}[obeytabs=true,tabsize=2,linenos=true,breaklines,firstnumber=537]{cpp}
                    auto const A{
                        MATRIX_ISWAP.tensoriel_product(
                            CMatrix<2>::identity()) *
                        CMatrix<2>::identity()
                            .tensoriel_product(MATRIX_ISWAP)};

                    move_1_instance(
                        std::array<bool, 3>{
                            m_board[i].first[source],
                            m_board[i].first[target], false},
                        i, A,
                        std::array<std::size_t, 3>{source, target, N * M + 1});
                }
                m_piece_board[target] = std::move(m_piece_board[source]);
                m_piece_board[source] = Piece();
            }
        }
    }
}

template <std::size_t N, std::size_t M>
CONSTEXPR void
Board<N, M>::move_split_slide(
    Coord const &s,
    Coord const &t1,
    Coord const &t2,
    std::function<
        bool(
            Board<N, M> const &,
            Coord const &,
\end{minted}
\end{frame}

\begin{frame}[fragile]
\hyperlink{conclusion}{Retour au début}\newline
\begin{minted}[obeytabs=true,tabsize=2,linenos=true,breaklines,firstnumber=567]{cpp}
            Coord const &,
            std::size_t,
            std::optional<Coord>)>
        check_path)
{
    std::size_t source = offset(s.n, s.m);
    std::size_t target1 = offset(t1.n, t1.m);
    std::size_t target2 = offset(t2.n, t2.m);
    std::size_t const size_board{std::size(m_board)};
    for (std::size_t i{0}; i < size_board; i++)
    {
        move_1_instance(
            std::array<bool, 5>{
                !check_path(*this, s, t2, i, t1),
                !check_path(*this, s, t1, i, t2),
                m_board[i].first[target2],
                m_board[i].first[target1],
                m_board[i].first[source]},
            i, MATRIX_SPLIT_SLIDE,
            std::array<std::size_t, 5>{
                N * M + 1,
                N * M + 1,
                target2,
                target1,
                source});
        /*La matrice split slide est créée pour
          être utlisé sans appliquer de porte cnot a
          un chemin, nos fonctions check_path renvoie
          un résultat où l'on a appliquer la porte cnot */
    }
\end{minted}
\end{frame}

\begin{frame}[fragile]
\hyperlink{conclusion}{Retour au début}\newline
\begin{minted}[obeytabs=true,tabsize=2,linenos=true,breaklines,firstnumber=597]{cpp}
    if (m_piece_board[target1].get_type() == TypePiece::EMPTY &&
        m_piece_board[target2].get_type() == TypePiece::EMPTY)
    {
        m_piece_board[target1] = m_piece_board[source];
        m_piece_board[target2] = m_piece_board[source];
        update_case(source);
        update_case(target1);
        update_case(target2);
    }
    else
    {
        m_piece_board[target1] = m_piece_board[source];
        m_piece_board[target2] = m_piece_board[source];
        update_case(source);
        update_case(target1);
        update_case(target2);
        update_board();
    }
}

template <std::size_t N, std::size_t M>
CONSTEXPR void
Board<N, M>::move_merge_slide(
    Coord const &s1,
    Coord const &s2,
    Coord const &t,
    std::function<
        bool(
            Board<N, M> const &,
            Coord const &,
\end{minted}
\end{frame}

\begin{frame}[fragile]
\hyperlink{conclusion}{Retour au début}\newline
\begin{minted}[obeytabs=true,tabsize=2,linenos=true,breaklines,firstnumber=627]{cpp}
            Coord const &,
            std::size_t,
            std::optional<Coord>)>
        check_path)
{
    std::size_t const source1 = offset(s1.n, s1.m);
    std::size_t const source2 = offset(s2.n, s2.m);
    std::size_t const target = offset(t.n, t.m);
    std::size_t const size_board{std::size(m_board)};
    for (std::size_t i{0}; i < size_board; i++)
    {
        move_1_instance(
            std::array<bool, 5>{
                !check_path(*this, s2, t, i, s1),
                !check_path(*this, s1, t, i, s2),
                m_board[i].first[source1],
                m_board[i].first[source2],
                m_board[i].first[target]},
            i, MATRIX_MERGE_SLIDE,
            std::array<std::size_t, 5>{
                N * M + 1,
                N * M + 1,
                source1,
                source2,
                target});
        /* La matrice merge slide est créée pour
           être utlisé sans appliquer de porte cnot au chemin,
           nos fonctions check_path renvoie un résultat où l'on
           a appliquer la porte cnot */
    }
\end{minted}
\end{frame}

\begin{frame}[fragile]
\hyperlink{conclusion}{Retour au début}\newline
\begin{minted}[obeytabs=true,tabsize=2,linenos=true,breaklines,firstnumber=657]{cpp}
    m_piece_board[target] = m_piece_board[source1];
    update_board();
    update_case(target);
    update_case(source1);
    update_case(source2);
}

template <std::size_t N, std::size_t M>
CONSTEXPR bool Board<N, M>::move_pawn(
    Coord const &s,
    Coord const &t,
    std::optional<bool> val_mes) noexcept
{
    bool move_exec{false};
    auto abs_diff{[](std::size_t x, std::size_t y) -> std::size_t
                  {
                      return (x >= y) ? x - y : y - x;
                  }};
    std::size_t diff_line{abs_diff(s.n, t.n)};
    if (diff_line == 2)
    {
        move_exec = move_pawn_two_step(s, t, val_mes);
        m_ep = Coord((s.n + t.n) / 2, s.m);
    }
    else if (diff_line == 1)
    {
        std::size_t diff_col{abs_diff(s.m, t.m)};
        if (diff_col == 1)
        {

\end{minted}
\end{frame}

\begin{frame}[fragile]
\hyperlink{conclusion}{Retour au début}\newline
\begin{minted}[obeytabs=true,tabsize=2,linenos=true,breaklines,firstnumber=687]{cpp}
            if (m_ep != std::nullopt && m_ep == t)
            {
                int step{
                    ((*this)(s.n, s.m).get_color() ==
                     Color::WHITE)
                        ? -1
                        : 1};
                Coord ep;
                ep.m = m_ep->m;
                ep.n = m_ep->n - step;
                move_exec = move_enpassant(s, t, ep, val_mes);
            }
            else
            {
                move_exec = capture_pawn(s, t, val_mes);
            }
        }
        else
        {
            move_exec = move_pawn_one_step(s, t, val_mes);
        }
        m_ep = std::nullopt;
    }
    return move_exec;
}

template <std::size_t N, std::size_t M>
CONSTEXPR void
Board<N, M>::move_promotion(
    Move const &move,
\end{minted}
\end{frame}

\begin{frame}[fragile]
\hyperlink{conclusion}{Retour au début}\newline
\begin{minted}[obeytabs=true,tabsize=2,linenos=true,breaklines,firstnumber=717]{cpp}
    std::optional<bool> val_mes)
{
    TypePiece p{move.promote.piece};
    Coord s{move.promote.src};
    Coord t{move.promote.arv};

    if (p == TypePiece::QUEEN ||
        p == TypePiece::KNIGHT ||
        p == TypePiece::BISHOP ||
        p == TypePiece::ROOK)
    {
        if (move_pawn(s, t, val_mes))
        {
            Piece piece{p, m_piece_board[offset(t.n, t.m)].get_color()};
            m_piece_board[offset(t.n, t.m)] = std::move(piece);
        }
    }
    else
    {
        throw std::runtime_error("Piece non valide pour faire une promotion");
    }
}

template <std::size_t N, std::size_t M>
CONSTEXPR void
Board<N, M>::move_classic(
    Coord const &s,
    Coord const &t,
    std::optional<bool> val_mes)
{
\end{minted}
\end{frame}

\begin{frame}[fragile]
\hyperlink{conclusion}{Retour au début}\newline
\begin{minted}[obeytabs=true,tabsize=2,linenos=true,breaklines,firstnumber=747]{cpp}
    TypePiece piece{(*this)(s.n, s.m).get_type()};
    switch (piece)
    {
    case TypePiece::KING:
        if constexpr (N == 8 && M == 8)
        {
            auto abs_diff{[](std::size_t x, std::size_t y) -> std::size_t
                          {
                              return (x >= y) ? x - y : y - x;
                          }};
            std::size_t diff_colum{abs_diff(s.m, t.m)};
            if (diff_colum == 2)
            {
                if (s.m > t.m)
                {
                    queen_side_castle(s, Coord(s.n, t.m - 2), val_mes);
                }
                else
                {
                    king_side_castle(s, Coord(s.n, t.m + 1), val_mes);
                }
            }
            else
            {
                move_classic_jump(s, t, val_mes);
            }
        }
        else
        {
            move_classic_jump(s, t, val_mes);
\end{minted}
\end{frame}

\begin{frame}[fragile]
\hyperlink{conclusion}{Retour au début}\newline
\begin{minted}[obeytabs=true,tabsize=2,linenos=true,breaklines,firstnumber=777]{cpp}
        }
        break;
    case TypePiece::KNIGHT:
        move_classic_jump(s, t, val_mes);
        break;
    case TypePiece::BISHOP:
        move_classic_slide(
            s, t,
            &check_path_diagonal_1_instance<N, M>,
            val_mes);
        break;
    case TypePiece::ROOK:
        move_classic_slide(
            s, t,
            &check_path_straight_1_instance<N, M>,
            val_mes);
        break;
    case TypePiece::QUEEN:
        move_classic_slide(
            s, t,
            &check_path_queen_1_instance<N, M>,
            val_mes);
        break;
    case TypePiece::PAWN:
        move_pawn(s, t, val_mes);
        return;
    case TypePiece::EMPTY:
    default:
        break;
    }
\end{minted}
\end{frame}

\begin{frame}[fragile]
\hyperlink{conclusion}{Retour au début}\newline
\begin{minted}[obeytabs=true,tabsize=2,linenos=true,breaklines,firstnumber=807]{cpp}
    m_ep = std::nullopt;
}

template <std::size_t N, std::size_t M>
CONSTEXPR void Board<N, M>::move_split(Coord const &s,
                                       Coord const &t1,
                                       Coord const &t2)
{
    TypePiece piece{(*this)(s.n, s.m).get_type()};
    switch (piece)
    {
    case TypePiece::KING:
    case TypePiece::KNIGHT:
        move_split_jump(s, t1, t2);
        break;
    case TypePiece::BISHOP:
        move_split_slide(s, t1, t2, &check_path_diagonal_1_instance<N, M>);
        break;
    case TypePiece::ROOK:
        move_split_slide(s, t1, t2, &check_path_straight_1_instance<N, M>);
        break;
    case TypePiece::QUEEN:
        move_split_slide(s, t1, t2, &check_path_queen_1_instance<N, M>);
        break;
    case TypePiece::PAWN:
    case TypePiece::EMPTY:
    default:
        break;
    }
    m_ep = std::nullopt;
\end{minted}
\end{frame}

\begin{frame}[fragile]
\hyperlink{conclusion}{Retour au début}\newline
\begin{minted}[obeytabs=true,tabsize=2,linenos=true,breaklines,firstnumber=837]{cpp}
}

template <std::size_t N, std::size_t M>
CONSTEXPR void Board<N, M>::move_merge(Coord const &s1,
                                       Coord const &s2,
                                       Coord const &t)
{
    TypePiece piece1{(*this)(s1.n, s1.m).get_type()};
    TypePiece piece2{(*this)(s2.n, s2.m).get_type()};
    if (piece1 != piece2)
    {
        return;
    }
    switch (piece1)
    {
    case TypePiece::KING:
    case TypePiece::KNIGHT:
        move_merge_jump(s1, s2, t);
        break;
    case TypePiece::BISHOP:
        move_merge_slide(s1, s2, t, &check_path_diagonal_1_instance<N, M>);
        break;
    case TypePiece::ROOK:
        move_merge_slide(s1, s2, t, &check_path_straight_1_instance<N, M>);
        break;
    case TypePiece::QUEEN:
        move_merge_slide(s1, s2, t, &check_path_queen_1_instance<N, M>);
        break;
    case TypePiece::PAWN:
    case TypePiece::EMPTY:
\end{minted}
\end{frame}

\begin{frame}[fragile]
\hyperlink{conclusion}{Retour au début}\newline
\begin{minted}[obeytabs=true,tabsize=2,linenos=true,breaklines,firstnumber=867]{cpp}
    default:
        break;
    }
    m_ep = std::nullopt;
}

template <std::size_t N, std::size_t M>
CONSTEXPR void Board<N, M>::move(Move const &movement, std::optional<bool> val_mes)
{
    if (std::empty(m_board))
    {
        throw std::runtime_error("Le plateau est vide impossible de réaliser l'opération move");
    }
    switch (movement.type)
    {
    case TypeMove::NORMAL:
        move_classic(movement.normal.src, movement.normal.arv, val_mes);
        break;
    case TypeMove::SPLIT:
        move_split(movement.split.src,
                   movement.split.arv1,
                   movement.split.arv2);
        break;
    case TypeMove::MERGE:
        move_merge(movement.merge.src1,
                   movement.merge.src2,
                   movement.merge.arv);
        break;
    case TypeMove::PROMOTE:
        move_promotion(movement, val_mes);
\end{minted}
\end{frame}

\begin{frame}[fragile]
\hyperlink{conclusion}{Retour au début}\newline
\begin{minted}[obeytabs=true,tabsize=2,linenos=true,breaklines,firstnumber=897]{cpp}
        break;
    default:
        break;
    }
    update_board_classic();
    //update_board();

\end{minted}
\end{frame}

\begin{frame}[fragile]
\label{getprobamove.tpp}
\frametitle{get-proba-move.tpp}
\hyperlink{conclusion}{Retour au début}\newline
\begin{minted}[obeytabs=true,tabsize=2,linenos=true, breaklines]{cpp}
// Lib standard
#include <cassert>
#include <utility>
#include <functional>
#include <cmath>
#include <optional>
#include <stdexcept>

// Inclusion projet
#include <Piece.hpp>
#include <TypePiece.hpp>
#include <Constexpr.hpp>
#include "Board.hpp"
#include <math_utility.hpp>

template <std::size_t N, std::size_t M>
CONSTEXPR double
Board<N, M>::get_proba_mesure(Coord const &p) const noexcept
{
    double res{0};
    std::size_t size_board = std::size(m_board);
    for (std::size_t i{0}; i < size_board; i++)
    {
        if (m_board[i].first[offset(p.n, p.m)])
        {
            res += std::pow(std::abs(m_board[i].second), 2);
\end{minted}
\end{frame}

\begin{frame}[fragile]
\hyperlink{conclusion}{Retour au début}\newline
\begin{minted}[obeytabs=true,tabsize=2,linenos=true,breaklines,firstnumber=27]{cpp}
        }
    }
    return res;
}

template <std::size_t N, std::size_t M>
CONSTEXPR double
Board<N, M>::get_proba_mesure_capture_slide(
    Coord const &s,
    Coord const &t,
    std::function<bool(Board<N, M> const &,
                       Coord const &,
                       Coord const &,
                       std::size_t,
                       std::optional<Coord>)>
        check_path) const noexcept
{
    double res{0};
    std::size_t position = offset(s.n, s.m);
    std::size_t size_board = std::size(m_board);
    for (std::size_t i{0}; i < size_board; i++)
    {
        if (m_board[i].first[position] &&
            check_path(*this, s, t, i, std::nullopt))
        {
            res += std::pow(std::abs(m_board[i].second), 2);
        }
    }
    return res;
}
\end{minted}
\end{frame}

\begin{frame}[fragile]
\hyperlink{conclusion}{Retour au début}\newline
\begin{minted}[obeytabs=true,tabsize=2,linenos=true,breaklines,firstnumber=57]{cpp}

template <std::size_t N, std::size_t M>
CONSTEXPR double
Board<N, M>::get_proba_mesure_castle(
    Coord const &king,
    Coord const &rook) const noexcept
{
    double res{0};
    std::size_t size_board = std::size(m_board);
    std::size_t position_king = offset(king.n, king.m);
    std::size_t position_rook = offset(rook.n, rook.m);
    for (std::size_t i{0}; i < size_board; i++)
    {
        if (m_board[i].first[position_king] &&
            m_board[i].first[position_rook] &&
            check_path_straight_1_instance(*this, king, rook, i, std::nullopt))
        {
            res += std::pow(std::abs(m_board[i].second), 2);
        }
    }
    return res;
}

template <std::size_t N, std::size_t M>
CONSTEXPR double Board<N, M>::get_proba_move(Move const &move) const noexcept
{
    Coord s;
    Coord t;
    switch (move.type)
    {
\end{minted}
\end{frame}

\begin{frame}[fragile]
\hyperlink{conclusion}{Retour au début}\newline
\begin{minted}[obeytabs=true,tabsize=2,linenos=true,breaklines,firstnumber=87]{cpp}
    case TypeMove::NORMAL:
        s = move.normal.src;
        t = move.normal.arv;
        break;
    case TypeMove::PROMOTE:
        s = move.promote.src;
        t = move.promote.arv;
        break;
    case TypeMove::MERGE:
        return 1.;
    case TypeMove::SPLIT:
        return 1.;
    default:
        return 1.;
    }
    TypePiece piece{(*this)(s.n, s.m).get_type()};
    TypePiece piece_target{(*this)(t.n, t.m).get_type()};
    if((piece == piece_target) && (*this)(s.n, s.m).same_color((*this)(t.n, t.m)))
    {
        return 1.;
    }
    switch (piece)
    {
    case TypePiece::KING:
        if constexpr (N == 8 && M == 8)
        {
            auto abs_diff{[](std::size_t x, std::size_t y) -> std::size_t
                          {
                              return (x >= y) ? x - y : y - x;
                          }};
\end{minted}
\end{frame}

\begin{frame}[fragile]
\hyperlink{conclusion}{Retour au début}\newline
\begin{minted}[obeytabs=true,tabsize=2,linenos=true,breaklines,firstnumber=117]{cpp}
            std::size_t diff_colum{abs_diff(s.m, t.m)};
            if (diff_colum == 2)
            {
                if (s.m > t.m)
                {
                    return get_proba_mesure_castle(s, Coord(s.n, t.m - 2));
                }
                else
                {
                    return get_proba_mesure_castle(s, Coord(s.n, t.m + 1));
                }
            }
            else
            {
                if (piece_target != TypePiece::EMPTY)
                {
                    if ((*this)(s.n, s.m).same_color((*this)(t.n, t.m)))
                    {
                        return 1. - get_proba_mesure(t);
                    }
                    else
                    {
                        return get_proba_mesure(s);
                    }
                }
                else
                {
                    return 1.;
                }
            }
\end{minted}
\end{frame}

\begin{frame}[fragile]
\hyperlink{conclusion}{Retour au début}\newline
\begin{minted}[obeytabs=true,tabsize=2,linenos=true,breaklines,firstnumber=147]{cpp}
        }
        else
        {
            if (piece_target != TypePiece::EMPTY)
            {
                if ((*this)(s.n, s.m).same_color((*this)(t.n, t.m)))
                {
                    return 1. - get_proba_mesure(t);
                }
                else
                {
                    return get_proba_mesure(s);
                }
            }
            else
            {
                return 1.;
            }
        }
        break;
    case TypePiece::KNIGHT:
        if (piece_target != TypePiece::EMPTY)
        {
            if ((*this)(s.n, s.m).same_color((*this)(t.n, t.m)))
            {
                return 1. - get_proba_mesure(t);
            }
            else
            {
                return get_proba_mesure(s);
\end{minted}
\end{frame}

\begin{frame}[fragile]
\hyperlink{conclusion}{Retour au début}\newline
\begin{minted}[obeytabs=true,tabsize=2,linenos=true,breaklines,firstnumber=177]{cpp}
            }
        }
        else
        {
            return 1.;
        }

        break;
    case TypePiece::BISHOP:
        if ((*this)(t.n, t.m).get_type() == TypePiece::EMPTY)
        {
            return 1.;
        }
        else
        {
            if (!(*this)(t.n, t.m).same_color((*this)(s.n, s.m)))
            {
                return get_proba_mesure_capture_slide(
                    s, t,
                    &check_path_diagonal_1_instance<N, M>);
            }
            else
            {
                return 1. - get_proba_mesure(t);
            }
        }
        break;
    case TypePiece::ROOK:
        if ((*this)(t.n, t.m).get_type() == TypePiece::EMPTY)
        {
\end{minted}
\end{frame}

\begin{frame}[fragile]
\hyperlink{conclusion}{Retour au début}\newline
\begin{minted}[obeytabs=true,tabsize=2,linenos=true,breaklines,firstnumber=207]{cpp}
            return 1.;
        }
        else
        {
            if (!(*this)(t.n, t.m).same_color((*this)(s.n, s.m)))
            {
                return get_proba_mesure_capture_slide(
                    s, t,
                    &check_path_straight_1_instance<N, M>);
            }
            else
            {
                return 1. - get_proba_mesure(t);
            }
        }
        break;
    case TypePiece::QUEEN:
       if ((*this)(t.n, t.m).get_type() == TypePiece::EMPTY)
        {
            return 1.;
        }
        else
        {
            if (!(*this)(t.n, t.m).same_color((*this)(s.n, s.m)))
            {
                return get_proba_mesure_capture_slide(
                    s, t,
                    &check_path_queen_1_instance<N, M>);
            }
            else
\end{minted}
\end{frame}

\begin{frame}[fragile]
\hyperlink{conclusion}{Retour au début}\newline
\begin{minted}[obeytabs=true,tabsize=2,linenos=true,breaklines,firstnumber=237]{cpp}
            {
                return 1. - get_proba_mesure(t);
            }
        }
        break;
    case TypePiece::PAWN:
        if ((*this)(t.n, t.m).get_type() == TypePiece::EMPTY ||
            ((*this)(t.n, t.m).get_type() == (*this)(s.n, s.m).get_type() &&
             (*this)(t.n, t.m).get_color() == (*this)(s.n, s.m).get_color()))
        {
            return 1.;
        }
        else
        {
            return 1. - get_proba_mesure(t);
        }
    case TypePiece::EMPTY:
    default:
        break;
    }
    return 1.;

\end{minted}
\end{frame}

\begin{frame}[fragile]
\label{Piece.hpp}
\frametitle{Piece.hpp}
\hyperlink{conclusion}{Retour au début}\newline
\begin{minted}[obeytabs=true,tabsize=2,linenos=true, breaklines]{cpp}
#ifndef PIECE_HPP
#define PIECE_HPP

#include <Color.hpp>
#include <Coord.hpp>
#include <concepts>
#include <vector>
#include <forward_list>
#include <Constexpr.hpp>
#include <Move.hpp>
#include "TypePiece.hpp"

template <std::size_t N, std::size_t M>
class Board;

/**
 * @brief Class conservant le type et la couleur d'une pièce.
 * Elle permet aussi de recuperer ses mouvement possible
 */
class Piece
{
public:
    constexpr inline Piece() noexcept;

    /**
     * @brief Construit une piece à l'aide de son type et de sa couleur
\end{minted}
\end{frame}

\begin{frame}[fragile]
\hyperlink{conclusion}{Retour au début}\newline
\begin{minted}[obeytabs=true,tabsize=2,linenos=true,breaklines,firstnumber=27]{cpp}
     *
     * @param[in] piece Le type de la piece
     * @param[in] color La couleur de la piece
     */
    constexpr inline Piece(TypePiece piece, Color color) noexcept;
    constexpr inline Piece(Piece const &) = default;
    constexpr inline Piece &operator=(Piece const &) = default;

    /**
     * @brief Deplacement d'un objet piece vers un autre objet piece
     */
    constexpr inline Piece(Piece &&) = default;
    constexpr inline Piece &operator=(Piece &&) = default;
    constexpr inline ~Piece() = default;

    /**
     * @brief Renvoie le type de la piece
     *
     * @return Une enumération TypePiece
     */
    constexpr inline TypePiece get_type() const noexcept;

    /**
     * @brief Renvoie la couleur d'un piece
     *
     * @return Color La couleur de la piece
     */
    constexpr inline Color get_color() const noexcept;

    /**
\end{minted}
\end{frame}

\begin{frame}[fragile]
\hyperlink{conclusion}{Retour au début}\newline
\begin{minted}[obeytabs=true,tabsize=2,linenos=true,breaklines,firstnumber=57]{cpp}
     * @brief Récupère la liste des mouvements classics pour une piece
     *
     * @tparam N Le nombre de ligne du plateau
     * @tparam M Le nombre de colonne du plateau
     * @param[in] board Le plateau
     * @param[in] pos la position de la piece sur le plateau
     * @warning Le resultat n'est pas garanti si la position ne concorde pas
     * avec la position de la piece sur le plateau
     * @return std::forward_list<Coord> La liste de coordonnées des positions
     * d'arrivées de chaque mouvement possible
     */
    template <std::size_t N, std::size_t M>
    CONSTEXPR std::forward_list<Coord>
    get_list_normal_move(Board<N, M> const &board,
                         Coord const &pos) const noexcept;

    /**
     * @brief Récupère la liste des mouvements split pour une piece
     *
     * @tparam N Le nombre de ligne du plateau
     * @tparam M Le nombre de colonne du plateau
     * @param[in] board Le plateau
     * @param[in] pos la position de la piece sur le plateau
     * @warning Le resultat n'est pas garanti si la position ne concorde pas
     * avec la position de la piece sur le plateau
     * @warning Ne renvoie qu'une liste de coordonnées uniques car
     * l'ensemble des mouvements possibles est toutes les couples
     * de coordonnées de la liste
     * @return std::forward_list<Coord> La liste de coordonnées
     * des positions d'arrivées possible sachant que chaque élément
\end{minted}
\end{frame}

\begin{frame}[fragile]
\hyperlink{conclusion}{Retour au début}\newline
\begin{minted}[obeytabs=true,tabsize=2,linenos=true,breaklines,firstnumber=87]{cpp}
     * représente une seule des deux cases d'arrivées d'un split move
     */
    template <std::size_t N, std::size_t M>
    CONSTEXPR std::forward_list<Coord>
    get_list_split_move(Board<N, M> const &board,
                        Coord const &pos) const noexcept;

    /**
     * @brief Teste si la piece est blanche
     *
     * @return bool true si la piece est blanche, false sinon
     */
    constexpr inline bool is_white() const noexcept;

    /**
     * @brief Teste si la piece est noire
     *
     * @return bool true si la piece est noire, false sinon
     */
    constexpr inline bool is_black() const noexcept;

    /**
     * @brief Teste si l'autre piece est de la meme couleur
     *
     * @param[in] other L'autre pièce à comparer
     *
     * @return bool true si la piece est de la meme couleur, false sinon
     */
    constexpr inline bool same_color(Piece const &other) const noexcept;

\end{minted}
\end{frame}

\begin{frame}[fragile]
\hyperlink{conclusion}{Retour au début}\newline
\begin{minted}[obeytabs=true,tabsize=2,linenos=true,breaklines,firstnumber=117]{cpp}
    /**
     * @brief Teste si les mouvement de la piece sont si la
     * pièce est un pion un mouvement de promotion
     * 
     * @note peut etre utiliser sans verifier si la pièce est un pion
     * 
     * @param[in] board Le plateau
     * @param[in] pos La position de la pièce
     * @return true si le mouvement possible est un mouvement
     * de promotion
     * @return false sinon
     */
    template <std::size_t N, std::size_t M>
    CONSTEXPR bool
    check_if_use_move_promote(Board<N, M> const &board,
                              Coord const &pos) const noexcept;

    /**
     * @brief Renvoie la liste de toutes les promotions 
     * 
     * @warning Ne verifie pas la validité du mouvement,
     * de la position ou du type de la pièce
     * 
     * @param[in] board Le plateau
     * @param[in] pos La position du pion
     * @return La liste de tout les mouvements de promotion
     */
    template <std::size_t N, std::size_t M>
    CONSTEXPR std::forward_list<Move>
    get_list_promote(Board<N, M> const &board,
\end{minted}
\end{frame}

\begin{frame}[fragile]
\hyperlink{conclusion}{Retour au début}\newline
\begin{minted}[obeytabs=true,tabsize=2,linenos=true,breaklines,firstnumber=147]{cpp}
                     Coord const &pos) const noexcept;

private:
    /**
     * @brief Enumération utilisé dans le template des fonctions
     * qui récupère les listes de mouvement pour modifier leurs
     * comportement lors de la prise de pièce.
     */
    enum class Move_Mode
    {
        /**
         * @brief Dans le cas du mouvement classique on
         * autorise la prise de pièce.
         */
        NORMAL,

        /**
         * @brief Dans le cas du split on interdit la
         * prise de pièce.
         */
        SPLIT
    };

    /**
     * @brief renvoie la valeur absolue d'un soustraction sur des types size_t
     *
     * @param[in] x Variable de type size_t
     * @param[in] y Variable de type size_t
     * @return std::size_t Retourne l'opération |x - y|
     */
\end{minted}
\end{frame}

\begin{frame}[fragile]
\hyperlink{conclusion}{Retour au début}\newline
\begin{minted}[obeytabs=true,tabsize=2,linenos=true,breaklines,firstnumber=177]{cpp}
    constexpr inline static std::size_t
    abs_substracte(std::size_t x, std::size_t y) noexcept;

    /**
     * @brief renvoie la distance entre deux coordonnées
     * à l'aide de la norme 2
     *
     * @param[in] x Première coordonnée
     * @param[in] y Seconde coordonnée
     * @return double le résultat de la norme 1
     */
    constexpr inline static double
    norm(Coord const &x, Coord const &y) noexcept;

    /**
     * @brief Récupère la liste de mouvement d'une pièce pour
     * les déplacements indiqué de manière récursive
     * c'est à dire tout les mouvements infini (ex : diagonale du fou)
     *
     * @tparam MOVE Le type de mouvement utiliser entre normal
     * et split (split ne prend pas les mouvements de capture de pièce)
     * @tparam N Le nombre de ligne du plateau
     * @tparam M Le nombre de colonne du plateau
     * @tparam SIZE La taille de la liste des mouvements possible
     * @param[in] board Le plateau
     * @param[in] pos La position de la pièce
     * @param[in] list_move La liste de tout les mouvements possible
     * @return std::forward_list<Coord> La liste de coordonnées
     * des positions d'arrivées possible
     */
\end{minted}
\end{frame}

\begin{frame}[fragile]
\hyperlink{conclusion}{Retour au début}\newline
\begin{minted}[obeytabs=true,tabsize=2,linenos=true,breaklines,firstnumber=207]{cpp}
    template <Move_Mode MOVE, std::size_t N, std::size_t M, std::size_t SIZE>
    CONSTEXPR std::forward_list<Coord>
    get_list_move_rec(Board<N, M> const &board, Coord const &pos,
                      std::array<int, SIZE> const &list_move) const noexcept;

    /**
     * @brief Récupère la liste de mouvement du roi
     *
     * @warning Ne récupère les mouvements du roque que pour un tableau
     * de taille 8 x 8
     *
     * @warning ne verifie pas si c'est bien un roi
     *
     * @tparam MOVE Le type de mouvement utiliser entre normal
     * et split (split ne prend pas les mouvements de capture de pièce)
     * @tparam N Le nombre de ligne du plateau
     * @tparam M Le nombre de colonne du plateau
     * @param[in] board Le plateau
     * @param[in] pos La position du roi
     * @return std::forward_list<Coord> La liste de coordonnées
     * des positions d'arrivées possible pour le roi
     */
    template <Move_Mode MOVE, std::size_t N, std::size_t M>
    CONSTEXPR std::forward_list<Coord>
    get_list_move_king(Board<N, M> const &board,
                       Coord const &pos) const noexcept;

    /**
     * @brief Récupère la liste de mouvement du cavalier
     *
\end{minted}
\end{frame}

\begin{frame}[fragile]
\hyperlink{conclusion}{Retour au début}\newline
\begin{minted}[obeytabs=true,tabsize=2,linenos=true,breaklines,firstnumber=237]{cpp}
     * @warning ne verifie pas si c'est bien un cavalier
     *
     * @tparam MOVE Le type de mouvement utiliser entre normal
     * et split (split ne prend pas les mouvements de capture de pièce)
     * @tparam N Le nombre de ligne du plateau
     * @tparam M Le nombre de colonne du plateau
     * @param[in] board Le plateau
     * @param[in] pos La position du cavalier
     * @return std::forward_list<Coord> La liste de coordonnées
     * des positions d'arrivées possible pour le cavalier
     */
    template <Move_Mode MOVE, std::size_t N, std::size_t M>
    CONSTEXPR std::forward_list<Coord>
    get_list_move_knight(Board<N, M> const &board,
                         Coord const &pos) const noexcept;

    /**
     * @brief Récupère la liste de mouvement du fou
     *
     * @warning ne verifie pas si c'est bien un fou
     *
     * @tparam MOVE Le type de mouvement utiliser entre normal
     * et split (split ne prend pas les mouvements de capture de pièce)
     * @tparam N Le nombre de ligne du plateau
     * @tparam M Le nombre de colonne du plateau
     * @param[in] board Le plateau
     * @param[in] pos La position du fou
     * @return std::forward_list<Coord> La liste de coordonnées
     * des positions d'arrivées possible pour le fou
     */
\end{minted}
\end{frame}

\begin{frame}[fragile]
\hyperlink{conclusion}{Retour au début}\newline
\begin{minted}[obeytabs=true,tabsize=2,linenos=true,breaklines,firstnumber=267]{cpp}
    template <Move_Mode MOVE, std::size_t N, std::size_t M>
    CONSTEXPR std::forward_list<Coord>
    get_list_move_bishop(Board<N, M> const &board,
                         Coord const &pos) const noexcept;

    /**
     * @brief Récupère la liste de mouvement de la tour
     *
     * @warning ne verifie pas si c'est bien une tour
     *
     * @tparam MOVE Le type de mouvement utiliser entre normal
     * et split (split ne prend pas les mouvements de capture de pièce)
     * @tparam N Le nombre de ligne du plateau
     * @tparam M Le nombre de colonne du plateau
     * @param[in] board Le plateau
     * @param[in] pos La position de la tour
     * @return std::forward_list<Coord> La liste de coordonnées
     * des positions d'arrivées possible pour la tour
     */
    template <Move_Mode MOVE, std::size_t N, std::size_t M>
    CONSTEXPR std::forward_list<Coord>
    get_list_move_rook(Board<N, M> const &board,
                       Coord const &pos) const noexcept;

    /**
     * @brief Récupère la liste de mouvement de la dame
     *
     * @warning ne verifie pas si c'est bien une dame
     *
     * @tparam MOVE Le type de mouvement utiliser entre normal
\end{minted}
\end{frame}

\begin{frame}[fragile]
\hyperlink{conclusion}{Retour au début}\newline
\begin{minted}[obeytabs=true,tabsize=2,linenos=true,breaklines,firstnumber=297]{cpp}
     * et split (split ne prend pas les mouvements de capture de pièce)
     * @tparam N Le nombre de ligne du plateau
     * @tparam M Le nombre de colonne du plateau
     * @param[in] board Le plateau
     * @param[in] pos La position de la dame
     * @return std::forward_list<Coord> La liste de coordonnées
     * des positions d'arrivées possible pour la dame
     */
    template <Move_Mode MOVE, std::size_t N, std::size_t M>
    CONSTEXPR std::forward_list<Coord>
    get_list_move_queen(Board<N, M> const &board,
                        Coord const &pos) const noexcept;

    /**
     * @brief Récupère la liste de mouvement du pion
     *
     * @warning ne verifie pas si c'est bien un pion
     *
     * @tparam N Le nombre de ligne du plateau
     * @tparam M Le nombre de colonne du plateau
     * @param[in] board Le plateau
     * @param[in] pos La position du pion
     * @return std::forward_list<Coord> La liste de coordonnées
     * des positions d'arrivées possible pour le pion
     */
    template <std::size_t N, std::size_t M>
    CONSTEXPR std::forward_list<Coord>
    get_list_move_pawn(Board<N, M> const &board,
                       Coord const &pos) const noexcept;

\end{minted}
\end{frame}

\begin{frame}[fragile]
\hyperlink{conclusion}{Retour au début}\newline
\begin{minted}[obeytabs=true,tabsize=2,linenos=true,breaklines,firstnumber=327]{cpp}
    /**
     * @brief Récupère la liste de mouvement pour
     * n'importe qu'elle pièce
     *
     * @tparam MOVE Le type de mouvement utiliser entre normal
     * et split (split ne prend pas les mouvements de capture de pièce)
     * @tparam N Le nombre de ligne du plateau
     * @tparam M Le nombre de colonne du plateau
     * @param[in] board Le plateau
     * @param[in] pos La position de la pièce
     * @param[in] list_move La liste de tout les mouvements possible
     * @return std::forward_list<Coord> La liste de coordonnées
     * des positions d'arrivées possible
     */
    template <Move_Mode MOVE, std::size_t N, std::size_t M>
    CONSTEXPR std::forward_list<Coord>
    get_list_move(Board<N, M> const &board,
                  Coord const &pos) const noexcept;

    /**
     * @brief La couleur de la pièce (WHITE/BLACK)
     */
    Color m_color;

    /**
     * @brief Le type de la piece
     * ( KING / QUEEN / ROOK / BISHOP / KNIGHT / PAWN )
     */
    TypePiece m_type;
};
\end{minted}
\end{frame}

\begin{frame}[fragile]
\hyperlink{conclusion}{Retour au début}\newline
\begin{minted}[obeytabs=true,tabsize=2,linenos=true,breaklines,firstnumber=357]{cpp}

#include "Piece.tpp"
#include "get_list_move.tpp"

/**
 * @brief L'objet représentant le roi blanc
 */
constexpr Piece W_KING{TypePiece::KING, Color::WHITE};

/**
 * @brief L'objet représentant le roi noir
 */
constexpr Piece B_KING{TypePiece::KING, Color::BLACK};

/**
 * @brief L'objet représentant la renne blanche
 */
constexpr Piece W_QUEEN{TypePiece::QUEEN, Color::WHITE};

/**
 * @brief L'objet représentant la renne noire
 */
constexpr Piece B_QUEEN{TypePiece::QUEEN, Color::BLACK};

/**
 * @brief L'objet représentant le fou blanc
 */
constexpr Piece W_BISHOP{TypePiece::BISHOP, Color::WHITE};

/**
\end{minted}
\end{frame}

\begin{frame}[fragile]
\hyperlink{conclusion}{Retour au début}\newline
\begin{minted}[obeytabs=true,tabsize=2,linenos=true,breaklines,firstnumber=387]{cpp}
 * @brief L'objet représentant le fou noir
 */
constexpr Piece B_BISHOP{TypePiece::BISHOP, Color::BLACK};

/**
 * @brief L'objet représentant la tour blanche
 */
constexpr Piece W_ROOK{TypePiece::ROOK, Color::WHITE};

/**
 * @brief L'objet représentant la tour noire
 */
constexpr Piece B_ROOK{TypePiece::ROOK, Color::BLACK};

/**
 * @brief L'objet représentant le cavalier blanc
 */
constexpr Piece W_KNIGHT{TypePiece::KNIGHT, Color::WHITE};

/**
 * @brief L'objet représentant le cavalier noir
 */
constexpr Piece B_KNIGHT{TypePiece::KNIGHT, Color::BLACK};

/**
 * @brief L'objet représentant le pion blanc
 */
constexpr Piece W_PAWN{TypePiece::PAWN, Color::WHITE};

/**
\end{minted}
\end{frame}

\begin{frame}[fragile]
\hyperlink{conclusion}{Retour au début}\newline
\begin{minted}[obeytabs=true,tabsize=2,linenos=true,breaklines,firstnumber=417]{cpp}
 * @brief L'objet représentant le pion noir
 */
constexpr Piece B_PAWN{TypePiece::PAWN, Color::BLACK};


\end{minted}
\end{frame}

\begin{frame}[fragile]
\label{Piece.tpp}
\frametitle{Piece.tpp}
\hyperlink{conclusion}{Retour au début}\newline
\begin{minted}[obeytabs=true,tabsize=2,linenos=true, breaklines]{cpp}
#include "Piece.hpp"
#include <Coord.hpp>
#include <check_path.hpp>
#include <cmath>
#include <forward_list>
#include <observer_ptr.hpp>
#include <math_utility.hpp>
#include <Constexpr.hpp>
#include <Move.hpp>

#define POW2(x) ((x) * (x))

constexpr inline Piece::Piece() noexcept
    : m_color(Color::BLACK),
      m_type(TypePiece::EMPTY)
{
}

constexpr inline Piece::Piece(TypePiece piece, Color color) noexcept
    : m_color(color),
      m_type(piece)
{
}

constexpr inline bool Piece::is_white() const noexcept
{
\end{minted}
\end{frame}

\begin{frame}[fragile]
\hyperlink{conclusion}{Retour au début}\newline
\begin{minted}[obeytabs=true,tabsize=2,linenos=true,breaklines,firstnumber=27]{cpp}
    return m_color == Color::WHITE;
}

constexpr inline bool Piece::is_black() const noexcept
{
    return m_color == Color::BLACK;
}

constexpr inline bool
Piece::same_color(Piece const &other) const noexcept
{
    return m_color == other.m_color;
}

constexpr inline TypePiece Piece::get_type() const noexcept
{
    return m_type;
}

constexpr inline Color Piece::get_color() const noexcept
{
    return m_color;
}

constexpr inline std::size_t
Piece::abs_substracte(std::size_t x, std::size_t y) noexcept
{
    if (x >= y)
    {
        return x - y;
\end{minted}
\end{frame}

\begin{frame}[fragile]
\hyperlink{conclusion}{Retour au début}\newline
\begin{minted}[obeytabs=true,tabsize=2,linenos=true,breaklines,firstnumber=57]{cpp}
    }
    return y - x;
}

constexpr inline double
Piece::norm(Coord const &x, Coord const &y) noexcept
{
    return sqrt(
        static_cast<double>(POW2(abs_substracte(x.n, y.n)) +
                            POW2(abs_substracte(x.m, y.m))));
}

template <std::size_t N, std::size_t M>
CONSTEXPR bool
Piece::check_if_use_move_promote(Board<N, M> const &board,
                                 Coord const &pos) const noexcept
{
    if constexpr (N >= 2)
    {
        if (board(pos.n, pos.m).get_type() == TypePiece::PAWN)
        {
            auto sign_color{
                [](Color color) -> int
                {
                    return (color == Color::WHITE) ? -1 : 1;
                }};
            std::size_t required_line{(get_color() == Color::WHITE) ? 1 : N - 2};
            if (pos.n == required_line)
            {
                return true;
\end{minted}
\end{frame}

\begin{frame}[fragile]
\hyperlink{conclusion}{Retour au début}\newline
\begin{minted}[obeytabs=true,tabsize=2,linenos=true,breaklines,firstnumber=87]{cpp}
            }
        }
    }
    return false;

\end{minted}
\end{frame}

\begin{frame}[fragile]
\label{getlistmove.tpp}
\frametitle{get-list-move.tpp}
\hyperlink{conclusion}{Retour au début}\newline
\begin{minted}[obeytabs=true,tabsize=2,linenos=true, breaklines]{cpp}
#include "Piece.hpp"
#include <Coord.hpp>
#include <check_path.hpp>
#include <cmath>
#include <forward_list>
#include <observer_ptr.hpp>
#include <math_utility.hpp>
#include <Constexpr.hpp>
#include <Move.hpp>

template <Piece::Move_Mode MOVE, std::size_t N, std::size_t M>
CONSTEXPR std::forward_list<Coord>
Piece::get_list_move_king(Board<N, M> const &board,
                          Coord const &pos) const noexcept
{
    CONSTEXPR int P{M + 2};
    std::array<int, 8> const
        move_king{{-P - 1, -P, -P + 1, -1, 1, P - 1, P, P + 1}};
    std::size_t current_pos{board.offset(pos.n, pos.m)};
    int posBox{board.m_S_mailbox[current_pos]};
    std::forward_list<Coord> list_move{};
    for (int const e : move_king)
    {
        int arv{board.m_L_mailbox[posBox + e]};
        if (arv >= 0)
        {
\end{minted}
\end{frame}

\begin{frame}[fragile]
\hyperlink{conclusion}{Retour au début}\newline
\begin{minted}[obeytabs=true,tabsize=2,linenos=true,breaklines,firstnumber=27]{cpp}
            std::size_t n{arv / M}, m{arv % M};
            Piece const &arvPiece{board(n, m)};
            bool eval {false};
            if CONSTEXPR (MOVE == Move_Mode::NORMAL)
            {
                eval = arvPiece.get_type() == TypePiece::EMPTY ||
                       !same_color(arvPiece) ||
                       board.get_proba(Coord(n, m)) < 1. - EPSILON;
            }
#if !defined(KING_NO_SPLIT)
            else
            {
                eval = arvPiece.get_type() == TypePiece::EMPTY;
            }
#endif
            if (eval)
            {
                list_move.push_front(Coord(n, m));
            }
        }
    }
    if CONSTEXPR (N == 8 && M == 8 && MOVE == Move_Mode::NORMAL)
    {
        if (board.m_q_castle[static_cast<int>(get_color())] &&
            check_path_straight(board, pos, Coord(pos.n, pos.m - 4)))
        {
            list_move.push_front(Coord(pos.n, pos.m - 2));
        }
        if (board.m_k_castle[static_cast<int>(get_color())] &&
            check_path_straight(board, pos, Coord(pos.n, pos.m + 3)))
\end{minted}
\end{frame}

\begin{frame}[fragile]
\hyperlink{conclusion}{Retour au début}\newline
\begin{minted}[obeytabs=true,tabsize=2,linenos=true,breaklines,firstnumber=57]{cpp}
        {
            list_move.push_front(Coord(pos.n, pos.m + 2));
        }
    }
    return list_move;
}

template <Piece::Move_Mode MOVE, std::size_t N, std::size_t M>
CONSTEXPR std::forward_list<Coord>
Piece::get_list_move_knight(Board<N, M> const &board,
                            Coord const &pos) const noexcept
{
    CONSTEXPR int P{M + 2};
    std::array<int, 8> const
        move_knight{{-2 * P - 1, -P - 2, -2 * P + 1, -P + 2,
                     P - 2, 2 * P - 1, 2 * P + 1, P + 2}};
    std::size_t current_pos{board.offset(pos.n, pos.m)};
    int posBox{board.m_S_mailbox[current_pos]};
    std::forward_list<Coord> list_move{};
    for (int const e : move_knight)
    {
        int arv{board.m_L_mailbox[posBox + e]};
        if (arv >= 0)
        {
            std::size_t n{arv / M}, m{arv % M};
            Piece const &arvPiece{board(n, m)};
            bool eval;
            if CONSTEXPR (MOVE == Move_Mode::NORMAL)
            {
                eval = arvPiece.get_type() == TypePiece::EMPTY ||
\end{minted}
\end{frame}

\begin{frame}[fragile]
\hyperlink{conclusion}{Retour au début}\newline
\begin{minted}[obeytabs=true,tabsize=2,linenos=true,breaklines,firstnumber=87]{cpp}
                       !same_color(arvPiece) ||
                       board.get_proba(Coord(n, m)) < 1. - EPSILON;
            }
            else
            {
                eval = arvPiece.get_type() == TypePiece::EMPTY;
            }
            if (eval)
            {
                list_move.push_front(Coord(n, m));
            }
        }
    }
    return list_move;
}

template <Piece::Move_Mode MOVE,
          std::size_t N,
          std::size_t M,
          std::size_t SIZE>
CONSTEXPR std::forward_list<Coord>
Piece::get_list_move_rec(
    Board<N, M> const &board,
    Coord const &pos,
    std::array<int, SIZE> const &list_move) const noexcept
{
    std::size_t current_pos{board.offset(pos.n, pos.m)};
    int const posBox{board.m_S_mailbox[current_pos]};
    std::forward_list<Coord> move{};
    for (int const e : list_move)
\end{minted}
\end{frame}

\begin{frame}[fragile]
\hyperlink{conclusion}{Retour au début}\newline
\begin{minted}[obeytabs=true,tabsize=2,linenos=true,breaklines,firstnumber=117]{cpp}
    {
        int arv{board.m_L_mailbox[posBox + e]};
        int posBox_tmp{posBox};
        while (arv >= 0)
        {
            std::size_t n{arv / M}, m{arv % M};
            Piece const &arvPiece{board(n, m)};
            if CONSTEXPR (MOVE == Move_Mode::NORMAL)
            {
                if (arvPiece.get_type() == TypePiece::EMPTY ||
                    board.get_proba(Coord(n, m)) < 1. - EPSILON)
                {
                    move.push_front(Coord(n, m));
                }
                else if (!same_color(arvPiece))
                {
                    move.push_front(Coord(n, m));
                    break;
                }
                else
                {
                    break;
                }
            }
            else
            {
                if (arvPiece.get_type() == TypePiece::EMPTY ||
                    (arvPiece.get_type() == board(pos.n, pos.m).get_type() &&
                     arvPiece.get_color() == board(pos.n, pos.m).get_color()))
                {
\end{minted}
\end{frame}

\begin{frame}[fragile]
\hyperlink{conclusion}{Retour au début}\newline
\begin{minted}[obeytabs=true,tabsize=2,linenos=true,breaklines,firstnumber=147]{cpp}
                    move.push_front(Coord(n, m));
                }
                else
                {
                    break;
                }
            }
            posBox_tmp = board.m_S_mailbox[arv];
            arv = board.m_L_mailbox[posBox_tmp + e];
        }
    }
    return move;
}

template <Piece::Move_Mode MOVE, std::size_t N, std::size_t M>
CONSTEXPR std::forward_list<Coord>
Piece::get_list_move_bishop(Board<N, M> const &board,
                            Coord const &pos) const noexcept
{
    CONSTEXPR int P{M + 2};
    std::array<int, 4> const move_bishop{{-P - 1, -P + 1, P - 1, P + 1}};
    return get_list_move_rec<MOVE>(board, pos, move_bishop);
}

template <Piece::Move_Mode MOVE, std::size_t N, std::size_t M>
CONSTEXPR std::forward_list<Coord>
Piece::get_list_move_rook(Board<N, M> const &board,
                          Coord const &pos) const noexcept
{
    CONSTEXPR int P{M + 2};
\end{minted}
\end{frame}

\begin{frame}[fragile]
\hyperlink{conclusion}{Retour au début}\newline
\begin{minted}[obeytabs=true,tabsize=2,linenos=true,breaklines,firstnumber=177]{cpp}
    std::array<int, 4> const move_rook{{-P, 1, P, -1}};
    return get_list_move_rec<MOVE>(board, pos, move_rook);
}

template <Piece::Move_Mode MOVE, std::size_t N, std::size_t M>
CONSTEXPR std::forward_list<Coord>
Piece::get_list_move_queen(Board<N, M> const &board,
                           Coord const &pos) const noexcept
{
    CONSTEXPR int P{M + 2};
    std::array<int, 8> const
        move_queen{{-P - 1, -P, -P + 1, -1, 1, P - 1, P, P + 1}};
    return get_list_move_rec<MOVE>(board, pos, move_queen);
}

template <std::size_t N, std::size_t M>
CONSTEXPR std::forward_list<Coord>
Piece::get_list_move_pawn(Board<N, M> const &board,
                          Coord const &pos) const noexcept
{
    CONSTEXPR int P{M + 2};
    auto sign_color{
        [](Color color)
        {
            return (color == Color::WHITE) ? -1 : 1;
        }};
    int sign{sign_color(get_color())};
    std::array<int, 3> move_pawn{sign * P, sign * P - 1, sign * P + 1};
    std::forward_list<Coord> list_move{};
    std::size_t current_pos{board.offset(pos.n, pos.m)};
\end{minted}
\end{frame}

\begin{frame}[fragile]
\hyperlink{conclusion}{Retour au début}\newline
\begin{minted}[obeytabs=true,tabsize=2,linenos=true,breaklines,firstnumber=207]{cpp}
    int posBox{board.m_S_mailbox[current_pos]};
    int arv{board.m_L_mailbox[posBox + move_pawn[0]]};
    if (arv >= 0)
    {
        std::size_t n{arv / N}, m{arv % N};
        if (board(n, m).get_type() == TypePiece::EMPTY ||
            board.get_proba(Coord(n, m)) < 1. - EPSILON)
        {
            list_move.push_front(Coord(n, m));
            if (((pos.n == 1 && get_color() == Color::BLACK) ||
                 (pos.n == N - 2 && get_color() == Color::WHITE)) &&
                (board(n + sign, m).get_type() == TypePiece::EMPTY ||
                 board.get_proba(Coord(n + sign, m)) < 1. - EPSILON))
            {
                list_move.push_front(Coord(n + sign, m));
            }
        }
    }
    for (std::size_t i{1}; i < 3; i++)
    {
        arv = board.m_L_mailbox[posBox + move_pawn[i]];
        if (arv >= 0)
        {
            std::size_t n{arv / N}, m{arv % N};
            if ((board(n, m).get_type() != TypePiece::EMPTY &&
                 !same_color(board(n, m))) ||
                (board.m_ep != std::nullopt &&
                 board.m_ep->n == pos.n &&
                 board.m_ep->m == pos.m))
            {
\end{minted}
\end{frame}

\begin{frame}[fragile]
\hyperlink{conclusion}{Retour au début}\newline
\begin{minted}[obeytabs=true,tabsize=2,linenos=true,breaklines,firstnumber=237]{cpp}
                list_move.push_front(Coord(n, m));
            }
        }
    }
    return list_move;
}

template <std::size_t N, std::size_t M>
CONSTEXPR std::forward_list<Move>
Piece::get_list_promote(Board<N, M> const &board, Coord const &pos) const noexcept
{
    std::forward_list<Coord> list_move{get_list_move_pawn(board, pos)};
    std::forward_list<Move> list_promote;
    std::array<TypePiece, 4> promote_choice{TypePiece::QUEEN,
                                            TypePiece::ROOK,
                                            TypePiece::BISHOP,
                                            TypePiece::KNIGHT};
    for (Coord const &e : list_move)
    {
        for (TypePiece p : promote_choice)
        {
            list_promote.push_front(Move_promote(pos, e, p));
        }
    }
    return list_promote;
}

template <Piece::Move_Mode MOVE, std::size_t N, std::size_t M>
CONSTEXPR std::forward_list<Coord>
Piece::get_list_move(Board<N, M> const &board, Coord const &pos) const noexcept
\end{minted}
\end{frame}

\begin{frame}[fragile]
\hyperlink{conclusion}{Retour au début}\newline
\begin{minted}[obeytabs=true,tabsize=2,linenos=true,breaklines,firstnumber=267]{cpp}
{
    switch (get_type())
    {
    case TypePiece::PAWN:
        if CONSTEXPR (MOVE == Move_Mode::NORMAL)
        {
            return get_list_move_pawn(board, pos);
        }
        else
        {
            return std::forward_list<Coord>{};
        }
    case TypePiece::KNIGHT:
        return get_list_move_knight<MOVE>(board, pos);
    case TypePiece::BISHOP:
        return get_list_move_bishop<MOVE>(board, pos);
    case TypePiece::ROOK:
        return get_list_move_rook<MOVE>(board, pos);
    case TypePiece::QUEEN:
        return get_list_move_queen<MOVE>(board, pos);
    case TypePiece::KING:
        return get_list_move_king<MOVE>(board, pos);
    case TypePiece::EMPTY:
    default: 
        return std::forward_list<Coord>{};
    }
}

template <std::size_t N, std::size_t M>
CONSTEXPR std::forward_list<Coord>
\end{minted}
\end{frame}

\begin{frame}[fragile]
\hyperlink{conclusion}{Retour au début}\newline
\begin{minted}[obeytabs=true,tabsize=2,linenos=true,breaklines,firstnumber=297]{cpp}
Piece::get_list_normal_move(Board<N, M> const &board,
                            Coord const &pos) const noexcept
{
    return get_list_move<Move_Mode::NORMAL>(board, pos);
}

template <std::size_t N, std::size_t M>
CONSTEXPR std::forward_list<Coord>
Piece::get_list_split_move(Board<N, M> const &board,
                           Coord const &pos) const noexcept
{
    return get_list_move<Move_Mode::SPLIT>(board, pos);

\end{minted}
\end{frame}

\begin{frame}[fragile]
\label{TypePiece.hpp}
\frametitle{TypePiece.hpp}
\hyperlink{conclusion}{Retour au début}\newline
\begin{minted}[obeytabs=true,tabsize=2,linenos=true, breaklines]{cpp}
#ifndef TYPE_PIECE_HPP
#define TYPE_PIECE_HPP

/**
 * @brief Enumère toutes les sortes de pièces
 */
enum class TypePiece
{
    EMPTY,
    PAWN,
    KNIGHT,
    BISHOP,
    ROOK,
    QUEEN,
    KING
};


\end{minted}
\end{frame}

\begin{frame}[fragile]
\label{Coord.hpp}
\frametitle{Coord.hpp}
\hyperlink{conclusion}{Retour au début}\newline
\begin{minted}[obeytabs=true,tabsize=2,linenos=true, breaklines]{cpp}
#ifndef COORD_HPP
#define COORD_HPP

#include <cstddef>

/**
 * @brief La structure qui représente une coordonnée
 * sur le plateau
 */
struct Coord
{
    Coord() = default;
    Coord(std::size_t i, std::size_t j) : n(i), m(j) {};
    Coord(Coord const&) = default;
    Coord& operator=(Coord const&) = default;
    Coord(Coord &&) = default;
    Coord& operator=(Coord &&) = default;


    /**
     * @brief l'indice sur les lignes
     */
    std::size_t n;

    /**
     * @brief l'indice sur les colonnes
\end{minted}
\end{frame}

\begin{frame}[fragile]
\hyperlink{conclusion}{Retour au début}\newline
\begin{minted}[obeytabs=true,tabsize=2,linenos=true,breaklines,firstnumber=27]{cpp}
     */
    std::size_t m;
};


struct Coord_hash
{
    std::size_t operator () (Coord const &coord) const;
};

bool operator==(Coord const &lhs, Coord const &rhs);


\end{minted}
\end{frame}

\begin{frame}[fragile]
\label{Coord.cpp}
\frametitle{Coord.cpp}
\hyperlink{conclusion}{Retour au début}\newline
\begin{minted}[obeytabs=true,tabsize=2,linenos=true, breaklines]{cpp}
#include "Coord.hpp"
#include <functional>

std::size_t Coord_hash::operator () (Coord const &coord) const
{
    std::size_t h1 = std::hash<std::size_t>()(coord.n);
    std::size_t h2 = std::hash<std::size_t>()(coord.m);

    return h1 ^ h2;
}


bool operator==(Coord const &lhs, Coord const &rhs)
{
    return lhs.n == rhs.n && lhs.m == rhs.m;

\end{minted}
\end{frame}

\begin{frame}[fragile]
\label{Move.hpp}
\frametitle{Move.hpp}
\hyperlink{conclusion}{Retour au début}\newline
\begin{minted}[obeytabs=true,tabsize=2,linenos=true, breaklines]{cpp}
#ifndef MOVE_HPP
#define MOVE_HPP

#include <Coord.hpp>
#include <TypePiece.hpp>

/**
 * @brief Enumération représentant tout les types de mouvements
 */
enum class TypeMove
{
  /**
   * @brief Représente le mouvement classic
   */
  NORMAL,

  /**
   * @brief Réprésente le mouvement de split de deux pièces
   */
  SPLIT,

  /**
   * @brief Représente le mouvement de fusion de deux pièces
   */
  MERGE,

\end{minted}
\end{frame}

\begin{frame}[fragile]
\hyperlink{conclusion}{Retour au début}\newline
\begin{minted}[obeytabs=true,tabsize=2,linenos=true,breaklines,firstnumber=27]{cpp}
  /**
   * @brief Représente le mouvement de promotion du pion
   *
   */
  PROMOTE
};

/**
 * @brief Stoque tout les mouvements possibles
 */
struct Move
{
  Move() = default;
  Move(Move const &) = default;
  Move &operator=(Move const &) = default;
  Move(Move &&) = default;
  Move &operator=(Move &&) = default;
  /**
   * @brief Indique quel est le mouvement stoqué par l'union
   */
  TypeMove type;

  /**
   * @brief Stoque au choix l'un des trois mouvements possible
   */
  union
  {
    /**
     * @brief Stoque les coordonnées pour un mouvement classic
     */
\end{minted}
\end{frame}

\begin{frame}[fragile]
\hyperlink{conclusion}{Retour au début}\newline
\begin{minted}[obeytabs=true,tabsize=2,linenos=true,breaklines,firstnumber=57]{cpp}
    struct
    {
      /**
       * @brief La coordonnée de depart
       */
      Coord src;

      /**
       * @brief La coordonnées d'arrivée
       */
      Coord arv;
    } normal;

    /**
     * @brief stoque l'ensemble des coordonnées pour un
     * mouvement splité
     */
    struct
    {
      /**
       * @brief La coordonnée de depart
       */
      Coord src;

      /**
       * @brief Une coordonnée d'arrivée
       */
      Coord arv1;

      /**
\end{minted}
\end{frame}

\begin{frame}[fragile]
\hyperlink{conclusion}{Retour au début}\newline
\begin{minted}[obeytabs=true,tabsize=2,linenos=true,breaklines,firstnumber=87]{cpp}
       * @brief Une seconde coordonnée d'arrivée
       *
       */
      Coord arv2;
    } split;

    /**
     * @brief  stoque l'ensemble des coordonnées pour un
     * mouvement de fusion
     */
    struct
    {
      /**
       * @brief La coordonnée de la première pièce à fusioner
       */
      Coord src1;

      /**
       * @brief La coordonnée de la seconde pièce à fusioner
       */
      Coord src2;

      /**
       * @brief La coordonnée de la position d'arrivée
       */
      Coord arv;
    } merge;

    /**
     * @brief stoque l'ensemble des coordonnées et la piece choisi
\end{minted}
\end{frame}

\begin{frame}[fragile]
\hyperlink{conclusion}{Retour au début}\newline
\begin{minted}[obeytabs=true,tabsize=2,linenos=true,breaklines,firstnumber=117]{cpp}
     * pour un mouvement de promotion
     */
    struct
    {
      /**
       * @brief La case de depart
       */
      Coord src;

      /**
       * @brief La case d'arrviée
       */
      Coord arv;

      /**
       * @brief La piece choisi pour le move
       */
      TypePiece piece;
    } promote;
  };

  bool operator==(Move const &m) const noexcept;
};

constexpr inline Move Move_classic(Coord const &src, Coord const &arv)
{
  Move move;
  move.type = TypeMove::NORMAL;
  move.normal.src = src;
  move.normal.arv = arv;
\end{minted}
\end{frame}

\begin{frame}[fragile]
\hyperlink{conclusion}{Retour au début}\newline
\begin{minted}[obeytabs=true,tabsize=2,linenos=true,breaklines,firstnumber=147]{cpp}
  return move;
}

constexpr inline Move Move_split(Coord const &src, Coord const &arv1, Coord const &arv2)
{
  Move move;
  move.type = TypeMove::SPLIT;
  move.split.src = src;
  move.split.arv1 = arv1;
  move.split.arv2 = arv2;
  return move;
}

constexpr inline Move Move_merge(Coord const &src1, Coord const &src2, Coord const &arv)
{
  Move move;
  move.type = TypeMove::MERGE;
  move.merge.src1 = src1;
  move.merge.src2 = src2;
  move.merge.arv = arv;
  return move;
}

constexpr inline Move Move_promote(Coord const &src, Coord const &arv, TypePiece promotion)
{
  Move move;
  move.type = TypeMove::PROMOTE;
  move.promote.src = src;
  move.promote.arv = arv;
  move.promote.piece = promotion;
\end{minted}
\end{frame}

\begin{frame}[fragile]
\hyperlink{conclusion}{Retour au début}\newline
\begin{minted}[obeytabs=true,tabsize=2,linenos=true,breaklines,firstnumber=177]{cpp}
  return move;
}


\end{minted}
\end{frame}

\begin{frame}[fragile]
\label{Move.cpp}
\frametitle{Move.cpp}
\hyperlink{conclusion}{Retour au début}\newline
\begin{minted}[obeytabs=true,tabsize=2,linenos=true, breaklines]{cpp}
#include <Coord.hpp>
#include <TypePiece.hpp>
#include "Move.hpp"

bool Move::operator==(Move const &m) const noexcept
{
  if (type == m.type)
  {
    switch (type)
    {
    case TypeMove::NORMAL:
      return normal.src == m.normal.src &&
             normal.arv == m.normal.arv;
    case TypeMove::SPLIT:
      return split.src == m.split.src &&
             split.arv1 == m.split.arv1 &&
             split.arv2 == m.split.arv2;
    case TypeMove::MERGE:
      return merge.src1 == m.merge.src1 &&
             merge.src2 == m.merge.src2 &&
             merge.arv == m.merge.arv;
    case TypeMove::PROMOTE:
      return promote.src == m.promote.src &&
             promote.arv == m.promote.arv &&
             promote.piece == m.promote.piece;
    default:
\end{minted}
\end{frame}

\begin{frame}[fragile]
\hyperlink{conclusion}{Retour au début}\newline
\begin{minted}[obeytabs=true,tabsize=2,linenos=true,breaklines,firstnumber=27]{cpp}
      return false;
    }
  }
  return false;

\end{minted}
\end{frame}

\begin{frame}[fragile]
\label{Qubit.hpp}
\frametitle{Qubit.hpp}
\hyperlink{conclusion}{Retour au début}\newline
\begin{minted}[obeytabs=true,tabsize=2,linenos=true, breaklines]{cpp}
#ifndef QUBIT_HPP
#define QUBIT_HPP

#include <cstddef>
#include <complex>
#include <array>
#include <CMatrix.hpp>
#include <Constexpr.hpp>

/**
 * @brief Représente un qubit
 * 
 * @tparam N La taille du Qubit
 */
template <std::size_t N>
class Qubit final
{
public:
    // Constructeur
    CONSTEXPR Qubit() = default;

    /**
     * @brief Construit un nouveau qubit à l'aide d'un tableau de booléen
     *
     * @param data le tableau de n booléen utilisé pour
     * construire un n-qubit ex : |10>
\end{minted}
\end{frame}

\begin{frame}[fragile]
\hyperlink{conclusion}{Retour au début}\newline
\begin{minted}[obeytabs=true,tabsize=2,linenos=true,breaklines,firstnumber=27]{cpp}
     */
    CONSTEXPR Qubit(std::array<bool, N> const &data);
    CONSTEXPR Qubit(std::array<std::complex<double>, _2POW(N)> &&init_list);

    // Copie
    CONSTEXPR Qubit(Qubit const &) = delete;
    CONSTEXPR Qubit &operator=(Qubit const &) = delete;

    // Mouvement
    CONSTEXPR Qubit(Qubit &&) = delete;
    CONSTEXPR Qubit &operator=(Qubit &&) = delete;

    // Destructeur
    CONSTEXPR ~Qubit() = default;

    template <std::size_t M>
    friend CONSTEXPR Qubit<M> operator*(
        CMatrix<_2POW(M)> const &lhs,
        Qubit<M> const &rhs);

    template <std::size_t M>
    friend std::ostream &operator<<(
        std::ostream &out,
        Qubit<M> const &qubit);

    template <std::size_t M>
    friend CONSTEXPR std::array<
        std::pair<
            std::array<bool, M>,
            std::complex<double>>,
\end{minted}
\end{frame}

\begin{frame}[fragile]
\hyperlink{conclusion}{Retour au début}\newline
\begin{minted}[obeytabs=true,tabsize=2,linenos=true,breaklines,firstnumber=57]{cpp}
        2>
    qubitToArray(Qubit<M> const &qubit);

private:
    std::array<std::complex<double>, _2POW(N)> m_data;
};

/**
 * @brief Transforme un qubit dans sa représentation sous forme de
 * vecteur en sa représentation classique.
 * Si le qubit contient deux cases non nuls, c'est une combinaisons
 * linéaires de deux qubits que l'on peut écrire avec la forme classique,
 * c'est-à-dire avec des "0" et des "1".
 * @warning Cette fonction ne permet de faire la transformation que
 * pour au plus deux cases non nuls dans le qubits
 * @tparam N La taille du qubit
 * @param[in] qubit Le qubit a transformer
 * @return Dans chaque cases du tableau, on a le qubits
 * sous sa représentation classique que l'on modélise par
 * un tableau de booléen, et un complexe qui est la scalaire.
 */
template <std::size_t N>
CONSTEXPR std::array<std::pair<std::array<bool, N>, std::complex<double>>, 2>
qubitToArray(Qubit<N> const &qubit);

/**
 * @brief Opérateur du produit entre un qubit et une matrice
 * @warning le produit n'est pas associatif
 * @tparam M La taille du qubit, la matrice est de taille 2^M
 * @param[in] lhs La matrice carré complexe  de taille 2^M
\end{minted}
\end{frame}

\begin{frame}[fragile]
\hyperlink{conclusion}{Retour au début}\newline
\begin{minted}[obeytabs=true,tabsize=2,linenos=true,breaklines,firstnumber=87]{cpp}
 * @param[in] rhs Le M-qubit
 * @return Qubit<M> renvoie le qubit resultant du produit
 */
template <std::size_t N>
CONSTEXPR Qubit<N> operator*(
    CMatrix<_2POW(N)> const &lhs,
    Qubit<N> const &rhs);

/**
 * @brief Fonction d'affichage des qubit sous forme de
 * liste de nombre complexe
 *
 * @tparam M La taille du qubit
 * @param[out] out Le flux de destination
 * @param[in] qubit Le qubit à afficher
 * @return std::ostream& le flux de destination
 */
template <std::size_t N>
std::ostream &operator<<(std::ostream &out, Qubit<N> const &qubit);

#include "Qubit.tpp"


\end{minted}
\end{frame}

\begin{frame}[fragile]
\label{Qubit.tpp}
\frametitle{Qubit.tpp}
\hyperlink{conclusion}{Retour au début}\newline
\begin{minted}[obeytabs=true,tabsize=2,linenos=true, breaklines]{cpp}
#include "Qubit.hpp"
#include <Matrix.hpp>
#include <algorithm>
#include <numeric>
#include <math_utility.hpp>
#include <Constexpr.hpp>
#include <Coord.hpp>

template <std::size_t N>
CONSTEXPR Qubit<N>::Qubit(
    std::array<std::complex<double>, _2POW(N)> &&init_list)
    : m_data(std::move(init_list))
{
}

template <std::size_t N>
CONSTEXPR Qubit<N>::Qubit(std::array<bool, N> const &data) : m_data()
{

    auto sum{[](std::array<bool, N> const &tab) -> int
             {
                 int acc{0};
                 for (std::size_t i{0}; i < std::size(tab); i++)
                 {
                     acc += _2POW(i) * tab[std::size(tab) - i - 1];
                 }
\end{minted}
\end{frame}

\begin{frame}[fragile]
\hyperlink{conclusion}{Retour au début}\newline
\begin{minted}[obeytabs=true,tabsize=2,linenos=true,breaklines,firstnumber=27]{cpp}
                 return acc;
             }};
    m_data[sum(data)] = 1;
}

template <std::size_t N>
CONSTEXPR Qubit<N> operator*(
    CMatrix<_2POW(N)> const &lhs,
    Qubit<N> const &rhs)
{
    std::array<std::complex<double>, _2POW(N)> tab{};
    for (std::size_t i{0}; i < lhs.numberLines(); i++)
    {
        for (std::size_t j{0}; j < lhs.numberColumns(); j++)
        {
            tab[i] += lhs(i, j) * rhs.m_data[j];
        }
    }

    return Qubit<N>(std::move(tab));
}

template <std::size_t N>
std::ostream &operator<<(
    std::ostream &out,
    Qubit<N> const &qubit)
{
    for (std::size_t i{0}; i < _2POW(N); i++)
    {
        out << qubit.m_data[i].real() << "+" << qubit.m_data[i].imag() << "i  ";
\end{minted}
\end{frame}

\begin{frame}[fragile]
\hyperlink{conclusion}{Retour au début}\newline
\begin{minted}[obeytabs=true,tabsize=2,linenos=true,breaklines,firstnumber=57]{cpp}
    }
    return out;
}

template <std::size_t N>
CONSTEXPR std::array<
    std::pair<std::array<bool, N>,
              std::complex<double>>,
    2>
qubitToArray(Qubit<N> const &qubit)
{
    std::array<
        std::pair<std::array<bool, N>,
                  std::complex<double>>,
        2>
        tab{};
    bool b{true};
    std::size_t compteur{N};
    std::size_t pow = _2POW(N);
    for (std::size_t i{0}; i < pow; i++)
    {
        using namespace std::complex_literals;
        if (complex_equal(qubit.m_data[i], 0i))
        {
            if (i == pow - 1)
            {
                return tab;
            }
            else
            {
\end{minted}
\end{frame}

\begin{frame}[fragile]
\hyperlink{conclusion}{Retour au début}\newline
\begin{minted}[obeytabs=true,tabsize=2,linenos=true,breaklines,firstnumber=87]{cpp}
                if (b)
                {
                    while (tab[1].first[compteur - 1])
                    {
                        tab[0].first[compteur - 1] = false;
                        tab[1].first[compteur - 1] = false;
                        compteur--;
                    }
                    tab[0].first[compteur - 1] = true;
                    tab[1].first[compteur - 1] = true;
                    compteur = N;
                }
                else
                {
                    while (tab[1].first[compteur - 1])
                    {
                        tab[1].first[compteur - 1] = false;
                        compteur--;
                    }
                    tab[1].first[compteur - 1] = true;
                    compteur = N;
                }
            }
        }
        else
        {
            if (b)
            {
                tab[0].second = qubit.m_data[i];
                b = false;
\end{minted}
\end{frame}

\begin{frame}[fragile]
\hyperlink{conclusion}{Retour au début}\newline
\begin{minted}[obeytabs=true,tabsize=2,linenos=true,breaklines,firstnumber=117]{cpp}
                if (i != pow - 1)
                {
                    while (tab[1].first[compteur - 1])
                    {
                        tab[1].first[compteur - 1] = false;
                        compteur--;
                    }
                    tab[1].first[compteur - 1] = true;
                    compteur = N;
                }
            }
            else
            {
                tab[1].second = qubit.m_data[i];
                return tab;
            }
        }
    }
    return tab;

\end{minted}
\end{frame}

\begin{frame}[fragile]
\label{Random.hpp}
\frametitle{Random.hpp}
\hyperlink{conclusion}{Retour au début}\newline
\begin{minted}[obeytabs=true,tabsize=2,linenos=true, breaklines]{cpp}
#ifndef RANDOM_HPP
#define RANDOM_HPP

namespace rnd
{
  /**
   * @brief Renvoie un entier dans l'intervale [a, b]
   * 
   * @param a, b Entier pour spécifier la plage
   * @return int Un entier entre a et b inclu
   */
  int randint(int a, int b);

  /**
   * @brief Renvoie un réel dans l'intervale [a, b)
   * 
   * @param a, b Réel pour spécifier la plage
   * @return double Un réel entre a et b inclu
   */
  double randreal(double a, double b);
}



\end{minted}
\end{frame}

\begin{frame}[fragile]
\label{Random.cpp}
\frametitle{Random.cpp}
\hyperlink{conclusion}{Retour au début}\newline
\begin{minted}[obeytabs=true,tabsize=2,linenos=true, breaklines]{cpp}
#include <random>
#include "Random.hpp"

namespace
{
  std::random_device rd;
}

namespace rnd
{
  int randint(int a, int b)
  {
    static std::uniform_int_distribution<> distrib(a, b);
    return distrib(rd);
  }

  double randreal(double a, double b)
  {
    static std::uniform_real_distribution<> distrib(a, b);
    return distrib(rd);
  }

\end{minted}
\end{frame}

\begin{frame}[fragile]
\label{mathutility.hpp}
\frametitle{math-utility.hpp}
\hyperlink{conclusion}{Retour au début}\newline
\begin{minted}[obeytabs=true,tabsize=2,linenos=true, breaklines]{cpp}
#ifndef MATH_UTILITY_HPP
#define MATH_UTILITY_HPP

#define EPSILON 10e-6

#include <complex>
#include <Constexpr.hpp>

namespace
{
    CONSTEXPR bool double_equal(double x, double y);
    CONSTEXPR bool complex_equal(
        std::complex<double> const &z1,
        std::complex<double> const &z2);
}

#include "math_utility.tpp"

\end{minted}
\end{frame}

\begin{frame}[fragile]
\label{mathutility.tpp}
\frametitle{math-utility.tpp}
\hyperlink{conclusion}{Retour au début}\newline
\begin{minted}[obeytabs=true,tabsize=2,linenos=true, breaklines]{cpp}
#include <cmath>
#include <complex>
#include <Constexpr.hpp>
#include "math_utility.hpp"

namespace
{
    CONSTEXPR bool double_equal(double x, double y)
    {
        return std::abs(x - y) <= EPSILON;
    }

    CONSTEXPR bool complex_equal(
        std::complex<double> const &z1,
        std::complex<double> const &z2)
    {
        return double_equal(z1.real(), z2.real()) &&
            double_equal(z1.imag(), z2.imag());
    }

\end{minted}
\end{frame}

\begin{frame}[fragile]
\label{checkpath.hpp}
\frametitle{check-path.hpp}
\hyperlink{conclusion}{Retour au début}\newline
\begin{minted}[obeytabs=true,tabsize=2,linenos=true, breaklines]{cpp}
#ifndef CHECK_PATH_HPP
#define CHECK_PATH_HPP

#include <Coord.hpp>
#include <Board.hpp>
#include <Constexpr.hpp>

template <std::size_t N, std::size_t M>
CONSTEXPR static bool
check_path_straight(
    Board<N, M> const &board,
    Coord const &dpt,
    Coord const &arv) noexcept;

template <std::size_t N, std::size_t M>
CONSTEXPR static bool
check_path_diagonal(
    Board<N, M> const &board,
    Coord const &dpt,
    Coord const &arv) noexcept;

/**
 * @brief Vérififie si il y a une pièce entre deux cases sur une instance
 * du plateau pour un mouvement orthogonale
 *
 * @tparam N Le nombre de ligne du plateau
\end{minted}
\end{frame}

\begin{frame}[fragile]
\hyperlink{conclusion}{Retour au début}\newline
\begin{minted}[obeytabs=true,tabsize=2,linenos=true,breaklines,firstnumber=27]{cpp}
 * @tparam M Le nombre de colonnes du plateau
 * @param[in] board Le plateau
 * @param[in] dpt Les coordonnées de la case de départ
 * @param[in] arv Les coordonnées de la case d'arrivé
 * @param[in] position L'instance du plateau
 * @return true si le mouvement est possible
 * @return false si le mouvement est bloqué,
 * on si ce n'est pas un mouvementorthogonal
 */
template <std::size_t N, std::size_t M>
CONSTEXPR bool
check_path_straight_1_instance(
    Board<N, M> const &board,
    Coord const &dpt,
    Coord const &arv,
    std::size_t position,
    std::optional<Coord> position_other_piece_merge);

/**
 * @brief Vérififie si il y a une pièce entre deux cases sur une
 * instance du plateau pour un mouvement diagonale.
 *
 * @tparam N Le nombre de ligne du plateau
 * @tparam M Le nombre de colonnes du plateau
 * @param[in] board Le plateau
 * @param[in] dpt Les coordonnées de la case de départ
 * @param[in] arv Les coordonnées de la case d'arrivé
 * @param[in] position L'instance du plateau
 * @return true si le mouvement est possible
 * @return false si le mouvement est bloqué, on si ce n'est pas
\end{minted}
\end{frame}

\begin{frame}[fragile]
\hyperlink{conclusion}{Retour au début}\newline
\begin{minted}[obeytabs=true,tabsize=2,linenos=true,breaklines,firstnumber=57]{cpp}
 * un mouvement diagonale.
 */
template <std::size_t N, std::size_t M>
CONSTEXPR bool check_path_diagonal_1_instance(
    Board<N, M> const &board,
    Coord const &dpt,
    Coord const &arv,
    std::size_t position,
    std::optional<Coord> position_other_piece_merge);

/**
 * @brief Vérififie si il y a une pièce entre deux cases
 * sur une instance du plateau pour un mouvement orthogonale ou diagonale.
 *
 * @tparam N Le nombre de ligne du plateau
 * @tparam M Le nombre de colonnes du plateau
 * @param[in] board Le plateau
 * @param[in] dpt Les coordonnées de la case de départ
 * @param[in] arv Les coordonnées de la case d'arrivé
 * @param[in] position L'instance du plateau
 * @return true si le mouvement est possible
 * @return false si le mouvement est bloqué, on si ce n'est pas
 * un mouvement orthogonal ou diagonale.
 */
template <std::size_t N, std::size_t M>
CONSTEXPR bool check_path_queen_1_instance(
    Board<N, M> const &board,
    Coord const &dpt,
    Coord const &arv,
    std::size_t position,
\end{minted}
\end{frame}

\begin{frame}[fragile]
\hyperlink{conclusion}{Retour au début}\newline
\begin{minted}[obeytabs=true,tabsize=2,linenos=true,breaklines,firstnumber=87]{cpp}
    std::optional<Coord> position_other_piece_merge);

#include "check_path.tpp"

\end{minted}
\end{frame}

\begin{frame}[fragile]
\label{checkpath.tpp}
\frametitle{check-path.tpp}
\hyperlink{conclusion}{Retour au début}\newline
\begin{minted}[obeytabs=true,tabsize=2,linenos=true, breaklines]{cpp}
#include <algorithm>
#include <Board.hpp>
#include <Coord.hpp>
#include <Constexpr.hpp>
#include "check_path.hpp"

template <std::size_t N, std::size_t M>
CONSTEXPR bool
check_path_straight(
    Board<N, M> const &board,
    Coord const &dpt,
    Coord const &arv) noexcept
{
    if (dpt.n == arv.n)
    {
        for (std::size_t i{std::min(dpt.m, arv.m) + 1};
             i < std::max(dpt.m, arv.m);
             i++)
        {
            if (board(dpt.n, i).get_type() != TypePiece::EMPTY ||
                board.get_proba(Coord(dpt.n, i)) < 1.)
            {
                return false;
            }
        }
        return true;
\end{minted}
\end{frame}

\begin{frame}[fragile]
\hyperlink{conclusion}{Retour au début}\newline
\begin{minted}[obeytabs=true,tabsize=2,linenos=true,breaklines,firstnumber=27]{cpp}
    }
    else if (dpt.m == arv.m)
    {
        for (std::size_t i{std::min(dpt.n, arv.n) + 1};
             i < std::max(dpt.n, arv.n);
             i++)
        {
            if (board(i, dpt.m).get_type() != TypePiece::EMPTY ||
                board.get_proba(Coord(i, dpt.m)) < 1.)
            {
                return false;
            }
        }
        return true;
    }
    return false;
}

template <std::size_t N, std::size_t M>
CONSTEXPR bool
check_path_diagonal(
    Board<N, M> const &board,
    Coord const &dpt,
    Coord const &arv) noexcept
{
    std::size_t const max_lines{std::max(dpt.n, arv.n)};
    std::size_t const min_lines{std::min(dpt.n, arv.n)};
    std::size_t const max_columns{std::max(dpt.m, arv.m)};
    std::size_t const min_columns{std::min(dpt.m, arv.m)};

\end{minted}
\end{frame}

\begin{frame}[fragile]
\hyperlink{conclusion}{Retour au début}\newline
\begin{minted}[obeytabs=true,tabsize=2,linenos=true,breaklines,firstnumber=57]{cpp}
    std::size_t const dist{max_lines - min_lines};

    if (dist == max_columns - min_columns)
    {
        for (std::size_t i{1}; i < dist; i++)
        {
            if (board(min_lines + i, min_columns + i).get_type() != TypePiece::EMPTY)
            {
                return false;
            }
        }
        return true;
    }
    return false;
}

template <std::size_t N, std::size_t M>
CONSTEXPR bool
check_path_straight_1_instance(
    Board<N, M> const &board,
    Coord const &dpt,
    Coord const &arv,
    std::size_t position,
    std::optional<Coord> position_other_piece_merge)
{
    if (dpt.n == arv.n)
    {
        for (std::size_t i{std::min(dpt.m, arv.m) + 1};
             i < std::max(dpt.m, arv.m);
             i++)
\end{minted}
\end{frame}

\begin{frame}[fragile]
\hyperlink{conclusion}{Retour au début}\newline
\begin{minted}[obeytabs=true,tabsize=2,linenos=true,breaklines,firstnumber=87]{cpp}
        {
            if (board.m_board[position].first[board.offset(dpt.n, i)])
            {
                if (position_other_piece_merge.has_value())
                {
                    if (!(board.offset(position_other_piece_merge.value().n,
                                       position_other_piece_merge.value().m) == board.offset(dpt.n, i)))
                    {
                        return false;
                    }
                }
                else
                {
                    return false;
                }
            }
        }
            return true;
        
    }
        else if (dpt.m == arv.m)
        {
            for (std::size_t i{std::min(dpt.n, arv.n) + 1};
                 i < std::max(dpt.n, arv.n);
                 i++)
            {
                if (board.m_board[position].first[board.offset(i, dpt.m)])
                {
                    if (position_other_piece_merge.has_value())
                    {
\end{minted}
\end{frame}

\begin{frame}[fragile]
\hyperlink{conclusion}{Retour au début}\newline
\begin{minted}[obeytabs=true,tabsize=2,linenos=true,breaklines,firstnumber=117]{cpp}
                        if (!(board.offset(position_other_piece_merge.value().n,
                                           position_other_piece_merge.value().m) == board.offset(i, dpt.m)))
                        {
                            return false;
                        }
                    }
                    else
                    {
                        return false;
                    }
                }
            }
            return true;
        }
        return false;
    }

    template <std::size_t N, std::size_t M>
    CONSTEXPR bool
    check_path_diagonal_1_instance(
        Board<N, M> const &board,
        Coord const &dpt,
        Coord const &arv,
        std::size_t position,
        std::optional<Coord> position_other_piece_merge)
    {
        std::size_t const max_lines{std::max(dpt.n, arv.n)};
        std::size_t const min_lines{std::min(dpt.n, arv.n)};
        std::size_t const max_columns{std::max(dpt.m, arv.m)};
        std::size_t const min_columns{std::min(dpt.m, arv.m)};
\end{minted}
\end{frame}

\begin{frame}[fragile]
\hyperlink{conclusion}{Retour au début}\newline
\begin{minted}[obeytabs=true,tabsize=2,linenos=true,breaklines,firstnumber=147]{cpp}

        std::size_t const dist{max_lines - min_lines};

        if (dist == max_columns - min_columns)
        {
            for (std::size_t i{1}; i < dist; i++)
            {
                if (board.m_board[position]
                        .first[board.offset(min_lines + i, min_columns + i)])
                {
                    if (position_other_piece_merge.has_value())
                    {
                        if (!(board.offset(position_other_piece_merge.value().n,
                                           position_other_piece_merge.value().m) ==
                              board.offset(min_lines + i,
                                           min_columns + i)))
                        {
                            return false;
                        }
                    }
                    else
                    {
                        return false;
                    }
                }
            }
            return true;
        }
        return false;
    }
\end{minted}
\end{frame}

\begin{frame}[fragile]
\hyperlink{conclusion}{Retour au début}\newline
\begin{minted}[obeytabs=true,tabsize=2,linenos=true,breaklines,firstnumber=177]{cpp}

    template <std::size_t N, std::size_t M>
    CONSTEXPR bool check_path_queen_1_instance(
        Board<N, M> const &board,
        Coord const &dpt,
        Coord const &arv,
        std::size_t position,
        std::optional<Coord> position_other_piece_merge)
    {
        return check_path_straight_1_instance<N, M>(board, dpt, arv, position, position_other_piece_merge) ||
               check_path_diagonal_1_instance<N, M>(board, dpt, arv, position, position_other_piece_merge);

\end{minted}
\end{frame}

\begin{frame}[fragile]
\label{Unitary.hpp}
\frametitle{Unitary.hpp}
\hyperlink{conclusion}{Retour au début}\newline
\begin{minted}[obeytabs=true,tabsize=2,linenos=true, breaklines]{cpp}
#ifndef UNITARY_HPP
#define UNITARY_HPP
#include <CMatrix.hpp>
#include <numbers>
#include <complex>
#include <Constexpr.hpp>

using namespace std::complex_literals;
using namespace std::numbers;

constexpr inline
CMatrix<4> MATRIX_ISWAP {
    {
        {1,  0,  0, 0},
        {0,  0, 1i, 0},
        {0, 1i,  0, 0},
        {0,  0,  0, 1}
    }
};


constexpr inline
CMatrix<4> MATRIX_SQRT_ISWAP {
    {
        {1,        0,         0, 0},
        {0, 1./sqrt2,  1i/sqrt2, 0},
\end{minted}
\end{frame}

\begin{frame}[fragile]
\hyperlink{conclusion}{Retour au début}\newline
\begin{minted}[obeytabs=true,tabsize=2,linenos=true,breaklines,firstnumber=27]{cpp}
        {0, 1i/sqrt2,  1./sqrt2, 0},
        {0,        0,         0, 1}
    }
};

constexpr inline
CMatrix<4> MATRIX_JUMP { MATRIX_ISWAP };

constexpr inline
CMatrix<8> MATRIX_SLIDE {
    {
        {1, 0, 0, 0, 0, 0, 0, 0},
        {0, 0, 1i, 0, 0, 0, 0, 0 },
        {0, 1i, 0, 0, 0, 0, 0, 0},
        {0, 0, 0, 1, 0, 0, 0, 0},
        {0, 0, 0, 0, 1, 0, 0, 0 },
        {0, 0, 0, 0, 0, 1, 0, 0},
        {0, 0, 0, 0, 0, 0, 1, 0},
        {0, 0, 0, 0, 0, 0, 0, 1}
    }
};

constexpr inline
CMatrix<8> MATRIX_ISWAP_8 {
    {
        {1, 0, 0, 0, 0, 0, 0, 0},
        {0, 0, 0, 0, 1i, 0, 0, 0 },
        {0, 0, 1, 0, 0, 0, 0, 0},
        {0, 0, 0, 0, 0, 0, 1i, 0},
        {0, 1i, 0, 0, 0, 0, 0, 0 },// erreur dans le papier?
\end{minted}
\end{frame}

\begin{frame}[fragile]
\hyperlink{conclusion}{Retour au début}\newline
\begin{minted}[obeytabs=true,tabsize=2,linenos=true,breaklines,firstnumber=57]{cpp}
        {0, 0, 0, 0, 0, 1, 0, 0},
        {0, 0, 0, 1i, 0, 0, 0, 0},
        {0, 0, 0, 0, 0, 0, 0, 1}
    }
};

constexpr inline
CMatrix<8> MATRIX_SQRT_ISWAP_8 {
    {
        {1,        0,        0, 0,  0,        0,        0, 0},
        {0, 1./sqrt2, 1i/sqrt2, 0, 0,        0,        0, 0},
        {0, 1i/sqrt2, 1./sqrt2, 0,  0,        0,        0, 0},
        {0,        0,        0, 1,  0,        0,        0, 0},
        {0,        0,        0, 0,  1,        0,        0, 0},
        {0,        0,        0, 0,  0, 1./sqrt2, 1i/sqrt2, 0},
        {0,        0,        0, 0,  0, 1i/sqrt2, 1./sqrt2, 0},
        {0,        0,        0, 0,  0,        0,        0, 1}
    }
};

constexpr inline
CMatrix<8> MATRIX_SPLIT {
    {
        {1,        0,         0,   0,  0,         0,         0, 0},
        {0,        0,         0,   0, 1i,         0,         0, 0},
        {0, 1i/sqrt2,  1./sqrt2,   0,  0,         0,         0, 0},
        {0,        0,         0,   0,  0, -1./sqrt2,  1i/sqrt2, 0},
        {0, 1i/sqrt2, -1./sqrt2,   0,  0,         0,         0, 0},
        {0,        0,         0,   0,  0,  1i/sqrt2, -1/sqrt2, 0},
        {0,        0,         0,   1i,  0,         0,         0, 0},
\end{minted}
\end{frame}

\begin{frame}[fragile]
\hyperlink{conclusion}{Retour au début}\newline
\begin{minted}[obeytabs=true,tabsize=2,linenos=true,breaklines,firstnumber=87]{cpp}
        {0,        0,         0,   0,  0,         0,         0, 1}
    }
};

constexpr inline
CMatrix<8> MATRIX_MERGE {
    {
        {1,   0,         0,         0,         0,         0,   0,  0},
        {0,   0, -1i/sqrt2,         0, -1i/sqrt2,         0,   0,  0},
        {0,   0,   1/sqrt2,         0,  -1/sqrt2,         0,   0,  0},
        {0,   0,         0,         0,         0,         0, -1i,  0},
        {0, -1i,         0,         0,         0,         0,    0, 0},
        {0,   0,         0,  -1/sqrt2,         0, -1i/sqrt2,    0, 0},
        {0,   0,         0, -1i/sqrt2,         0,  -1/sqrt2,    0, 0},
        {0,   0,         0,         0,         0,         0,    0, 1}
    }
};

constexpr inline
CMatrix<32> MATRIX_SPLIT_SLIDE {
    CMatrix<16>{
        MATRIX_SPLIT, CMatrix<8>{},
        CMatrix<8>{}, MATRIX_ISWAP_8
    },                                  CMatrix<16>{},
    CMatrix<16>{},                      CMatrix<16>{
                                            CMatrix<2>::identity()
                                                .tensoriel_product(MATRIX_ISWAP), CMatrix<8>{},
                                            CMatrix<8>{}, CMatrix<8>::identity()
                                        }
};
\end{minted}
\end{frame}

\begin{frame}[fragile]
\hyperlink{conclusion}{Retour au début}\newline
\begin{minted}[obeytabs=true,tabsize=2,linenos=true,breaklines,firstnumber=117]{cpp}

CONSTEXPR inline
CMatrix<32> MATRIX_MERGE_SLIDE {
    CMatrix<16>{
        MATRIX_MERGE, CMatrix<8>{},
        CMatrix<8>{}, conj(CMatrix<2>::identity().tensoriel_product(MATRIX_ISWAP).transposed())
    },                                                  
    CMatrix<16>{},
    CMatrix<16>{},                                      
    CMatrix<16>{
        conj(MATRIX_ISWAP_8.transposed()), CMatrix<8>{},
        CMatrix<8>{},                CMatrix<8>::identity()
    }
};




\end{minted}
\end{frame}

\begin{frame}[fragile]
\label{CMatrix.hpp}
\frametitle{CMatrix.hpp}
\hyperlink{conclusion}{Retour au début}\newline
\begin{minted}[obeytabs=true,tabsize=2,linenos=true, breaklines]{cpp}
#ifndef CMATRIX_HPP
#define CMATRIX_HPP
#include <Matrix.hpp>
#include <complex>
#include <Constexpr.hpp>
#include <cstddef>

/**
 * @brief Objet représentant les matrices carrées complexe
 * de dimension N
 * 
 * @tparam N La dimension de la matrice
 */
template <std::size_t N>
using CMatrix = Matrix<std::complex<double>, N>;

namespace
{
    template <std::size_t N>
    CONSTEXPR CMatrix<N> conj(CMatrix<N> const & matrix)
    {
        CMatrix<N> matrix_conj{};
        for(std::size_t i = 0; i<N; i++)
        {
            for(std::size_t j = 0; j<N; j++)
            {
\end{minted}
\end{frame}

\begin{frame}[fragile]
\hyperlink{conclusion}{Retour au début}\newline
\begin{minted}[obeytabs=true,tabsize=2,linenos=true,breaklines,firstnumber=27]{cpp}
                matrix_conj(i, j) = std::conj(matrix(i, j));
            }
        }
        return matrix_conj;
    }
}

/**
 * @brief Réalise l'opération 2^n
 * @warning Est valable uniquement sur les types entiers
 */
#define _2POW(n) (1 << (n))


\end{minted}
\end{frame}

\begin{frame}[fragile]
\label{Complexprinter.hpp}
\frametitle{Complex-printer.hpp}
\hyperlink{conclusion}{Retour au début}\newline
\begin{minted}[obeytabs=true,tabsize=2,linenos=true, breaklines]{cpp}
#ifndef COMPLEX_PRINTER_HPP
#define COMPLEX_PRINTER_HPP

#include <complex>
#include <concepts>
#include <string>

namespace std
{
    /**
     * @brief convertit un nombre complexe en chaine de caractère
     */
    using namespace std::literals;
    template <std::floating_point T>
    std::string to_string(std::complex<T> value)
    {
        return std::string{std::to_string(std::real(value)) 
                            + ((std::imag(value) >= 0) ? '+' : '\0') 
                            + std::to_string(std::imag(value)) + 'i'};
    }

    /**
     * @brief Implémente l'opérateur de flux sortant
     * sur les nombres complexes
     * 
     * @tparam T le type de représentation floatante du complexe
\end{minted}
\end{frame}

\begin{frame}[fragile]
\hyperlink{conclusion}{Retour au début}\newline
\begin{minted}[obeytabs=true,tabsize=2,linenos=true,breaklines,firstnumber=27]{cpp}
     * @param os Le flux de sorti
     * @param value Le nombre complexe a transférer
     * @return std::ostream& retourne le flux passé en paramètre
     */
    template <std::floating_point T>
    std::ostream& operator<<(std::ostream& os, std::complex<T> value)
    {
        os << std::real(value) << ((std::imag(value) >= 0) ? '+' : '\0' ) << std::imag(value) << 'i';
        return os;
    }
}


\end{minted}
\end{frame}

\begin{frame}[fragile]
\label{Constexpr.hpp}
\frametitle{Constexpr.hpp}
\hyperlink{conclusion}{Retour au début}\newline
\begin{minted}[obeytabs=true,tabsize=2,linenos=true, breaklines]{cpp}
#ifndef CONSTEXPR_HPP
#define CONSTEXPR_HPP

/**
 * @brief Utilisé pour utiliser ou non constexpr
 */
#define CONSTEXPR constexpr


\end{minted}
\end{frame}

\begin{frame}[fragile]
\label{ConsoleInterface.hpp}
\frametitle{ConsoleInterface.hpp}
\hyperlink{conclusion}{Retour au début}\newline
\begin{minted}[obeytabs=true,tabsize=2,linenos=true, breaklines]{cpp}
#ifndef CONSOLE_INTERFACE_HPP
#define CONSOLE_INTERFACE_HPP

#include <ostream>
#include <Board.hpp>

template <std::size_t N, std::size_t M>
std::ostream& operator<<(std::ostream& os, Board<N, M> const& board);

#include "ConsoleInterface.tpp"

\end{minted}
\end{frame}

\begin{frame}[fragile]
\label{ConsoleInterface.tpp}
\frametitle{ConsoleInterface.tpp}
\hyperlink{conclusion}{Retour au début}\newline
\begin{minted}[obeytabs=true,tabsize=2,linenos=true, breaklines]{cpp}
#include <iostream>
#include <Board.hpp>

const char *getUnicodeChar(TypePiece piece, Color color)
{
  using enum TypePiece;
  if (color == Color::WHITE)
  {
    switch (piece)
    {
    case KING:
      return "K";
    case QUEEN:
      return "Q";
    case ROOK:
      return "R";
    case BISHOP:
      return "B";
    case KNIGHT:
      return "N";
    case PAWN:
      return "P";
    case EMPTY:
    default:
      return " ";
    }
\end{minted}
\end{frame}

\begin{frame}[fragile]
\hyperlink{conclusion}{Retour au début}\newline
\begin{minted}[obeytabs=true,tabsize=2,linenos=true,breaklines,firstnumber=27]{cpp}
  }
  else
  {
    switch (piece)
    {
    case KING:
      return "k";
    case QUEEN:
      return "q";
    case ROOK:
      return "r";
    case BISHOP:
      return "b";
    case KNIGHT:
      return "n";
    case PAWN:
      return "p";
    case EMPTY:
    default:
      return " ";
    }
  }
}

template <std::size_t N, std::size_t M>
std::ostream &operator<<(std::ostream &os, Board<N, M> const &board)
{
  for (std::size_t i{0}; i < M; i++)
  {
    os << "|-----";
\end{minted}
\end{frame}

\begin{frame}[fragile]
\hyperlink{conclusion}{Retour au début}\newline
\begin{minted}[obeytabs=true,tabsize=2,linenos=true,breaklines,firstnumber=57]{cpp}
  }
  os << "|\n";
  for (std::size_t i{0}; i < N; i++)
  {
    for (std::size_t j{0}; j < M; j++)
    {
      TypePiece piece{board(i, j).get_type()};
      Color color{board(i, j).get_color()};
      os << "|  " << getUnicodeChar(piece, color) << "  ";
    }
    os << "|\n";
    for (std::size_t j{0}; j < M; j++)
    {
      if (board(i, j).get_type() != TypePiece::EMPTY)
      {
        double proba{board.get_proba(Coord(i, j))};
        int p{static_cast<int>(std::round(100. * proba))};
        if (p != 100)
        {
          os << "| ";
          if (p < 10){
            os << '0';
          }
          os << p << "% ";
        }
        else
        {
          os << "|     ";
        }
      }
\end{minted}
\end{frame}

\begin{frame}[fragile]
\hyperlink{conclusion}{Retour au début}\newline
\begin{minted}[obeytabs=true,tabsize=2,linenos=true,breaklines,firstnumber=87]{cpp}
      else
      {
        os << "|     ";
      }
    }
    os << "|\n";
    for (std::size_t j{0}; j < M; j++)
    {
      os << "|-----";
    }
    os << "|\n";
  }
  (void)board;
  return os;

\end{minted}
\end{frame}

\begin{frame}[fragile]
\label{finalsolver.hpp}
\frametitle{final-solver.hpp}
\hyperlink{conclusion}{Retour au début}\newline
\begin{minted}[obeytabs=true,tabsize=2,linenos=true, breaklines]{cpp}
#ifndef FINAL_SOLVER_HPP
#define FINAL_SOLVER_HPP
#include <Color.hpp>
#include <Constexpr.hpp>
#include <memory>
#include <stack>

struct Node
{
    Node(Move const &y, std::stack<Node> &&o = std::stack<Node>{});
    Move m;
    std::stack<Node> other_way;
};

Node::Node(Move const &y, std::stack<Node> &&o) : m(y), other_way(o)
{
}

namespace Final
{
    template <std::size_t N, std::size_t M>
    CONSTEXPR bool brut_force_classic_chess(Board<N, M> const &board, std::size_t profondeur, Color c);

    template <std::size_t N, std::size_t M>
    CONSTEXPR bool brut_force_quantum_chess(Board<N, M> const &board, std::size_t profondeur, Color c);

\end{minted}
\end{frame}

\begin{frame}[fragile]
\hyperlink{conclusion}{Retour au début}\newline
\begin{minted}[obeytabs=true,tabsize=2,linenos=true,breaklines,firstnumber=27]{cpp}
    template <std::size_t N, std::size_t M>
    CONSTEXPR std::pair<double, std::stack<Node>> res_pos(Board<N, M> &board, std::size_t profondeur);
    template<std::size_t N>
    CONSTEXPR void print_stack(std::stack<Node> &stack);
};

/*struct C_hash{
std::size_t operator () (std::pair<TypePiece, Coord> const &c) const;
};
bool operator==(std::pair<TypePiece, Coord> const &lhs, std::pair<TypePiece, Coord> const &rhs);
*/

#include "final_solver.tpp"


\end{minted}
\end{frame}

\begin{frame}[fragile]
\label{finalsolver.tpp}
\frametitle{final-solver.tpp}
\hyperlink{conclusion}{Retour au début}\newline
\begin{minted}[obeytabs=true,tabsize=2,linenos=true, breaklines]{cpp}
#include <Board.hpp>
#include <Piece.hpp>
#include <Move.hpp>
#include <iostream>
#include <Color.hpp>
#include <unordered_map>
#include <functional>
#include <stack>

/*std::size_t C_hash::operator()(std::pair<TypePiece, Coord> const &c) const
{
    std::size_t h1 = std::hash<std::size_t>()(c.second.n);
    std::size_t h2 = std::hash<std::size_t>()(c.second.m);

    return h1 ^ h2;
}

bool operator==(std::pair<TypePiece, Coord> const &lhs, std::pair<TypePiece, Coord> const &rhs)
{
    return lhs.second.n == rhs.second.n && lhs.second.m == rhs.second.m && lhs.first == rhs.first;
}*/
template <std::size_t N, std::size_t M>
CONSTEXPR bool Final::brut_force_classic_chess(Board<N, M> &board, std::size_t profondeur, Color c)
{
    if (board.winning_position(c))
    {
\end{minted}
\end{frame}

\begin{frame}[fragile]
\hyperlink{conclusion}{Retour au début}\newline
\begin{minted}[obeytabs=true,tabsize=2,linenos=true,breaklines,firstnumber=27]{cpp}
        return true;
    }
    else
    {
        Color col = Color::WHITE;
        if (c == Color::WHITE)
        {
            col = Color::BLACK;
        }
        if (board.winning_position(col))
        {
            return false;
        }
        else
        {
            if (profondeur == 0)
            {
                // std::cout<<"profondeur pas assez élevé"<<std::endl;
                return false;
            }
            else
            {

                if (board.get_current_player() == c)
                {
                    for (std::size_t i{0}; i < N; i++)
                    {
                        for (std::size_t j{0}; j < M; j++)
                        {
                            if (board(i, j).get_type() != TypePiece::EMPTY && board(i, j).get_color() == board.get_current_player())
\end{minted}
\end{frame}

\begin{frame}[fragile]
\hyperlink{conclusion}{Retour au début}\newline
\begin{minted}[obeytabs=true,tabsize=2,linenos=true,breaklines,firstnumber=57]{cpp}
                            {
                                std::forward_list<Coord> l = board.get_list_normal_move(Coord(i, j));
                                for (auto const &e : l)
                                {
                                    Move m = Move_classic(Coord(i, j), e);
                                    Board<N, M> b = board;
                                    b.change_player();
                                    b.move(m);
                                    if (brut_force_classic_chess(b, profondeur - 1, c))
                                    {
                                        return true;
                                    }
                                }
                            }
                        }
                    }
                    return false;
                }
                else
                {
                    for (std::size_t i{0}; i < N; i++)
                    {
                        for (std::size_t j{0}; j < M; j++)
                        {
                            if (board(i, j).get_type() != TypePiece::EMPTY && board(i, j).get_color() == board.get_current_player())
                            {
                                std::forward_list<Coord> l = board.get_list_normal_move(Coord(i, j));
                                for (auto const &e : l)
                                {
                                    Move m = Move_classic(Coord(i, j), e);
\end{minted}
\end{frame}

\begin{frame}[fragile]
\hyperlink{conclusion}{Retour au début}\newline
\begin{minted}[obeytabs=true,tabsize=2,linenos=true,breaklines,firstnumber=87]{cpp}
                                    Board<N, M> b = board;
                                    b.change_player();
                                    b.move(m);
                                    if (!brut_force_classic_chess(b, profondeur - 1, c))
                                    {
                                        return false;
                                    }
                                }
                            }
                        }
                    }
                    return true;
                }
            }
        }
    }
    std::cout << "erreur" << std::endl;
    return false;
}

template <std::size_t N, std::size_t M>
CONSTEXPR bool Final::brut_force_quantum_chess(Board<N, M> &board, std::size_t profondeur, Color c)
{
    bool found_win_pos{false};
    if (board.winning_position(c))
    {
        return true;
    }
    else
    {
\end{minted}
\end{frame}

\begin{frame}[fragile]
\hyperlink{conclusion}{Retour au début}\newline
\begin{minted}[obeytabs=true,tabsize=2,linenos=true,breaklines,firstnumber=117]{cpp}
        if (profondeur == 0)
        {
            return false;
        }
        else
        {
            board.all_move(
                [&board,
                 profondeur,
                 c,
                 &found_win_pos](Move const &m) mutable -> bool
                {
                    Board<N, M> board_cpy = board;
                    board_cpy.move(m);
                    board_cpy.change_player();
                    if (c == board.get_current_player())
                    {
                        if (brut_force_quantum_chess(board_cpy, profondeur - 1, c))
                        {
                            found_win_pos = true;
                            return true;
                        }
                    }
                    else
                    {
                        if (!brut_force_quantum_chess(board_cpy, profondeur - 1, c))
                        {
                            return true;
                        }
                        found_win_pos = true;
\end{minted}
\end{frame}

\begin{frame}[fragile]
\hyperlink{conclusion}{Retour au début}\newline
\begin{minted}[obeytabs=true,tabsize=2,linenos=true,breaklines,firstnumber=147]{cpp}
                    }
                    return false;
                });
        }
    }
    return found_win_pos;
}

template <std::size_t N, std::size_t M>
double evaluer(const Board<N, M> &board, Color c)
{

    if (board.winning_position(c))
    {
        return 1.;
    }
    else
    {
        if (board.winning_position(opponent_color(c)))
        {
            return -1.;
        }
        else
        {
            return 0.;
        }
    }
}
template<std::size_t N>
void print_move(std::ostream &output, Move m)
\end{minted}
\end{frame}

\begin{frame}[fragile]
\hyperlink{conclusion}{Retour au début}\newline
\begin{minted}[obeytabs=true,tabsize=2,linenos=true,breaklines,firstnumber=177]{cpp}
{
    std::string data1{{static_cast<char>('a' + m.normal.src.m), static_cast<char>('0' + N - m.normal.src.n)}};
    std::string data2{{static_cast<char>('a' + m.normal.arv.m), static_cast<char>('0' + N - m.normal.arv.n)}};
    if (m.type == TypeMove::NORMAL)
    {
        output << "N " << data1 << data2 ;
    }
    else if (m.type == TypeMove::PROMOTE)
    {
        output << "P " << data1 << data2 << " :: " << std::to_string(static_cast<int>(m.promote.piece)) ;
    }
    else
    {
        std::string data3{{static_cast<char>('a' + m.split.arv2.m), static_cast<char>('0' + N - m.split.arv2.n)}};
        if (m.type == TypeMove::SPLIT)
        {
            output << "S " << data1 << '(' << data2 << data3 << ')' ;
        }
        else
        {
            output << "M " << '(' << data1 << data2 << ')' << data3 ;
        }
    }
}
// Fonction récursive pour l'élagage alpha-bêta
template <std::size_t N, std::size_t M>
CONSTEXPR double alphaBeta(Board<N, M> const &board, std::size_t profondeur, double alpha, double beta, bool estMax, Color c, std::stack<Node> &stack)
{
    if (profondeur == 0 || board.winning_position(c) || board.winning_position(opponent_color(c)))
    {
\end{minted}
\end{frame}

\begin{frame}[fragile]
\hyperlink{conclusion}{Retour au début}\newline
\begin{minted}[obeytabs=true,tabsize=2,linenos=true,breaklines,firstnumber=207]{cpp}
        Move m = Move_classic(Coord(0, 0), Coord(0, 0));
        stack.push(m); // Cette information n'a pas d'intérêt dans la pile mais évite les erreurs de segmentation lorsqu'on dépile
        return evaluer(board, c);
    }
    else
    {
        Move meilleur_move;
        if (estMax)
        {
            double meilleurScore = {-2.};
            board.all_move(
                [&board,
                 profondeur,
                 c,
                 &beta,
                 &alpha,
                 &meilleur_move,
                 &meilleurScore,
                 &stack](Move const &m) mutable -> bool

                {
                    double res;
                    bool one_move{true};
                    double proba_move{board.get_proba_move(m)};
                    if (double_equal(proba_move, 1.))
                    {
                        Board<N, M> board_cpy = board;
                        board_cpy.change_player();
                        board_cpy.move(m, true);
                        double eval{evaluer(board_cpy, c)};
\end{minted}
\end{frame}

\begin{frame}[fragile]
\hyperlink{conclusion}{Retour au début}\newline
\begin{minted}[obeytabs=true,tabsize=2,linenos=true,breaklines,firstnumber=237]{cpp}
                        if (double_equal(eval, 1.))
                        {
                            meilleur_move = m;
                            meilleurScore = 1.;
                            return true;
                        }
                        else
                        {
                            if (double_equal(eval, -1.))
                            {
                                if (-1. > meilleurScore)
                                {
                                    meilleurScore = -1.;
                                    meilleur_move = m;
                                }
                            }
                            else
                            {
                                std::stack<Node> new_stack{};
                                res = alphaBeta(board_cpy, profondeur - 1, alpha, beta, false, c, new_stack);
                                if (res > meilleurScore)
                                {
                                    meilleurScore = res;
                                    meilleur_move = m;
                                    stack = std::move(new_stack);
                                }
                            }
                        }
                    }
                    else
\end{minted}
\end{frame}

\begin{frame}[fragile]
\hyperlink{conclusion}{Retour au début}\newline
\begin{minted}[obeytabs=true,tabsize=2,linenos=true,breaklines,firstnumber=267]{cpp}
                    {
                        if (double_equal(proba_move, 0.))
                        {
                            Board<N, M> board_cpy = board;
                            board_cpy.change_player();
                            board_cpy.move(m, false);
                            double eval{evaluer(board_cpy, c)};
                            if (double_equal(eval, 1.))
                            {
                                meilleur_move = m;
                                meilleurScore = 1.;
                                return true;
                            }
                            else
                            {
                                if (double_equal(eval, -1.))
                                {
                                    if (-1. > meilleurScore)
                                    {
                                        meilleurScore = -1.;
                                        meilleur_move = m;
                                    }
                                }
                                else
                                {
                                    std::stack<Node> new_stack{};
                                    res = alphaBeta(board_cpy, profondeur - 1, alpha, beta, false, c, new_stack);
                                    if (res > meilleurScore)
                                    {
                                        meilleurScore = res;
\end{minted}
\end{frame}

\begin{frame}[fragile]
\hyperlink{conclusion}{Retour au début}\newline
\begin{minted}[obeytabs=true,tabsize=2,linenos=true,breaklines,firstnumber=297]{cpp}
                                        meilleur_move = m;
                                        stack = std::move(new_stack);
                                    }
                                }
                            }
                        }
                        else
                        {
                            one_move = false;
                            Board<N, M> board_cpy1 = board;
                            board_cpy1.change_player();
                            board_cpy1.move(m, true);
                            Board<N, M> board_cpy2 = board;
                            board_cpy2.change_player();
                            board_cpy2.move(m, false);
                            double eval1{evaluer(board_cpy1, c)};
                            std::stack<Node> new_stack1{};
                            std::stack<Node> new_stack2{};
                            if (double_equal(eval1, 0.))
                            {
                                eval1 = alphaBeta(board_cpy1, profondeur - 1, -2., 2., false, c, new_stack1);
                            }
                            double eval2{evaluer(board_cpy2, c)};
                            if (double_equal(eval2, 0.))
                            {
                                eval2 = alphaBeta(board_cpy2, profondeur - 1, -2., 2., false, c, new_stack2);
                            }
                            res = eval1 * proba_move + eval2 * (1 - proba_move);
                            if (res > meilleurScore)
                            {
\end{minted}
\end{frame}

\begin{frame}[fragile]
\hyperlink{conclusion}{Retour au début}\newline
\begin{minted}[obeytabs=true,tabsize=2,linenos=true,breaklines,firstnumber=327]{cpp}
                                meilleurScore = res;
                                meilleur_move = m;
                                stack = std::move(new_stack1);
                                if (stack.empty())
                                {
                                    stack = std::move(new_stack2);
                                }
                                else
                                {
                                    Node inter = stack.top();
                                    stack.pop();
                                    inter.other_way = std::move(new_stack2);
                                }
                            }
                        }
                    }
                    if (one_move)
                    {
                        alpha = std::max(alpha, meilleurScore);
                        if (beta <= alpha)
                        {
                            return true; // Élagage bêta
                        }
                    }
                    return false;
                });
            stack.push(Node(meilleur_move));
            return meilleurScore;
        }
        else
\end{minted}
\end{frame}

\begin{frame}[fragile]
\hyperlink{conclusion}{Retour au début}\newline
\begin{minted}[obeytabs=true,tabsize=2,linenos=true,breaklines,firstnumber=357]{cpp}
        {
            double meilleurScore{2.};
            board.all_move(
                [&board,
                 profondeur,
                 c,
                 &alpha,
                 &meilleurScore,
                 &meilleur_move,
                 &beta,
                 &stack](Move const &m) mutable -> bool

                {
                    double res;
                    bool one_move{true};
                    double proba_move{board.get_proba_move(m)};
                    if (double_equal(proba_move, 1.))
                    {
                        Board<N, M> board_cpy = board;
                        board_cpy.change_player();
                        board_cpy.move(m, true);
                        double eval{evaluer(board_cpy, c)};
                        if (double_equal(eval, -1.))
                        {
                            meilleur_move = m;
                            meilleurScore = -1.;
                            return true;
                        }
                        else
                        {
\end{minted}
\end{frame}

\begin{frame}[fragile]
\hyperlink{conclusion}{Retour au début}\newline
\begin{minted}[obeytabs=true,tabsize=2,linenos=true,breaklines,firstnumber=387]{cpp}
                            if (double_equal(eval, 1.))
                            {
                                if (1. < meilleurScore)
                                {
                                    meilleurScore = 1.;
                                    meilleur_move = m;
                                }
                            }
                            else
                            {
                                std::stack<Node> new_stack{};
                                res = alphaBeta(board_cpy, profondeur - 1, alpha, beta, true, c, new_stack);
                                if (res < meilleurScore)
                                {
                                    meilleurScore = res;
                                    meilleur_move = m;
                                    stack = std::move(new_stack);
                                }
                            }
                        }
                    }
                    else
                    {
                        if (double_equal(proba_move, 0.))
                        {
                            Board<N, M> board_cpy = board;
                            board_cpy.change_player();
                            board_cpy.move(m, false);
                            double eval{evaluer(board_cpy, c)};
                            if (double_equal(eval, -1.))
\end{minted}
\end{frame}

\begin{frame}[fragile]
\hyperlink{conclusion}{Retour au début}\newline
\begin{minted}[obeytabs=true,tabsize=2,linenos=true,breaklines,firstnumber=417]{cpp}
                            {
                                meilleur_move = m;
                                meilleurScore = -1.;
                                return true;
                            }
                            else
                            {
                                if (double_equal(eval, 1.))
                                {
                                    if (1. < meilleurScore)
                                    {
                                        meilleurScore = 1.;
                                        meilleur_move = m;
                                    }
                                }
                                else
                                {
                                    std::stack<Node> new_stack{};
                                    res = alphaBeta(board_cpy, profondeur - 1, alpha, beta, true, c, new_stack);
                                    if (res < meilleurScore)
                                    {
                                        meilleurScore = res;
                                        meilleur_move = m;
                                    }
                                }
                            }
                        }
                        else
                        {
                            one_move = false;
\end{minted}
\end{frame}

\begin{frame}[fragile]
\hyperlink{conclusion}{Retour au début}\newline
\begin{minted}[obeytabs=true,tabsize=2,linenos=true,breaklines,firstnumber=447]{cpp}
                            Board<N, M> board_cpy1 = board;
                            board_cpy1.change_player();
                            board_cpy1.move(m, true);
                            Board<N, M> board_cpy2 = board;
                            board_cpy2.change_player();
                            board_cpy2.move(m, false);
                            double eval1{evaluer(board_cpy1, c)};
                            std::stack<Node> new_stack1{};
                            std::stack<Node> new_stack2{};
                            if (double_equal(eval1, 0.))
                            {
                                eval1 = alphaBeta(board_cpy1, profondeur - 1, -2., 2., true, c, new_stack1);
                            }
                            double eval2{evaluer(board_cpy2, c)};
                            if (double_equal(eval2, 0.))
                            {
                                eval2 = alphaBeta(board_cpy2, profondeur - 1, -2., 2., true, c, new_stack2);
                            }
                            res = eval1 * proba_move + eval2 * (1 - proba_move);
                            if (res < meilleurScore)
                            {
                                meilleurScore = res;
                                meilleur_move = m;

                                stack = std::move(new_stack1);
                                if (stack.empty())
                                {
                                    stack = std::move(new_stack2);
                                }
                                else
\end{minted}
\end{frame}

\begin{frame}[fragile]
\hyperlink{conclusion}{Retour au début}\newline
\begin{minted}[obeytabs=true,tabsize=2,linenos=true,breaklines,firstnumber=477]{cpp}
                                {
                                    Node inter = stack.top();
                                    stack.pop();
                                    inter.other_way = std::move(new_stack2);
                                    stack.push(inter);
                                }
                            }
                        }
                    }
                    if (one_move)
                    {
                        beta = std::min(beta, meilleurScore);
                        if (beta <= alpha)
                        {
                            return true; // Élagage bêta
                        }
                    }
                    return false;
                });
            stack.push(Node(meilleur_move));
            return meilleurScore;
        }
    }
}

template <std::size_t N, std::size_t M>
CONSTEXPR std::pair<double, std::stack<Node>> Final::res_pos(Board<N, M> &board, std::size_t profondeur)
{
    Color c{board.get_current_player()};
    std::stack<Node> stack{};
\end{minted}
\end{frame}

\begin{frame}[fragile]
\hyperlink{conclusion}{Retour au début}\newline
\begin{minted}[obeytabs=true,tabsize=2,linenos=true,breaklines,firstnumber=507]{cpp}
    return std::make_pair(alphaBeta(board, profondeur, -2., 2., true, c, stack), stack);
}
template <std::size_t N>
CONSTEXPR void Final::print_stack( std::stack<Node> & stack)
{
    std::vector<std::stack<Node>> tab{};
    tab.push_back(stack);
    std::size_t compteur{0};
    while (compteur < std::size(tab))
    {
        compteur = 0;
        for (auto &e : tab)
        {
            if (e.empty())
            {
                compteur++;
                std::cout << "      ";
            }
            else
            {
                Node elt = e.top();
                e.pop();
                if(!elt.other_way.empty())
                {
                    tab.push_back(elt.other_way);
                }
                print_move<N>(std::cout, elt.m);
            }
        }
        std::cout<<std::endl;
\end{minted}
\end{frame}

\begin{frame}[fragile]
\hyperlink{conclusion}{Retour au début}\newline
\begin{minted}[obeytabs=true,tabsize=2,linenos=true,breaklines,firstnumber=537]{cpp}
    }

\end{minted}
\end{frame}

\begin{frame}[fragile]
\label{Matrix.hpp}
\frametitle{Matrix.hpp}
\hyperlink{conclusion}{Retour au début}\newline
\begin{minted}[obeytabs=true,tabsize=2,linenos=true, breaklines]{cpp}
#ifndef MATRIX_HPP
#define MATRIX_HPP

#include <array>
#include <initializer_list>
#include <concepts>
#include <type_traits>
#include <ostream>

template <typename T>
concept arithmetic = std::is_arithmetic_v<T>;

template <typename T>
concept MathObj = requires(T lhs, T rhs) {
	lhs + rhs == rhs + lhs;
	lhs - rhs;
	lhs *rhs;
};

template <MathObj T, std::size_t LINES, std::size_t COLUMNS = LINES>
class Matrix final
{
	static_assert(LINES > 0 || COLUMNS > 0,
				  "The matrix cannot have a zero size");

public:
\end{minted}
\end{frame}

\begin{frame}[fragile]
\hyperlink{conclusion}{Retour au début}\newline
\begin{minted}[obeytabs=true,tabsize=2,linenos=true,breaklines,firstnumber=27]{cpp}
	Matrix() = default;
	constexpr Matrix(T const initialValue);
	constexpr Matrix(std::initializer_list<std::initializer_list<T>> const &matrix);

	using SubBloc = Matrix<T, LINES / 2, COLUMNS / 2>;
	constexpr Matrix(SubBloc const &a, SubBloc const &b, SubBloc const &c, SubBloc const &d);

	constexpr ~Matrix() noexcept = default;

	Matrix(Matrix const &matrix) = default;
	Matrix &operator=(Matrix const &matrix) = default;

	Matrix(Matrix &&matrix) noexcept = default;
	Matrix &operator=(Matrix &&matrix) = default;

	constexpr std::size_t numberLines() const noexcept;
	constexpr std::size_t numberColumns() const noexcept;

	constexpr T const &operator()(std::size_t const lines,
								  std::size_t const columns) const noexcept;
	constexpr T &operator()(std::size_t const lines, std::size_t const columns);

	constexpr Matrix<T, LINES, COLUMNS>
	operator+=(Matrix<T, LINES, COLUMNS> const &matrix) noexcept;

	template <typename U, std::size_t N, std::size_t M>
	constexpr friend Matrix<U, N, M>
	operator+(Matrix<U, N, M> matrix_left, Matrix<U, N, M> const &matrix_right) noexcept;

	constexpr Matrix<T, LINES, COLUMNS>
\end{minted}
\end{frame}

\begin{frame}[fragile]
\hyperlink{conclusion}{Retour au début}\newline
\begin{minted}[obeytabs=true,tabsize=2,linenos=true,breaklines,firstnumber=57]{cpp}
	operator-=(Matrix<T, LINES, COLUMNS> const &matrix) noexcept;

	template <MathObj U, std::size_t N, std::size_t M>
	constexpr friend Matrix<U, N, M>
	operator-(Matrix<U, N, M> matrix_left, Matrix<U, N, M> const &matrix_right) noexcept;

	constexpr Matrix<T, LINES, COLUMNS> &operator*=(int const multiplier) noexcept;

	template <MathObj U, std::size_t N, std::size_t M, arithmetic V>
	constexpr friend Matrix<U, N, M> operator*(Matrix<U, N, M> matrix,
											   V const multiplier) noexcept;

	template <MathObj U, std::size_t N, std::size_t M, arithmetic V>
	constexpr friend Matrix<U, N, M> operator*(V const multiplier,
											   Matrix<U, N, M> matrix) noexcept;

	template <MathObj U, std::size_t LINES_M1, std::size_t SHARED,
			  std::size_t COLUMNS_M2>
	constexpr friend Matrix<U, LINES_M1, COLUMNS_M2>
	operator*(Matrix<U, LINES_M1, SHARED> const &matrix_left,
			  Matrix<U, SHARED, LINES_M1> const &matrix_right) noexcept;

	template <MathObj U, std::size_t N, std::size_t M>
	friend std::ostream &operator<<(std::ostream &out,
									Matrix<U, N, M> const &matrix);

	constexpr Matrix<T, COLUMNS, LINES> transposed() const noexcept;

	constexpr Matrix<T, 1, COLUMNS>
	line(std::size_t const indexLine) const noexcept;
\end{minted}
\end{frame}

\begin{frame}[fragile]
\hyperlink{conclusion}{Retour au début}\newline
\begin{minted}[obeytabs=true,tabsize=2,linenos=true,breaklines,firstnumber=87]{cpp}

	constexpr Matrix<T, LINES, 1>
	column(std::size_t const indexColumn) const noexcept;

	constexpr static Matrix<T, LINES, COLUMNS> identity() noexcept;

	template <std::size_t N, std::size_t M>
	constexpr Matrix<T, LINES*N, COLUMNS*M> tensoriel_product(Matrix<T, N, M> const& rhs) const noexcept;

private:
	constexpr std::size_t offset(std::size_t const lines,
								 std::size_t const columns) const noexcept;

	constexpr std::array<std::size_t, COLUMNS> strSizeMax() const;
	constexpr std::size_t strSize(T const value) const;

	constexpr std::array<T, COLUMNS * LINES>
	array_matrix_with_initializer_list(std::initializer_list<std::initializer_list<T>> const &matrix);

	std::array<T, LINES * COLUMNS> m_matrix{0};
};

#include "Matrix.tpp"


\end{minted}
\end{frame}

\begin{frame}[fragile]
\label{Matrix.tpp}
\frametitle{Matrix.tpp}
\hyperlink{conclusion}{Retour au début}\newline
\begin{minted}[obeytabs=true,tabsize=2,linenos=true, breaklines]{cpp}
#include <cassert>
#include <algorithm>
#include <string>
#include <ostream>
#include <iomanip>
#include <cmath>
#include "Matrix.hpp"

template <MathObj T, std::size_t LINES, std::size_t COLUMNS>
constexpr Matrix<T, LINES, COLUMNS>::Matrix(T const initialValue)
{
	std::fill(std::begin(m_matrix), std::end(m_matrix), initialValue);
}

template <MathObj T, std::size_t LINES, std::size_t COLUMNS>
constexpr Matrix<T, LINES, COLUMNS>::Matrix(std::initializer_list<std::initializer_list<T>> const & matrix) :
m_matrix()
{
	m_matrix = array_matrix_with_initializer_list(matrix);
}

template <MathObj T, std::size_t LINES, std::size_t COLUMNS>
constexpr Matrix<T, LINES, COLUMNS>::Matrix(SubBloc const& a, SubBloc const& b, SubBloc const& c, SubBloc const& d)
 : m_matrix()
{
	static_assert( (LINES % 2 == 0 || COLUMNS % 2 == 0) 
\end{minted}
\end{frame}

\begin{frame}[fragile]
\hyperlink{conclusion}{Retour au début}\newline
\begin{minted}[obeytabs=true,tabsize=2,linenos=true,breaklines,firstnumber=27]{cpp}
						&& "Le constructeur par bloc n'est valable que pour des matrices de taille pair" );
	constexpr std::size_t N { LINES/2 };
	constexpr std::size_t M { COLUMNS/2 };

	for (std::size_t i {0}; i < N; i++)
	{
		for (std::size_t j {0}; j < M; j++)
		{
			m_matrix[offset(i, j)] = a(i, j);
			m_matrix[offset(i, j + M)] = b(i, j);
			m_matrix[offset(i + N, j)] = c(i, j);
			m_matrix[offset(i + N, j + M)] = d(i, j);
		}
	}
}

template <MathObj T, std::size_t LINES, std::size_t COLUMNS>
constexpr std::size_t
Matrix<T, LINES, COLUMNS>::numberLines() const noexcept
{
	return LINES;
}

template <MathObj T, std::size_t LINES, std::size_t COLUMNS>
constexpr std::size_t
Matrix<T, LINES, COLUMNS>::numberColumns() const noexcept
{
	return COLUMNS;
}

\end{minted}
\end{frame}

\begin{frame}[fragile]
\hyperlink{conclusion}{Retour au début}\newline
\begin{minted}[obeytabs=true,tabsize=2,linenos=true,breaklines,firstnumber=57]{cpp}
template <MathObj T, std::size_t LINES, std::size_t COLUMNS>
constexpr std::size_t
Matrix<T, LINES, COLUMNS>::offset(std::size_t const lines,
								  std::size_t const columns) const noexcept
{
	assert(lines < numberLines() && "Column requested invalid.");
	assert(columns < numberColumns() && "Line requested invalid.");
	return lines * numberColumns() + columns;
}

template <MathObj T, std::size_t LINES, std::size_t COLUMNS>
constexpr T const &
Matrix<T, LINES, COLUMNS>::operator()(std::size_t const lines,
									  std::size_t const columns) const noexcept
{
	return m_matrix[offset(lines, columns)];
}

template <MathObj T, std::size_t LINES, std::size_t COLUMNS>
constexpr T &
Matrix<T, LINES, COLUMNS>::operator()(std::size_t const lines,
									  std::size_t const columns)
{
	return m_matrix[offset(lines, columns)];
}

template <MathObj T, std::size_t LINES, std::size_t COLUMNS>
constexpr Matrix<T, LINES, COLUMNS>
Matrix<T, LINES, COLUMNS>::operator+=(Matrix<T, LINES, COLUMNS> const &matrix) noexcept
{
\end{minted}
\end{frame}

\begin{frame}[fragile]
\hyperlink{conclusion}{Retour au début}\newline
\begin{minted}[obeytabs=true,tabsize=2,linenos=true,breaklines,firstnumber=87]{cpp}
	for (std::size_t i{0}; i < numberLines(); i++)
	{
		for (std::size_t j{0}; j < numberColumns(); j++)
		{
			(*this)(i, j) += matrix(i, j);
		}
	}
	return *this;
}

template <MathObj T, std::size_t LINES, std::size_t COLUMNS>
constexpr Matrix<T, LINES, COLUMNS>
operator+(Matrix<T, LINES, COLUMNS> matrix_left,
		  Matrix<T, LINES, COLUMNS> const &matrix_right) noexcept
{
	return matrix_left += matrix_right;
}

template <MathObj T, std::size_t LINES, std::size_t COLUMNS>
constexpr Matrix<T, LINES, COLUMNS>
Matrix<T, LINES, COLUMNS>::operator-=(Matrix<T, LINES, COLUMNS> const &matrix) noexcept
{
	assert(matrix.numberLines() == numberLines() && "The subtraction of two matrixes, requires that they be of the "
													"same size");
	assert(matrix.numberColumns() == numberColumns() && "The subtraction of two matrixes, requires that they be of the "
														"same size");

	for (std::size_t i{0}; i < numberLines(); i++)
	{
		for (std::size_t j{0}; j < numberColumns(); j++)
\end{minted}
\end{frame}

\begin{frame}[fragile]
\hyperlink{conclusion}{Retour au début}\newline
\begin{minted}[obeytabs=true,tabsize=2,linenos=true,breaklines,firstnumber=117]{cpp}
		{
			(*this)(i, j) -= matrix(i, j);
		}
	}
	return *this;
}

template <MathObj T, std::size_t LINES, std::size_t COLUMNS>
constexpr Matrix<T, LINES, COLUMNS>
operator-(Matrix<T, LINES, COLUMNS> matrix_left,
		  Matrix<T, LINES, COLUMNS> const &matrix_right) noexcept
{
	return matrix_left -= matrix_right;
}

template <MathObj T, std::size_t LINES, std::size_t COLUMNS>
constexpr Matrix<T, LINES, COLUMNS> &
Matrix<T, LINES, COLUMNS>::operator*=(int const multiplier) noexcept
{
	for (std::size_t i{0}; i < numberLines(); i++)
	{
		for (std::size_t j{0}; j < numberColumns(); j++)
		{
			(*this)(i, j) *= multiplier;
		}
	}
	return *this;
}

template <MathObj T, std::size_t LINES, std::size_t COLUMNS, arithmetic V>
\end{minted}
\end{frame}

\begin{frame}[fragile]
\hyperlink{conclusion}{Retour au début}\newline
\begin{minted}[obeytabs=true,tabsize=2,linenos=true,breaklines,firstnumber=147]{cpp}
constexpr Matrix<T, LINES, COLUMNS>
operator*(Matrix<T, LINES, COLUMNS> matrix, V const multiplier) noexcept
{
	return matrix *= multiplier;
}

template <MathObj T, std::size_t LINES, std::size_t COLUMNS, arithmetic V>
constexpr Matrix<T, LINES, COLUMNS>
operator*(V const multiplier, Matrix<T, LINES, COLUMNS> matrix) noexcept
{
	return matrix * multiplier;
}

template <MathObj T, std::size_t LINES_M1, std::size_t SHARED,
		  std::size_t COLUMNS_M2>
constexpr Matrix<T, LINES_M1, COLUMNS_M2>
operator*(Matrix<T, LINES_M1, SHARED> const &matrix_left,
		  Matrix<T, SHARED, COLUMNS_M2> const &matrix_right) noexcept
{

	Matrix<T, LINES_M1, COLUMNS_M2> destination{};
	for (std::size_t i{0}; i < matrix_left.numberLines(); i++)
	{
		for (std::size_t j{0}; j < matrix_right.numberColumns(); j++)
		{
			T summe = 0;
			for (std::size_t k{0}; k < matrix_left.numberColumns(); k++)
			{
				summe += matrix_left(i, k) * matrix_right(k, j);
			}
\end{minted}
\end{frame}

\begin{frame}[fragile]
\hyperlink{conclusion}{Retour au début}\newline
\begin{minted}[obeytabs=true,tabsize=2,linenos=true,breaklines,firstnumber=177]{cpp}
			destination(i, j) = summe;
		}
	}
	return destination;
}

template <MathObj T, std::size_t LINES, std::size_t COLUMNS>
constexpr std::size_t Matrix<T, LINES, COLUMNS>::strSize(T const value) const
{
	std::size_t size{1};
	if constexpr (std::is_integral<T>())
	{
		if (value != 0)
		{
			size = static_cast<std::size_t>(std::log10(value)) + 1;
		}
	}
	else
	{
		std::string str{std::to_string(value)};
		size = std::size(str);
	}
	return size;
}

template <MathObj T, std::size_t LINES, std::size_t COLUMNS>
constexpr std::array<std::size_t, COLUMNS> Matrix<T, LINES, COLUMNS>::strSizeMax() const
{
	std::array<std::size_t, COLUMNS> maxSize {};
	for (std::size_t i {0}; i < LINES; i++)
\end{minted}
\end{frame}

\begin{frame}[fragile]
\hyperlink{conclusion}{Retour au début}\newline
\begin{minted}[obeytabs=true,tabsize=2,linenos=true,breaklines,firstnumber=207]{cpp}
	{
		for (std::size_t j {0}; j < COLUMNS; j++)
		{
			std::size_t size = strSize(m_matrix[offset(i, j)]);
			if (size > maxSize[j])
			{
				maxSize[j] = size;
			}
		}
	}
	return maxSize;
}

template <MathObj T, std::size_t LINES, std::size_t COLUMNS>
std::ostream &
operator<<(std::ostream &out, Matrix<T, LINES, COLUMNS> const &matrix)
{
	out << std::setfill(' ');
	if constexpr (LINES == 1)
	{
		out << "( ";
		std::for_each(std::begin(matrix.m_matrix), std::end(matrix.m_matrix), 
		[&out](T const& v) -> void
		{
			out << v << ' ';
		});
		out << ')' << std::endl;
		return out;
	}
	else
\end{minted}
\end{frame}

\begin{frame}[fragile]
\hyperlink{conclusion}{Retour au début}\newline
\begin{minted}[obeytabs=true,tabsize=2,linenos=true,breaklines,firstnumber=237]{cpp}
	{
		out << "/";
	}
	std::array<std::size_t, COLUMNS> maxSize {matrix.strSizeMax()};
	for (std::size_t i{0}; i < matrix.numberLines(); i++)
	{
		if (i == matrix.numberLines() -1)
		{
			out << '\\';
		}
		else if (i > 0)
		{
			out << '|';
		}
		for (std::size_t j{0}; j < matrix.numberColumns(); j++)
		{
			std::size_t const size = maxSize[j] - matrix.strSize(matrix(i, j));
			if (size > 0)
			{
				out << std::setw(static_cast<int>(size)+1);
			}
			out << ' ' << matrix(i, j);
		}
		if (i == 0)
		{
			out << " \\" << std::endl;
		}
		else if (i < matrix.numberLines() -1)
		{
			out << " |" << std::endl;
\end{minted}
\end{frame}

\begin{frame}[fragile]
\hyperlink{conclusion}{Retour au début}\newline
\begin{minted}[obeytabs=true,tabsize=2,linenos=true,breaklines,firstnumber=267]{cpp}
		}
	}
	if constexpr (LINES == 1)
	{
		out << ")" << std::endl;
	}
	else
	{
		out << " / " << std::endl;
	}
	return out;
}

template <MathObj T, std::size_t LINES, std::size_t COLUMNS>
constexpr Matrix<T, COLUMNS, LINES>
Matrix<T, LINES, COLUMNS>::transposed() const noexcept
{
	Matrix<T, COLUMNS, LINES> trans{};
	for (std::size_t i{0}; i < numberLines(); i++)
	{
		for (std::size_t j{0}; j < numberColumns(); j++)
		{
			trans(j, i) = (*this)(i, j);
		}
	}
	return trans;
}

template <MathObj T, std::size_t LINES, std::size_t COLUMNS>
constexpr Matrix<T, 1, COLUMNS>
\end{minted}
\end{frame}

\begin{frame}[fragile]
\hyperlink{conclusion}{Retour au début}\newline
\begin{minted}[obeytabs=true,tabsize=2,linenos=true,breaklines,firstnumber=297]{cpp}
Matrix<T, LINES, COLUMNS>::line(std::size_t const indexLine) const noexcept
{
	assert(indexLine < numberLines() && "index outside the matrix");
	Matrix<T, 1, COLUMNS> cpyLine{};
	for (std::size_t i{0}; i < numberColumns(); i++)
	{
		cpyLine(0, i) = (*this)(indexLine, i);
	}
	return cpyLine;
}

template <MathObj T, std::size_t LINES, std::size_t COLUMNS>
constexpr Matrix<T, LINES, 1>
Matrix<T, LINES, COLUMNS>::column(
	std::size_t const indexColumn) const noexcept
{
	assert(indexColumn < numberColumns() && "index outside the matrix");
	Matrix<T, LINES, 1> cpyLine{};
	for (std::size_t i{0}; i < numberLines(); i++)
	{
		cpyLine(i, 0) = (*this)(i, indexColumn);
	}
	return cpyLine;
}

template <MathObj T, std::size_t LINES, std::size_t COLUMNS>
constexpr std::array<T, COLUMNS*LINES>
Matrix<T, LINES, COLUMNS>::array_matrix_with_initializer_list(
	std::initializer_list<std::initializer_list<T>> const & matrix
)
\end{minted}
\end{frame}

\begin{frame}[fragile]
\hyperlink{conclusion}{Retour au début}\newline
\begin{minted}[obeytabs=true,tabsize=2,linenos=true,breaklines,firstnumber=327]{cpp}
{
	std::array<T, COLUMNS*LINES> array_matrix = {};
	auto it_mat_dst { std::begin(array_matrix) };
	for (std::initializer_list<T> const & e : matrix)
	{
		assert(std::size(e) <= COLUMNS && "Value entry out of matrix");
		std::move(std::begin(e), std::end(e), it_mat_dst);
		it_mat_dst += COLUMNS;
	}
	return array_matrix;
}

template <MathObj T, std::size_t LINES, std::size_t COLUMNS>
constexpr Matrix<T, LINES, COLUMNS> Matrix<T, LINES, COLUMNS>::identity() noexcept
{
	static_assert(LINES == COLUMNS && "not identity matrix in non square matrix");
	Matrix<T, LINES, COLUMNS> Id {};
	for (std::size_t i { 0 }; i < LINES; i++)
	{
		Id(i, i) = static_cast<T>(1);
	}
	return Id;
}

template <MathObj T, std::size_t LINES, std::size_t COLUMNS>
template <std::size_t N, std::size_t M>
constexpr Matrix<T, LINES*N, COLUMNS*M> Matrix<T, LINES, COLUMNS>::tensoriel_product(Matrix<T, N, M> const& rhs) const noexcept
{
	Matrix<T, LINES*N, COLUMNS*M> res {};
	for (std::size_t i {0}; i < LINES; i++)
\end{minted}
\end{frame}

\begin{frame}[fragile]
\hyperlink{conclusion}{Retour au début}\newline
\begin{minted}[obeytabs=true,tabsize=2,linenos=true,breaklines,firstnumber=357]{cpp}
	{
		for (std::size_t j {0}; j < COLUMNS; j++)
		{
			for (std::size_t k {0}; k < N; k++)
			{
				for (std::size_t l {0}; l < M; l++)
				{
					res(k + i*N, l + j * M) = m_matrix[offset(i, j)] * rhs(k, l);
				}
			}
		}
	}
	return res;

\end{minted}
\end{frame}

\begin{frame}[fragile]
\label{testmovepromotion.cpp}
\frametitle{test-move-promotion.cpp}
\hyperlink{conclusion}{Retour au début}\newline
\begin{minted}[obeytabs=true,tabsize=2,linenos=true, breaklines]{cpp}
#include <Board.hpp>
#include <Piece.hpp>
#include <Coord.hpp>
#include <TypePiece.hpp>
#include <Move.hpp>
#include <math_utility.hpp>

int test_move_promotion(int argc, char **argv)
{
    Board<3> board{
        {W_ROOK, Piece(), Piece()},
        {B_PAWN, Piece(), Piece()},
        {Piece(), Piece(), Piece()}};
    Move m = Move_promote(Coord(1, 0), Coord(2, 0), TypePiece::BISHOP);
    board.move_promotion(m);
    return !(board(2, 0).get_type() == TypePiece::BISHOP &&
             double_equal(board.get_proba(Coord(1, 0)), 0.) &&
             board(1, 0).get_type() == TypePiece::EMPTY);

\end{minted}
\end{frame}

\begin{frame}[fragile]
\label{testmoveenpassantwithcapture.cpp}
\frametitle{test-move-enpassant-with-capture.cpp}
\hyperlink{conclusion}{Retour au début}\newline
\begin{minted}[obeytabs=true,tabsize=2,linenos=true, breaklines]{cpp}
#include <Board.hpp>
#include <Piece.hpp>
#include <Coord.hpp>
#include <TypePiece.hpp>
#include <Move.hpp>
#include <math_utility.hpp>

int test_move_enpassant_with_capture(int argc, char **argv)
{
    Board<5> board{
        {W_ROOK, Piece(), Piece(), Piece(), Piece()},
        {B_PAWN, Piece(), Piece(), Piece(), Piece()},
        {Piece(), Piece(), Piece(), Piece(), Piece()},
        {Piece(), W_PAWN, Piece(), Piece(), Piece()},
        {Piece(), Piece(), Piece(), W_BISHOP, Piece()}};
    Move m = Move_split(Coord(4, 3), Coord(3, 2), Coord(2, 1));
    Move m1 = Move_classic(Coord(3, 1), Coord(1, 1));
    Move m2 = Move_classic(Coord(1, 0), Coord(2, 1));
    board.move(m);
    board.move(m1);
    board.move(m2);
    bool equal = double_equal(board.get_proba(Coord(3, 2)), 0.5);
    return !(board(2, 1).get_type() == TypePiece::PAWN &&
             equal &&
             board(1, 1).get_type() == TypePiece::EMPTY &&
             board(1, 0).get_type() == TypePiece::EMPTY);

\end{minted}
\end{frame}

\begin{frame}[fragile]
\label{testmovepawntwostep.cpp}
\frametitle{test-move-pawn-two-step.cpp}
\hyperlink{conclusion}{Retour au début}\newline
\begin{minted}[obeytabs=true,tabsize=2,linenos=true, breaklines]{cpp}
#include <math_utility.hpp>
#include <Board.hpp>
#include <Piece.hpp>
#include <Coord.hpp>
#include <TypePiece.hpp>
#include <Move.hpp>

int test_move_pawn_two_step(int argc, char **argv)
{
    Board<5> board{
        {W_ROOK, Piece(), Piece(), Piece(), Piece()},
        {B_PAWN, Piece(), Piece(), Piece(), Piece()},
        {Piece(), Piece(), Piece(), Piece(), Piece()},
        {Piece(), Piece(), Piece(), Piece(), Piece()},
        {Piece(), Piece(), Piece(), Piece(), Piece()}};
    Move m = Move_classic(Coord(1, 0), Coord(3, 0));
    board.move(m);
    return !(board(3, 0).get_type() == TypePiece::PAWN &&
             double_equal(board.get_proba(Coord(1, 0)), 0.) &&
             board(1, 0).get_type() == TypePiece::EMPTY);

\end{minted}
\end{frame}

\begin{frame}[fragile]
\label{testmovepromotionwithcapture.cpp}
\frametitle{test-move-promotion-with-capture.cpp}
\hyperlink{conclusion}{Retour au début}\newline
\begin{minted}[obeytabs=true,tabsize=2,linenos=true, breaklines]{cpp}
#include <Board.hpp>
#include <Piece.hpp>
#include <Coord.hpp>
#include <TypePiece.hpp>
#include <Move.hpp>
#include <math_utility.hpp>

int test_move_promotion_with_capture(int argc, char **argv)
{
    Board<3> board{
        {W_ROOK, Piece(), Piece()},
        {B_PAWN, Piece(), Piece()},
        {Piece(), W_QUEEN, Piece()}};
    Move m = Move_promote(Coord(1, 0), Coord(2, 1), TypePiece::BISHOP);
    board.move(m);
    return !(board(2, 1).get_type() == TypePiece::BISHOP &&
             board(2, 1).get_color() == Color::BLACK &&
             double_equal(board.get_proba(Coord(1, 0)), 0.) &&
             board(1, 0).get_type() == TypePiece::EMPTY);

\end{minted}
\end{frame}

\begin{frame}[fragile]
\label{testmovepromotionwithsplitbefore.cpp}
\frametitle{test-move-promotion-with-split-before.cpp}
\hyperlink{conclusion}{Retour au début}\newline
\begin{minted}[obeytabs=true,tabsize=2,linenos=true, breaklines]{cpp}
#include <Board.hpp>
#include <Piece.hpp>
#include <Coord.hpp>
#include <TypePiece.hpp>
#include <Move.hpp>
#include <math_utility.hpp>

int test_move_promotion_with_split_before(int argc, char **argv)
{
    Board<3> board{
        {W_ROOK, Piece(), Piece()},
        {B_PAWN, Piece(), Piece()},
        {Piece(), W_QUEEN, Piece()}};
    Move m1 = Move_split(Coord(2, 1), Coord(2, 0), Coord(2, 2));
    Move m2 = Move_promote(Coord(1, 0), Coord(2, 0), TypePiece::BISHOP);
    board.move(m1);
    board.move(m2);
    bool res = !((board(2, 0).get_type() == TypePiece::BISHOP &&
                  board(2, 0).get_color() == Color::BLACK &&
                  board(1, 0).get_type() == TypePiece::EMPTY &&
                  double_equal(board.get_proba(Coord(2, 2)), 1.)) ||
                 (board(2, 0).get_type() == TypePiece::QUEEN &&
                  board(2, 0).get_color() == Color::WHITE &&
                  board(1, 0).get_type() == TypePiece::PAWN &&
                  double_equal(board.get_proba(Coord(2, 2)), 0.)));
    return res;

\end{minted}
\end{frame}

\begin{frame}[fragile]
\label{testmovesplitwithmesure.cpp}
\frametitle{test-move-split-with-mesure.cpp}
\hyperlink{conclusion}{Retour au début}\newline
\begin{minted}[obeytabs=true,tabsize=2,linenos=true, breaklines]{cpp}
#include <Board.hpp>
#include <Piece.hpp>
#include <Coord.hpp>
#include <TypePiece.hpp>
#include <Move.hpp>
#include <math_utility.hpp>

int test_move_split_with_mesure(int argc, char* *argv)
{
    Board<5> board{
        {W_ROOK, Piece(), Piece(),Piece(), Piece()},
        {Piece(), Piece(), Piece(), Piece(), Piece()},
        {Piece(), Piece(), Piece(),Piece(),Piece()},
        {Piece(), Piece(),  Piece(), Piece(),Piece()},
        {Piece(), Piece(), Piece(),B_ROOK,Piece()}};
    Move m = Move_split(Coord(4,3), Coord(4, 1), Coord(4,4));
    Move m1 = Move_split(Coord(0, 0), Coord(0, 1), Coord(0 ,4));
    Move m2 = Move_classic(Coord(0,1), Coord(4,1));
    Move m3 = Move_classic(Coord(4, 4), Coord(0, 4));
    board.move(m);
    board.move(m1);
    board.move(m2);
    board.move(m3);
    std::size_t acc{0};
    for(std::size_t i{0}; i<5; i++)
    {
\end{minted}
\end{frame}

\begin{frame}[fragile]
\hyperlink{conclusion}{Retour au début}\newline
\begin{minted}[obeytabs=true,tabsize=2,linenos=true,breaklines,firstnumber=27]{cpp}
        for(std::size_t j{0}; j<5; j++)
        {
        if(board(i, j).get_type() == TypePiece::ROOK)
        {
            acc++;
        }
        }
    }
    return acc != 2 && acc != 1;

\end{minted}
\end{frame}

\begin{frame}[fragile]
\label{testsplitinnonemptycase.cpp}
\frametitle{test-split-in-non-empty-case.cpp}
\hyperlink{conclusion}{Retour au début}\newline
\begin{minted}[obeytabs=true,tabsize=2,linenos=true, breaklines]{cpp}
#include <Board.hpp>
#include <Piece.hpp>
#include <Move.hpp>
#include <Coord.hpp>
#include <math_utility.hpp>

int test_split_in_non_empty_case(int argc, char** argv)
{
    Board<1, 4> board {
        {W_ROOK, Piece(), Piece(), Piece()}
    };

    Move m1 {Move_split(Coord(0, 0), Coord(0, 1), Coord(0, 3))};
    Move m2 {Move_split(Coord(0, 1), Coord(0, 2), Coord(0, 3))};

    board.move(m1);
    board.move(m2);
    
    bool ret {
        board(0, 1).get_type() == TypePiece::ROOK &&
        board(0, 2).get_type() == TypePiece::ROOK &&
        board(0, 3).get_type() == TypePiece::ROOK &&
        double_equal(board.get_proba(Coord(0, 1)), 0.5) &&
        double_equal(board.get_proba(Coord(0, 2)), 0.25) &&
        double_equal(board.get_proba(Coord(0, 3)), 0.25)
    };
\end{minted}
\end{frame}

\begin{frame}[fragile]
\hyperlink{conclusion}{Retour au début}\newline
\begin{minted}[obeytabs=true,tabsize=2,linenos=true,breaklines,firstnumber=27]{cpp}
    return !ret;

\end{minted}
\end{frame}

\begin{frame}[fragile]
\label{testslidemergemovethroughasplitpiece.cpp}
\frametitle{test-slide-merge-move-through-a-split-piece.cpp}
\hyperlink{conclusion}{Retour au début}\newline
\begin{minted}[obeytabs=true,tabsize=2,linenos=true, breaklines]{cpp}
#include <Board.hpp>
#include <Piece.hpp>
#include <Coord.hpp>
#include <TypePiece.hpp>
#include <Move.hpp>
#include <math_utility.hpp>

int test_slide_merge_move_through_a_split_piece(int argc, char **argv)
{
    Board<2, 5> board{
        {W_ROOK, Piece(), Piece(), Piece(), Piece()},
        {Piece(), Piece(), B_ROOK, Piece(), Piece()}};
    Move m = Move_split(Coord(1, 2), Coord(0, 2), Coord(1, 4));
    Move m1 = Move_split(Coord(0, 0), Coord(0, 1), Coord(0, 4));
    Move m2 = Move_merge(Coord(0, 1), Coord(0, 4), Coord(0, 3));
    board.move(m);
    board.move(m1);
    board.move(m2);
    bool res{board(0, 2).get_type() == TypePiece::ROOK &&
             board(0, 2).get_color() == Color::BLACK &&
             board(0, 1).get_type() == TypePiece::ROOK &&
             board(0, 1).get_color() == Color::WHITE &&
             double_equal(board.get_proba(Coord(0, 1)), 0.5) &&
             board(0, 3).get_type() == TypePiece::ROOK &&
             board(0, 3).get_color() == Color::WHITE &&
             double_equal(board.get_proba(Coord(0, 3)), 0.5) &&
\end{minted}
\end{frame}

\begin{frame}[fragile]
\hyperlink{conclusion}{Retour au début}\newline
\begin{minted}[obeytabs=true,tabsize=2,linenos=true,breaklines,firstnumber=27]{cpp}
             board(0, 4).get_type() == TypePiece::EMPTY};
    return !res;

\end{minted}
\end{frame}

\begin{frame}[fragile]
\label{testgetprobamoveslidecapture.cpp}
\frametitle{test-get-proba-move-slide-capture.cpp}
\hyperlink{conclusion}{Retour au début}\newline
\begin{minted}[obeytabs=true,tabsize=2,linenos=true, breaklines]{cpp}
#include <Board.hpp>
#include <Piece.hpp>
#include <Coord.hpp>
#include <TypePiece.hpp>
#include <Move.hpp>
#include <math_utility.hpp>

int test_get_proba_move_slide_capture(int argc, char **argv)
{
    Board<3, 5> board{
        {W_ROOK, Piece(), Piece(), Piece(), Piece()},
        {Piece(), W_ROOK,  W_ROOK,  W_ROOK, Piece()},
        {Piece(),Piece(),Piece(),  Piece(), B_ROOK}};
    Move m1 = Move_split(Coord(1, 1), Coord(0, 1), Coord(2, 1));
    Move m2 = Move_split(Coord(1, 2), Coord(0, 2), Coord(2, 2));
    Move m3 = Move_split(Coord(1, 3), Coord(0, 3), Coord(2, 3));
    Move m4 = Move_split(Coord(2, 4), Coord(0, 4), Coord(1, 4));
    Move m5 = Move_classic(Coord(0, 4), Coord(0, 0));
    board.move(m1);
    board.move(m2);
    board.move(m3);
    board.move(m4);
    bool res{double_equal(board.get_proba_move(m5), 1./16)};
    return !res;

\end{minted}
\end{frame}

\begin{frame}[fragile]
\label{testgetprobamoveslideonsamecolor.cpp}
\frametitle{test-get-proba-move-slide-on-same-color.cpp}
\hyperlink{conclusion}{Retour au début}\newline
\begin{minted}[obeytabs=true,tabsize=2,linenos=true, breaklines]{cpp}
#include <Board.hpp>
#include <Piece.hpp>
#include <Coord.hpp>
#include <TypePiece.hpp>
#include <Move.hpp>
#include <math_utility.hpp>

int test_get_proba_move_slide_on_same_color(int argc, char **argv)
{
    Board<3, 5> board{
        {Piece(), Piece(), Piece(), Piece(), Piece()},
        {B_QUEEN, W_ROOK,  W_ROOK,  W_ROOK, Piece()},
        {Piece(),Piece(),Piece(),  Piece(), B_ROOK}};
    Move m1 = Move_split(Coord(1, 1), Coord(0, 1), Coord(2, 1));
    Move m2 = Move_split(Coord(1, 2), Coord(0, 2), Coord(2, 2));
    Move m3 = Move_split(Coord(1, 3), Coord(0, 3), Coord(2, 3));
    Move m4 = Move_split(Coord(2, 4), Coord(0, 4), Coord(1, 4));
    Move m = Move_split(Coord(1, 0), Coord(0, 0), Coord(2, 0));
    Move m5 = Move_classic(Coord(0, 4), Coord(0, 0));
    board.move(m1);
    board.move(m2);
    board.move(m3);
    board.move(m4);
    board.move(m);
    bool res{double_equal(board.get_proba_move(m5), 1./2)};
    return !res;

\end{minted}
\end{frame}

\begin{frame}[fragile]
\label{testgetprobamoveslidewithouttarget.cpp}
\frametitle{test-get-proba-move-slide-without-target.cpp}
\hyperlink{conclusion}{Retour au début}\newline
\begin{minted}[obeytabs=true,tabsize=2,linenos=true, breaklines]{cpp}
#include <Board.hpp>
#include <Piece.hpp>
#include <Coord.hpp>
#include <TypePiece.hpp>
#include <Move.hpp>
#include <math_utility.hpp>

int test_get_proba_move_slide_without_target(int argc, char **argv)
{
    Board<3, 5> board{
        {Piece(), Piece(), Piece(), Piece(), Piece()},
        {Piece(), W_ROOK,  W_ROOK,  W_ROOK, Piece()},
        {Piece(),Piece(),Piece(),  Piece(), B_ROOK}};
    Move m1 = Move_split(Coord(1, 1), Coord(0, 1), Coord(2, 1));
    Move m2 = Move_split(Coord(1, 2), Coord(0, 2), Coord(2, 2));
    Move m3 = Move_split(Coord(1, 3), Coord(0, 3), Coord(2, 3));
    Move m4 = Move_split(Coord(2, 4), Coord(0, 4), Coord(1, 4));
    Move m5 = Move_classic(Coord(0, 4), Coord(0, 0));
    board.move(m1);
    board.move(m2);
    board.move(m3);
    board.move(m4);
    bool res{double_equal(board.get_proba_move(m5), 1.)};
    return !res;

\end{minted}
\end{frame}

\begin{frame}[fragile]
\label{testgetprobamovejumpsamecolor.cpp}
\frametitle{test-get-proba-move-jump-same-color.cpp}
\hyperlink{conclusion}{Retour au début}\newline
\begin{minted}[obeytabs=true,tabsize=2,linenos=true, breaklines]{cpp}
#include <Board.hpp>
#include <Piece.hpp>
#include <Coord.hpp>
#include <TypePiece.hpp>
#include <Move.hpp>
#include <math_utility.hpp>

int test_get_proba_move_jump_same_color(int argc, char **argv)
{
    Board<3> board{
        { W_KING, Piece(), Piece()},
        {Piece(), Piece(), Piece()},
        {Piece(), Piece(), W_KNIGHT}};
    Move m1 = Move_split(Coord(2, 2), Coord(0, 1), Coord(1, 0));
    Move m5 = Move_classic(Coord(0, 0), Coord(0, 1));
    board.move(m1);
    bool res{double_equal(board.get_proba_move(m5), 1./2)};
    return !res;

\end{minted}
\end{frame}

\begin{frame}[fragile]
\label{testgetprobamovejumpcapture.cpp}
\frametitle{test-get-proba-move-jump-capture.cpp}
\hyperlink{conclusion}{Retour au début}\newline
\begin{minted}[obeytabs=true,tabsize=2,linenos=true, breaklines]{cpp}
#include <Board.hpp>
#include <Piece.hpp>
#include <Coord.hpp>
#include <TypePiece.hpp>
#include <Move.hpp>
#include <math_utility.hpp>

int test_get_proba_move_jump_capture(int argc, char **argv)
{
    Board<3> board{
        { W_KING, Piece(), Piece()},
        {Piece(), Piece(), Piece()},
        {Piece(), Piece(), B_KNIGHT}};
    Move m1 = Move_split(Coord(2, 2), Coord(0, 1), Coord(1, 0));
    Move m5 = Move_classic(Coord(0, 0), Coord(0, 1));
    board.move(m1);
    bool res{double_equal(board.get_proba_move(m5), 1.)};
    return !res;

\end{minted}
\end{frame}

\begin{frame}[fragile]
\label{testcaptureslidemovetrue.cpp}
\frametitle{test-capture-slide-move-true.cpp}
\hyperlink{conclusion}{Retour au début}\newline
\begin{minted}[obeytabs=true,tabsize=2,linenos=true, breaklines]{cpp}
#include <Board.hpp>
#include <Piece.hpp>
#include <Coord.hpp>
#include <TypePiece.hpp>
#include <Move.hpp>
#include <Color.hpp>
#include <math_utility.hpp>

int test_capture_slide_move_true(int argc, char **argv)
{
    Board<3, 5> board{
        {W_ROOK, Piece(), Piece(), Piece(), Piece()},
        {Piece(), W_ROOK,  W_ROOK,  W_ROOK, Piece()},
        {Piece(),Piece(),Piece(),  Piece(), B_ROOK}};
    Move m1 = Move_split(Coord(1, 1), Coord(0, 1), Coord(2, 1));
    Move m2 = Move_split(Coord(1, 2), Coord(0, 2), Coord(2, 2));
    Move m3 = Move_split(Coord(1, 3), Coord(0, 3), Coord(2, 3));
    Move m4 = Move_split(Coord(2, 4), Coord(0, 4), Coord(1, 4));
    Move m5 = Move_classic(Coord(0, 4), Coord(0, 0));
    board.move(m1);
    board.move(m2);
    board.move(m3);
    board.move(m4);
    board.move(m5, true);
    bool res{board(0, 0).get_type() == TypePiece::ROOK && board(0, 0).get_color() == Color::BLACK };
    return !res;

\end{minted}
\end{frame}

\begin{frame}[fragile]
\label{testcaptureslidemovefalse.cpp}
\frametitle{test-capture-slide-move-false.cpp}
\hyperlink{conclusion}{Retour au début}\newline
\begin{minted}[obeytabs=true,tabsize=2,linenos=true, breaklines]{cpp}
#include <Board.hpp>
#include <Piece.hpp>
#include <Coord.hpp>
#include <TypePiece.hpp>
#include <Move.hpp>
#include <Color.hpp>
#include <math_utility.hpp>

int test_capture_slide_move_false(int argc, char **argv)
{
    Board<3, 5> board{
        {W_ROOK, Piece(), Piece(), Piece(), Piece()},
        {Piece(), W_ROOK,  W_ROOK,  W_ROOK, Piece()},
        {Piece(),Piece(),Piece(),  Piece(), B_ROOK}};
    Move m1 = Move_split(Coord(1, 1), Coord(0, 1), Coord(2, 1));
    Move m2 = Move_split(Coord(1, 2), Coord(0, 2), Coord(2, 2));
    Move m3 = Move_split(Coord(1, 3), Coord(0, 3), Coord(2, 3));
    Move m4 = Move_split(Coord(2, 4), Coord(0, 4), Coord(1, 4));
    Move m5 = Move_classic(Coord(0, 4), Coord(0, 0));
    board.move(m1);
    board.move(m2);
    board.move(m3);
    board.move(m4);
    board.move(m5, false);
    bool res{board(0, 4).get_type() == TypePiece::ROOK && board(0, 4).get_color() == Color::BLACK };
    return !res;

\end{minted}
\end{frame}

\begin{frame}[fragile]
\label{testmovemergeaftersplit.cpp}
\frametitle{test-move-merge-after-split.cpp}
\hyperlink{conclusion}{Retour au début}\newline
\begin{minted}[obeytabs=true,tabsize=2,linenos=true, breaklines]{cpp}
#include <Board.hpp>
#include <Piece.hpp>
#include <Coord.hpp>
#include <TypePiece.hpp>
#include <Move.hpp>
#include <math_utility.hpp>

int test_move_merge_after_split(int argc, char **argv)
{
    Board<4> board{
        {B_KING, Piece(), Piece(), Piece()},
        {B_ROOK, Piece(),  Piece(), Piece()},
        {Piece(),Piece(),Piece(),  Piece()},
        {W_QUEEN, Piece(), Piece(), W_KING}};
    Move m1 = Move_split(Coord(3, 0), Coord(3, 1), Coord(3, 2));
    Move m2 = Move_merge(Coord(3, 1), Coord(3, 2), Coord(3, 0));
    board.move(m1);
    board.move(m2, true);
    bool res{double_equal(board.get_proba(Coord(3,0)), 1.)};
    return !res;

\end{minted}
\end{frame}

\begin{frame}[fragile]
\label{testsplitaftersplit.cpp}
\frametitle{test-split-after-split.cpp}
\hyperlink{conclusion}{Retour au début}\newline
\begin{minted}[obeytabs=true,tabsize=2,linenos=true, breaklines]{cpp}
#include <Board.hpp>
#include <Piece.hpp>
#include <Coord.hpp>
#include <TypePiece.hpp>
#include <Move.hpp>
#include <math_utility.hpp>

int test_split_after_split(int argc, char **argv)
{
    Board<4> board{
        {Piece(), Piece(), Piece(), Piece()},
        {Piece(), B_KING,  B_ROOK, Piece()},
        {Piece(),Piece(),Piece(),  Piece()},
        {W_QUEEN, Piece(), Piece(), W_KING}};
    Move m1 = Move_split(Coord(3, 0), Coord(2, 1), Coord(3, 2));
    Move m2 = Move_split(Coord(3, 2), Coord(2, 1), Coord(1, 0));
    board.move(m1);
    board.move(m2, true);
    bool res{double_equal(board.get_proba(Coord(3,3)), 1.)&&
    double_equal(board.get_proba(Coord(2,1)), 1./2) &&
    double_equal(board.get_proba(Coord(1,0)), 1./2)};
    return !res;

\end{minted}
\end{frame}

\begin{frame}[fragile]
\label{testsplitaftersplit2.cpp}
\frametitle{test-split-after-split2.cpp}
\hyperlink{conclusion}{Retour au début}\newline
\begin{minted}[obeytabs=true,tabsize=2,linenos=true, breaklines]{cpp}
#include <Board.hpp>
#include <Piece.hpp>
#include <Coord.hpp>
#include <TypePiece.hpp>
#include <Move.hpp>
#include <math_utility.hpp>

int test_split_after_split2(int argc, char **argv)
{
    Board<4> board{
        {Piece(), Piece(), Piece(), Piece()},
        {Piece(), B_KING,  B_ROOK, Piece()},
        {Piece(),Piece(),Piece(),  Piece()},
        {W_QUEEN, Piece(), Piece(), W_KING}};
    Move m1 = Move_split(Coord(3, 0), Coord(2, 1), Coord(3, 2));
    Move m2 = Move_split(Coord(3, 2), Coord(1, 0), Coord(2, 1));
    board.move(m1);
    board.move(m2, true);
    bool res{double_equal(board.get_proba(Coord(3,3)), 1.)&&
    double_equal(board.get_proba(Coord(2,1)), 1./4) &&
    double_equal(board.get_proba(Coord(1,0)), 1./4) &&
    double_equal(board.get_proba(Coord(3,2)), 1./2) };
    return !res;

\end{minted}
\end{frame}

\begin{frame}[fragile]
\label{testmultiplesplit.cpp}
\frametitle{test-multiple-split.cpp}
\hyperlink{conclusion}{Retour au début}\newline
\begin{minted}[obeytabs=true,tabsize=2,linenos=true, breaklines]{cpp}
#include <Board.hpp>
#include <Piece.hpp>
#include <Coord.hpp>
#include <TypePiece.hpp>
#include <Move.hpp>
#include <math_utility.hpp>

int test_multiple_split(int argc, char **argv)
{
    Board<3> board{
        {Piece(), Piece(), B_BISHOP},
        {Piece(), B_KING,  Piece()},
        {Piece(),Piece(),Piece()}};
    Move m1 = Move_split(Coord(1, 1), Coord(2, 1), Coord(2, 2));
    Move m2 = Move_split(Coord(2, 2), Coord(1, 1), Coord(1,2));
    Move m3 = Move_classic(Coord(1, 2), Coord(2, 2));
    Move m4 = Move_split(Coord(1, 1), Coord(0, 0), Coord(0, 1));
    board.move(m1);
    board.move(m2);
    board.move(m3);
    board.move(m4);
    bool res{double_equal(board.get_proba(Coord(0,2)), 1.)};
    return !res;

\end{minted}
\end{frame}

\begin{frame}[fragile]
\label{testsplittakes.cpp}
\frametitle{test-split-takes.cpp}
\hyperlink{conclusion}{Retour au début}\newline
\begin{minted}[obeytabs=true,tabsize=2,linenos=true, breaklines]{cpp}
#include <Board.hpp>
#include <Piece.hpp>
#include <Coord.hpp>
#include <TypePiece.hpp>
#include <Move.hpp>
#include <math_utility.hpp>

int test_split_takes(int argc, char **argv)
{
    Board<3> board{
        {Piece(), Piece(), W_BISHOP},
        {Piece(), B_KING,  Piece()},
        {Piece(),Piece(),Piece()}};
    Move m1 = Move_split(Coord(1, 1), Coord(2, 1), Coord(2, 0));
    Move m2 = Move_split(Coord(2, 1), Coord(1, 1), Coord(2, 2));
    Move m3 = Move_classic(Coord(0, 2), Coord(1, 1));
    Move m4 = Move_classic(Coord(1, 1), Coord(2, 0));
    board.move(m1);
    board.move(m2);
    board.move(m3);
    board.move(m4);
    bool res{double_equal(board.get_proba(Coord(2,0)), 1.)};
    return !res;

\end{minted}
\end{frame}
}